\documentclass[10pt]{amsart}
\usepackage[a4paper,margin=23mm]{geometry}
\usepackage{amsmath,amssymb,amsthm,mathtools,hyperref,graphicx,xspace,letltxmacro,xpatch}
\usepackage[english]{babel}
\newtheorem{theo}{Theorem}[section]
\newtheorem{lemm}[theo]{Lemma}\newtheorem{prop}[theo]{Proposition}\newtheorem{coro}[theo]{Corollary}\newtheorem{conj}[theo]{Conjecture}
\theoremstyle{definition}\newtheorem{defi}[theo]{Definition}\newtheorem{rema}[theo]{Remark}\newtheorem{exam}[theo]{Example}\newtheorem{nota}[theo]{Notation}
\newcommand{\ie}{i.e.}\newcommand{\eg}{e.g.}\newcommand{\resp}{resp.}\newcommand{\skpt}{\par\smallskip\noindent}
\emergencystretch=3em\pagestyle{plain}
%~Mouliné par MaN_auto v.0.27.3 2023-06-20 11:00:03

\usepackage{amssymb}
\usepackage{mathtools}


%%%MATHOP, MATHRM%%%%%%%%%%%%%%%%%%%%%%%%%%%%%%
\DeclareMathOperator{\DO}{DO}
\DeclareMathOperator{\spann}{span}
\DeclareMathOperator{\ad}{ad}
\DeclareMathOperator{\diag}{diag}
\DeclareMathOperator{\supp}{supp}
\DeclareMathOperator{\Op}{Op}

\DeclareMathOperator{\spec}{spec}
\DeclareMathOperator{\dom}{dom}
\DeclareMathOperator{\vol}{vol}
\DeclareMathOperator{\dist}{dist}

\DeclareMathOperator{\Id}{Id}
\DeclareMathOperator{\id}{id}
\newcommand{\Mat}{\mathrm{Mat}}
\newcommand{\de}{\, \mathrm{d}}

\newcommand{\RB}{\mathrm{B}}
\newcommand{\RC}{\mathrm{C}}
\newcommand{\RD}{\mathrm{D}}
\newcommand{\RH}{\mathrm{H}}
\newcommand{\RL}{\mathrm{L}}
\newcommand{\rmI}{\mathrm{I}}
\newcommand{\rmi}{\mathrm{i}}
\newcommand{\rme}{\mathrm{e}}
\newcommand{\rmII}{\mathrm{II}}
\newcommand{\CAP}{\mathrm{CAP}}
\newcommand{\EP}{\mathrm{EP}}
\renewcommand{\mathsf}{\mathrm}
\newcommand{\sfH}{\mathsf{H}}
\newcommand{\sfB}{\mathsf{B}}
\newcommand{\sfL}{\mathsf{L}}
\newcommand{\sfC}{\mathsf{C}}

%%%%%%%%%MATHBB%%%%%%%%%%%%%%%%%%
\newcommand{\R}{\mathbb R}
\newcommand{\C}{\mathbb C}
\newcommand{\N}{\mathbb N}
\newcommand{\BBS}{\mathbb S}
\newcommand{\Z}{\mathbb Z}


%%%%%%%%%%MATHCAL%%%%%%%%%%%%%%%%%%

\newcommand{\CA}{\mathcal{A}}
\newcommand{\CB}{\mathcal{B}}
\newcommand{\CC}{\mathcal{C}}
\newcommand{\CD}{\mathcal{D}}
\newcommand{\CE}{\mathcal{E}}
\newcommand{\CG}{\mathcal{G}}
\newcommand{\CI}{\mathcal{I}}
\newcommand{\CJ}{\mathcal{J}}
\newcommand{\CK}{\mathcal{K}}
\newcommand{\CL}{\mathcal{L}}
\newcommand{\CM}{\mathcal{M}}
\newcommand{\CN}{\mathcal{N}}
\newcommand{\CO}{\mathcal{O}}
\newcommand{\CR}{\mathcal{R}}
\newcommand{\CS}{\mathcal{S}}
\newcommand{\CT}{\mathcal{T}}
\newcommand{\CU}{\mathcal{U}}
\newcommand{\CV}{\mathcal{V}}
\newcommand{\CX}{\mathcal{X}}
\newcommand{\CY}{\mathcal{Y}}
\newcommand{\CZ}{\mathcal{Z}}

\newcommand{\COD}{{\CO\CD}}
\newcommand{\CNR}{{\CN\!\CR}}

%%%%%%%%BOLDSYMBOL%%%%%%%%%%%%%%%%%%%%%%%
\newcommand{\BPsi}{\boldsymbol \Psi}

\newcommand{\btheta}{\boldsymbol \theta}
\newcommand{\bzeta}{\boldsymbol \zeta}
\newcommand{\bxi}{\boldsymbol \xi}
\newcommand{\boeta}{\boldsymbol \eta}
\newcommand{\bone}{\boldsymbol 1}
\newcommand{\bzero}{\boldsymbol 0}
%%%%%%%%%%%%%%MATHBF%%%%%%%%%%%%%%%%%%%
\newcommand{\ba}{\mathbf a}
\newcommand{\bg}{\mathbf g}
\newcommand{\bb}{\mathbf b}
\newcommand{\BA}{\mathbf A}
\newcommand{\BB}{\mathbf B}
\newcommand{\BD}{\mathbf D}
\newcommand{\BE}{\mathbf E}
\newcommand{\BS}{\mathbf S}
\newcommand{\BR}{\mathbf R}
\newcommand{\BP}{\mathbf P}
\newcommand{\BT}{\mathbf T}
\newcommand{\BU}{\mathbf U}
\newcommand{\BX}{\mathbf X}
\newcommand{\BY}{\mathbf Y}
\newcommand{\bx}{\mathbf x}
\newcommand{\bbb}{\mathbf b}
\newcommand{\be}{\mathbf e}
\newcommand{\bk}{\mathbf k}
\newcommand{\bp}{\mathbf p}
\newcommand{\bm}{\mathbf m}
\newcommand{\by}{\mathbf y}
\newcommand{\br}{\mathbf r}
\newcommand{\bu}{\mathbf u}
\newcommand{\bv}{\mathbf v}
\newcommand{\bn}{\mathbf n}

%%%%%%%%%%MATHFRAK%%%%%%%%%%%%%%%%%%

\newcommand{\Fc}{\mathfrak c}
\newcommand{\FA}{\mathfrak A}
\newcommand{\FB}{\mathfrak B}
\newcommand{\FD}{\mathfrak D}
\newcommand{\FE}{\mathfrak{E}}
\newcommand{\FH}{\mathfrak H}
\newcommand{\FK}{\mathfrak K}
\newcommand{\FL}{\mathfrak L}
\newcommand{\FS}{\mathfrak S}
\newcommand{\FT}{\mathfrak T}
\newcommand{\FU}{\mathfrak U}
\newcommand{\FV}{\mathfrak V}


\renewcommand{\phi}{\varphi}
\newcommand{\del}{\partial}
\newcommand{\resc}{s}
\newcommand{\rco}{\chi}
\newcommand{\pd}{uncoupleable}
\newcommand{\mpu}{competely uncoupleable}
\newcommand{\eps}{\varepsilon}
\renewcommand{\epsilon}{\varepsilon}


\newcommand{\act}{\triangleright}



\renewcommand{\subseteq}{\subset}
\renewcommand{\supseteq}{\supset}



\newcommand{\bigo}[1]{O\left(#1 \right)}
\newcommand{\smallo}[1]{o\left(#1 \right)}





%\renewcommand{\tilde}{\widetilde}
\newcommand{\flo}[1]{\left\lfloor#1\right\rfloor}
\newcommand{\ceil}[1]{\left\lceil#1\right\rceil}

\newcommand{\abs}[1]{\left\lvert #1 \right\rvert}
\newcommand{\set}[1]{\left\{ #1 \right\}}
\newcommand{\norm}[1]{\left\| #1 \right\|}

\newcommand{\snorm}[3]{\left\| #1\right\|_{#3}^{(#2)}}


\newcommand{\ang}[1]{\langle #1 \rangle}
\newcommand{\er}{r}



\makeatletter
\def\editors#1{%
\def\editor@name{#1}
\if@francais
\def\editor@string{Recommand\'e par les \'editeurs \editor@name.}
\else
\def\editor@string{Recommended by Editors \editor@name.}
\fi}
\makeatother


%%%%%%%%%%%%%%%%%%%%%%%%%%%%%%%%%%%%%%%%%%%%
\graphicspath{{./figures/}}

\newcommand*{\mk}{\mkern -1mu}
\newcommand*{\Mk}{\mkern -2mu}
\newcommand*{\mK}{\mkern 1mu}
\newcommand*{\MK}{\mkern 2mu}

%\hypersetup{urlcolor=purple, linkcolor=blue, citecolor=red}

\newcommand*{\romanenumi}{\renewcommand*{\theenumi}{\roman{enumi}}}
\newcommand*{\Romanenumi}{\renewcommand*{\theenumi}{\Roman{enumi}}}
\newcommand*{\alphenumi}{\renewcommand*{\theenumi}{\alph{enumi}}}
\newcommand*{\Alphenumi}{\renewcommand*{\theenumi}{\Alph{enumi}}}
\let\oldtilde\tilde
\renewcommand*{\tilde}[1]{\mathchoice{\widetilde{#1}}{\widetilde{#1}}{\oldtilde{#1}}{\oldtilde{#1}}}
\let\oldhat\hat
\renewcommand*{\hat}[1]{\mathchoice{\widehat{#1}}{\widehat{#1}}{\oldhat{#1}}{\oldhat{#1}}}
\let\oldforall\forall
\renewcommand*{\forall}{\mathrel{\oldforall}}
\let\oldtocsection=\tocsection

\let\oldtocsubsection=\tocsubsection

%\let\oldtocsubsubsection=\tocsubsubsection
%%%%%%%%%%%%%%%%%%%%%%%%%%%%%%%%%%%%%%%%%%%%%%%%%%


\NewDocumentEnvironment{enonce*}{ O{} m O{} }{\par\smallskip\noindent\textbf{#2\if\relax\detokenize{#3}\relax\else\ (#3)\fi.}\ }{\par\smallskip}
\NewDocumentEnvironment{enonce}{ O{} m O{} }{\par\smallskip\refstepcounter{theo}\noindent\textbf{#2~\thetheo\if\relax\detokenize{#3}\relax\else\ (#3)\fi.}\ }{\par\smallskip}

\let\exem\relax\let\endexem\relax\newtheorem{exem}[theo]{Example}
\title{Bethe–Sommerfeld half-lines with quantitative overlap for periodic elliptic systems}
\author{Lagacé, Jean, Morozov, Sergey, Parnovski, Leonid, Pfirsch, Bernhard, Shterenberg, Roman}
\date{doi:10.5802/ahl.184; CC BY 4.0}
\begin{document}
\maketitle
\noindent\textit{Editorial note.} Complete original mathematical body and original publisher bibliography retained. General main theorem selected once; Dirac specializations and intermediate reduction theorems not counted separately. Only original class environments locally rendered.\par\medskip
\noindent\textit{Allowed standard prerequisites (editorial).}\begin{itemize}
\item Standard functional analysis, measure theory and the author-cited classical theorems. Standard prior results stay cited; complete source-local proofs and definitions retained.
\end{itemize}\medskip
\tableofcontents


\section{Introduction}



\subsection{A Gauge transform}

During the last fifteen years, substantial progress has been made in the spectral theory of periodic and almost periodic scalar operators. An important tool that was developed during this period and was used to obtain asymptotic spectral results was the method of gauge transform (see, e.g., \cite{Ivrii2018,MorParSht2014,ParSob2010,
ParSht2012,
ParSht2016,ParSht2019,Sobolev2005,Sobolev2006}), sometimes also called the quantum Birkhoff normal form~\cite{LaurentVuNgoc}. This method, which heavily uses commutator estimates, was originally created for classical pseudo-differential operators (see e.g.~\cite{Rozenbljum1978,Weinstein1977}) but was then modified to the periodic case by Sobolev~\cite{Sobolev2005,Sobolev2006} and to the almost periodic setting by Parnovski and Shterenberg~\cite{ParSht2012}. The aim of this paper is to describe the method of gauge transform on an abstract level and then apply this abstract scheme to a concrete example --- elliptic systems of operators (including Dirac operators).

Here is the basic setting: suppose that we are given an operator
\begin{equation}\label{LP1}
A = A_0 + B,
\end{equation}
where $A_0$ is a diagonal operator in a given basis and $B$ is a perturbation, which is assumed to be small in some sense. The standard example which the reader may want to keep in mind is
\begin{equation}
A_0 = \diag\left(a_1(-\Delta)^{\alpha/2},\,\dotsc,\,a_m(-\Delta)^{\alpha/2}\right),
\end{equation}
where $\alpha>0$, $0 \ne a_j \in \R$, and $B$ is a pseudo-differential perturbation of order smaller than $\alpha$ with periodic or almost periodic coefficients. For instance, a Dirac operator with an almost periodic potential can be brought to such a form by a unitary transformation. In many applications we will furthermore require $A$ to be self-adjoint, even though our general scheme may not always require it.

We want to find an operator $A'$ that is unitarily equivalent to $A$ and is simpler --- either diagonal or, failing this, has a form
\begin{equation}\label{LP2}
A' = U A U^{-1} = A_0' + B',
\end{equation}
where $A_0'$ is diagonal, $U$ is unitary, and $B'$ is a perturbation that is smaller than $B$. The notion of ``smallness'' assumes that we have a small parameter, and $B'$ has this small parameter entering in a higher power than $B$. The most common example of application to PDEs assumes that the order of $B'$ is smaller than the order of $B$ (so the role of the small parameter is played by the inverse of the energy), but in some cases the small parameter can be chosen to be a coupling constant, see~\cite{ParSht2019}. The operators $A$ and $A'$ have the same spectrum and the hope is that it is easier to describe the spectrum of $A'$, both quantitatively and qualitatively. As an example of the spectral properties we want to study, we list the following two types of problems:
\begin{enumerate}
\item Obtaining asymptotic expansions for the so-called \emph{integrated density of states} $N(A;\lambda)$ as the spectral parameter $\lambda$ tends to $\pm \infty$;
\item If $B$ has periodic coefficients, to prove that whenever $A$ is unbounded above (resp. below), its spectrum contains a semi-axis $[\lambda_0,\infty)$ (resp. $(-\infty,\lambda_0])$. Such an operator $A$ is said to satisfy the
\emph{Bethe--Sommerfeld property}.
\end{enumerate}

If we seek the unitary operator $U$ in~\eqref{LP2} in the form $U=\exp(i \Psi)$, then we have
\begin{equation}\label{LP3}
A' = A_0 + B+i [A_0,\Psi] + i[B,\Psi]-\frac12 [[A_0,\Psi],\Psi]-\frac12 [[B,\Psi],\Psi]+R,
\end{equation}
where $R$ consists of further terms given by formally expanding the series for the exponentials $\exp(i\Psi)$. Our hope is to solve the equation
\begin{equation}\label{LP4}
B + i[A_0,\Psi]=0
\end{equation}
for $\Psi$, so that the second and third terms of~\eqref{LP3} cancel each other. Ideally, the rest of the terms (starting from the fourth one) would indeed be smaller than $B$. In most cases, however, these two wishes turn out to be infeasible.

The main obstacle is that solutions $\Psi$ to equation~\eqref{LP4} involve a denominator that could be small for some $B$ (for example, to have any hope of solving~\eqref{LP4}, the diagonal part of $B$ has to be absent). Therefore, we usually have to modify our procedure and divide the perturbation $B$ into two parts --- good (or non-resonant) part $B^{\CNR}$ for which the equation
\begin{equation}\label{LP5}
B^{\CNR} + i[A_0,\Psi]=0
\end{equation}
has a nice solution $\Psi^{\CNR}$ and bad (or resonant) part $B^{\CR}=B-B^{\CNR}$ which we will be unable to destroy using our procedure. Thus, at the end we will have
\begin{equation}\label{LP6}
A' = A_0' + B^{\CR}+B',
\end{equation}
where $B'$ is smaller (in order, say) than $B$. Of course, we also hope that the resonant part $B^{\CR}$ is better in some sense than the initial perturbation $B$; in many applications, the operator $B^{\CR}$ acts in subspaces of our Hilbert space generated by ``specially designated and geometrically defined'' areas of the phase space.


After we have reduced our operator to the improved form~\eqref{LP6}, in principle we can repeat the same procedure --- finitely, or even infinitely many times. The latter process is much more difficult to realise, and we will not give examples of it in this paper. However, in many settings we indeed have to run this procedure several times (more than once) in order to achieve the desired ``smallness'' of the remainder. In other words, we construct the ``improved'' operator in the form
\begin{equation}\label{LP7}
A_n = \exp(i \Psi_n)\dots\exp(i \Psi_2)\exp(i \Psi_1)A\exp(-i \Psi_1)\exp(-i \Psi_2)\dots\exp(-i \Psi_n).
\end{equation}
We call this method {\emph {the serial gauge transform}}. Sometimes, it is more convenient to look for the improved operator in the form
\begin{equation}\label{LP8}
A^{(n)} = \exp(i (\Psi_n+\dotso+ \Psi_2+\Psi_1))A\exp(-i (\Psi_1+\Psi_2+\dotso +\Psi_n)),
\end{equation}
which we call {\emph {the parallel gauge transform}}. In both situations, the operators $\Psi_j$ are solutions of equations similar in form to~\eqref{LP5}.

Another important distinction between different variations of the gauge transform is as follows. In order to prove that the order of the remainder $B'$ is smaller than the order of $B$, we have to estimate the orders of various commutators. Sometimes, it is enough to have the basic estimate: the order of the commutator is not greater than the sum of the orders of its entries. This estimate holds without any restrictions, but for it to be effective we need to have some a priori inequalities between the orders of the principal term $A_0$ and the perturbation $B$; we call this approach the weak gauge transform. On the other hand, quite often we can improve our estimate on commutators: for example, in the classical scalar pseudo-differential calculus, the order of the commutator can be estimated by the sum of the orders of the entries minus one. If we have such an estimate, we can guarantee that the order of $B'$ is indeed smaller than the order of $B$, assuming nothing other than that the order of $A_0$ is larger than the order of $B$. This approach is called the strong gauge transform. In this paper, we will define the weak and strong gauge transforms rigorously and give a general abstract setting in which they can be applied. We discuss the advantages and drawbacks of both types of gauge transforms and finish with a couple of concrete applications.

The first application is to obtain asymptotic expansions for the density of states of elliptic almost periodic operator systems. Under some technical conditions described later, we may either obtain complete or limited expansions as the spectral parameter goes to $\pm\infty$. The other application is to prove that some elliptic periodic systems have the Bethe--Sommerfeld property. This will be done under the same conditions that allow us to obtain a complete asymptotic expansion for the density of states. In either of these cases, some Dirac operators are examples of those to which we can apply our results.

\subsection{Description of the results for elliptic systems and the Dirac operator}

While describing the precise class of operators $A$ for which we obtain spectral asymptotics requires definitions that are made later, we can make these results explicit for Dirac operators in dimension~$2$ and $3$ perturbed by classical pseudo-differential almost periodic operators right away. The two-dimensional Dirac operator with mass $M$ acts in $\RL^2(\R^2;\C^2)$ and is given by
\begin{equation}
\BA_{2,M} := -i(\sigma_1 \del_{x_1} + \sigma_2 \del_{x_2}) + \sigma_3 M,
\end{equation}
where $\sigma_1, \sigma_2,
\sigma_3$ are the Pauli matrices
\begin{align}
\sigma_1=
\begin{pmatrix}
0 & 1\\
1 & 0
\end{pmatrix},\qquad
\sigma_2=
\begin{pmatrix}
0 & -i\\
i & 0
\end{pmatrix},\qquad \text{and} \qquad
\sigma_3=
\begin{pmatrix}
1 & 0\\
0 & -1
\end{pmatrix}.
\end{align}

The three-dimensional Dirac operator with mass $M$ acts in $\RL^2(\R^3;\C^4)$ and is given by
\begin{equation}
\BA_{3,M} := -i \left(\gamma_1 \del_{x_1} + \gamma_2 \del_{x_2} + \gamma_3
\del_{x_3} \right) + \Gamma M,
\end{equation}
where the matrices $\gamma_j$, $\Gamma$ are the Dirac matrices (see~\cite{Upmeier})\footnote{Many authors would write $\alpha_j$ for $\gamma_j$ and $\beta$ for $\Gamma$, see e.g.~\cite{thaller}. We keep our convention in line with higher-dimensional generalisations and to avoid some notational conflicts later on.}
\begin{equation}
\begin{aligned}
\gamma_j=
\begin{pmatrix}
\mathbf 0 & \sigma_j\\
\sigma_j & \mathbf 0
\end{pmatrix}, \quad \text{and} \quad\
\Gamma=
\begin{pmatrix}
\Id_2 & \mathbf 0\\
\mathbf 0 & -\Id_2
\end{pmatrix}.
\end{aligned}
\end{equation}
We obtain asymptotic expansions for the density of states of operators of the type $\BA = \BA_{d,M} + \BB$ under the assumption that $\BB$ is a ``generic'' almost periodic pseudo-differential perturbation. The precise meaning of generic is given in Section~\ref{sec:besicovitch}. The density of states for elliptic differential operators $A$ that are not semi-bounded can be defined by the formula
\begin{equation}
N(\lambda;A) := \lim_{L\,\to\,\infty} \frac{N(\lambda;A^{(L)})}{(2L)^d}.
\end{equation}
Here, $A^{(L)}$ is the restriction of $A$ to the cube $[-L,L]^d$ with periodic boundary conditions, and $N(\lambda;A^{(L)})$ is the counting function for the discrete eigenvalues of $A^{(L)}$ in the interval $[0,\lambda)$ when $\lambda > 0$ and $(\lambda,0]$ when $\lambda < 0$. Later, we will give several equivalent definitions of $N(\lambda)$ which are more convenient to work with and allow pseudo-differential perturbations.

\begin{theo}\label{thm:dirac2dids}
Let $\BA = \BA_{2,M} + \BB$, where $\BB$ is a generic symmetric pseudo-differential operator with almost periodic coefficients of order $\beta < 1$ acting in $\RL^2(\R^2;\C^2)$. Then, there is a complete asymptotic expansion for the density of states of $\BA$ in the sense that for every $K > -2$, there is a finite set $L \subset (0,2+K)$ and constants $C_{j}^\pm$, $C_{j,\log}^\pm$, $j \in L \cup \set 0$ such that
\begin{equation}
N(\pm\lambda;\BA) = C_0^\pm \lambda^2 + \sum_{j\,\in\,L}\left(C_{j}^\pm
\lambda^{2 - j} + C_{j,\log}^\pm \lambda^{2-j} \log \lambda \right)+ \bigo{\lambda^{-K}}
\end{equation}
as $\lambda \to \infty$.
\end{theo}
Theorem~\ref{thm:aexpconcrete} is a more general version of Theorem~\ref{thm:dirac2dids}. It is applicable to elliptic systems of pseudodifferential operators whose principal symbol has only simple eigenvalues.

We obtain a restricted expansion for the three-dimensional case.
\begin{theo}\label{thm:dirac3dids}
Let $\BA = \BA_{3,M} + \BB$, where $\BB$ is a generic operator of the form
\[
\BB = B_1 \gamma_1 + B_2 \gamma_2 + B_3 \gamma_3 + B_\Gamma \Gamma + B_{\Id}
\Id_4,
\]
where each $B_j$, $j \in\set{1,2,3,\Gamma,\Id}$ is a scalar symmetric pseudo-differential operator with almost periodic coefficients of order $\beta$, $0 \le \beta \le 1/2$. Then, writing $\gamma^* = \max \set{\beta - 1, 2 \beta - 1}$ there is a finite set $L \subset (0,1 - \gamma^*)$ and constants $C_{j,q}^\pm$, $j \in L \cup \set{0}$, $q \in \set{0,1,2}$ such that
\begin{equation}
N(\pm\lambda;A) = C_0^\pm \lambda^3 + \sum_{j \in L} \sum_{q = 0}^2 C_{j,q}^\pm
\lambda^{3 - j} \log^q(\lambda) + \bigo{\lambda^{2 +
\gamma^*}}
\end{equation}
as $\lambda \to \infty$.
\end{theo}
This time, it is Theorem~\ref{thm:aexpconcretecut} which is a more general version of Theorem~\ref{thm:dirac3dids}. It is applicable to elliptic systems of pseudodifferential operators whose principal symbol has multiple eigenvalues under some more restrictive conditions on the perturbation.

We also obtain that two-dimensional Dirac operators satisfy the Bethe--Sommerfeld property.
\begin{theo}\label{thm:dirac2dbs}
Let $\BA = \BA_{2,M} + \BB$, where $\BB$ is a symmetric pseudo-differential operator of order $\beta < 1$ with periodic coefficients. Then, there exists $\lambda_0 > 0$ such that the spectrum of $\BA$ contains intervals $(-\infty,-\lambda_0]$ and $[\lambda_0,\infty)$.
\end{theo} 
This theorem also has a more general version in Theorem~\ref{thm:bs}. It is applicable to systems whose principal symbol has only simple eigenvalues.

In Section~\ref{sec:dirac}, we also describe generalisations of these results to higher dimensional Dirac operators, and give some technical conditions under which we can get complete asymptotic expansions or the Bethe--Sommerfeld property for the three-dimensional Dirac operator.

\subsection{Description of the main results and plan of the paper}

In the first half of our paper, we discuss the gauge transform in an abstract setting. The setting is developed while keeping in mind particular applications to almost periodic operators. As such, the space on which the operators act looks like an abstract version of a Besicovitch space. In the second half, we will discuss the specific applications of the results obtained in the first half to elliptic systems of pseudo-differential almost periodic operators; in particular, in the last section we will show that Dirac operators are a specific example of them. An interesting part of the application of our methods to systems is that we need to intertwine and alternate the use of the weak and strong gauge transforms, whereas in the past only one type was used at a time. In order to help the reader familiar with previous literature on the method of gauge transform, we have kept the notation as close as possible to the one used in~\cite{MorParSht2014,ParSob2010}.

\subsubsection*{Plan and results of Part I}
In Section~\ref{sec:operators} we define an algebra of operators $\BS^\infty$ acting on a non separable Hilbert space which should be thought of as an abstract version of a Besicovitch space. For some set $\Xi$ this algebra will be concretely realised on $\ell^2(\Xi)$ through a group action on its basis elements. This algebra is filtered as an algebra of pseudo-differential operators on $\ell^2(\Xi)$, and it has similar properties to those of classical pseudo-differential operators in the PDE sense. Their natural domains are generalisations of Sobolev spaces. This section contains many technical but very useful lemmas describing boundedness properties, adjoints, compositions and commutators of operators in $\BS^\infty$. One of the main differences with classical pseudo-differential operators is illustrated in Proposition~\ref{prop:normorder}, which plays the role of the Calder\'on--Vaillancourt theorem in our setting. It essentially says that we can directly correlate symbol norms of operators with the norms of individual summands in Paley--Wiener type decompositions.

In Section~\ref{sec:perturbation}, we turn our attention to some natural subspaces of $\BS^\infty$ -- operators that are either elliptic or diagonal. Just as in the classical setting, our definition of elliptic operators allows us to characterise natural domains of self-adjointness for operators in the algebra $\BS^\infty$. The three main results of this section illustrate the three most important properties of elliptic operators. In Proposition~\ref{prop:globell}, they are shown to admit a parametrix, and are therefore invertible up to a controllable error. Lemma~\ref{lem:relative boundedness} is used repeatedly throughout the paper and shows that lower order perturbations of elliptic operators are relatively bounded, with explicit bounds. Finally, in Proposition~\ref{prop:closedness and self-adjointness}, we show that elliptic operators are closed and self-adjoint if symmetric.

In Section~\ref{sec:ids}, we consider the situation where operators in $\BS^\infty$ are affiliated to a $\rmI_\infty$ or $\rmII_\infty$ factor. This is common in the study of almost-periodic operators and their generalisations. We define a general notion of density of states measures (DSM) in $\BS^\infty$ as traces in the affiliated $\rmI_\infty$ or $\rmII_\infty$ factor. We give a variational description of the DSM of an interval $J$ even in situations where the operator is not bounded below. This is used to show the principal results of this section: small perturbations of elliptic self-adjoint operators do not change their density of states much. The definition of ``smallness'' of the perturbation is made clear in that section. In Lemma~\ref{lem:monotonicity}, we control to what extent perturbations of smaller order can affect the DSM, whereas in Lemma~\ref{lem:spectralperturb} it is perturbations that are spectrally supported away from the interval $J$ that are shown to have a small effect.

In Section~\ref{sec:gt}, we describe the abstract gauge transform scheme, which is split into two cases: the weak and strong gauge transforms. In both cases, we describe the resonant regions geometrically as subsets of the index set $\Xi$. The serial scheme for the weak gauge transform is described in Lemmas~\ref{lem:psi} and~\ref{lem:onestepweakgauge} and Corollary~\ref{cor:onestepweakgauge}, whereas the parallel scheme is described in Proposition~\ref{prop:symbolestweakgauge}. In both cases, only trivial estimates on the commutator are used. In Lemma~\ref{lem:sgt}, we describe conditions under which a stronger scheme can be used. Since conditions for the strong transform to be applicable are varying in nature, we do not attempt at completely classifying them.

Finally, in Section~\ref{sec:mapo}, we describe the case where the symbols are functions into $\Mat^{m\times m}(\C)$ rather than $\C$. We describe how this can be reduced to the abstract scalar case and introduce a new class of operator systems: uncoupled operators. Our goal is to show that under some specific conditions, elliptic systems are unitarily equivalent to uncoupled operators up to a remainder which we can control. In that light, the main results of this section are Theorems~\ref{thm:systemonestep} and~\ref{thm:gtsystem} which give explicit conditions under which one can use the weak gauge transform to conjugate elliptic symmetric operators into almost uncoupled ones. The remainders are small (in the sense of Section~\ref{sec:ids}) perturbations.


\subsubsection*{Plan and results of Part II}

In the second part, we apply the results of Part I to concrete systems of elliptic pseudo-differential operators with periodic and almost periodic perturbations. More specifically, we study operator systems of the form $\BA =
\BA_0 + \BB$, defined on a dense domain in $\RL^2(\R^d;\C^m)$ where $\BA_0$ is defined as in~\eqref{LP2}, and $\BB$ is a pseudo-differential perturbation of order $\beta < \alpha$. In Section~\ref{sec:besicovitch}, we give a description of these operators in term of Besicovitch space, and we make the relevant definitions concerning periodic operators. In Sections~\ref{sec:prelim} and~\ref{sec:asyexp}, we obtain asymptotic expansions for the IDS. In Section~\ref{sec:periodic}, we describe how periodic operators enter in our framework while exhibiting more structure. In Sections~\ref{sec:bs} and~\ref{sec:cg} we prove that some elliptic systems of operators have the Bethe--Sommerfeld property using some combinatorial geometric arguments. Finally in Section~\ref{sec:dirac} we expose how Dirac operators may fit in our setting.

Since the precise description of the results requires some notations and language defined in Part I, we postpone their description to the beginning of Part II.

\subsection*{Acknowledgements}
The authors are grateful to Marcello Seri and the anonymous referee for their careful reading and comments on this paper, which greatly improved the exposition. 

\section*{\textsc{Part I : An abstract gauge transform scheme}}

\section{Generalised almost-periodic operators} \label{sec:operators}

In this section, we define an algebra of generalised almost-periodic operators. We start by defining the space on which those operators are defined. We also define generalised Sobolev spaces which are their natural domains. We then describe the algebraic properties of the generalised almost-periodic operators, and obtain version of the Calder\'on--Vaillancourt theorem in our context in Proposition~\ref{prop:normorder}.

\subsection{Generalised Sobolev spaces}


Let $\Xi$ be an infinite, possibly uncountable set equipped with a weight function $\ang \cdot : \Xi \to [1,\infty)$. We will often call $\Xi$ {\emph {the index set}}. For $\gamma \in \R$ we define the spaces
\begin{align}
\begin{split} \RH^{\gamma}(\Xi) &{:=} \bigg\{
\begin{matrix}x: \Xi \to \C\\ \displaystyle x:\xi\mapsto x_\xi
\end{matrix}:\quad \sum_{\xi\,\in\,\Xi}\ang{\xi}^{2\gamma}\abs{x_\xi}^2 < \infty\bigg\}
\end{split}
\\
\intertext{and}
{\RH}^\infty(\Xi) &{:=}\bigcap\limits_{\gamma\,\in\,\R}{\RH}^\gamma(\Xi).
\end{align}
In particular, every $x\in \RH^{\gamma}(\Xi)$ vanishes at all but countably many $\xi\in\Xi$. Every $\RH^{\gamma}(\Xi)$ is a Hilbert space with inner product
\begin{equation}\label{eq:inner product}
\left(x,y \right)_{\RH^\gamma(\Xi)} := \sum_{\xi\,\in\,\Xi}\ang{\xi}^{2\gamma}x_\xi\,\overline{y_\xi}.
\end{equation}
It is easy to see that $\RH^0(\Xi) = \ell^2(\Xi)$ with the standard orthonormal basis indexed bijectively from $\Xi$ as
\begin{equation}
\CE := \left\{\be_{\xi}: \xi \in \Xi\right\}, \qquad \be_{\xi} :\eta\in\Xi\mapsto
\begin{cases}
1 & \text{if }\eta =\xi,\\ 0 & \text{if }\eta\neq\xi
\end{cases},
\end{equation}
and that $\RH^{\gamma_1}(\Xi) \subset \RH^{\gamma_2}(\Xi)$ for all $\gamma_1 >\gamma_2\in\R$. It follows from their definition that $\RH^\gamma(\Xi)$ is the completion of $\spann(\CE)$ under the norm generated by~\eqref{eq:inner product}, where $\spann(\CE)$ consists of finite linear combinations from $\CE$, i.e. elements where $x_\xi = 0$ except for finitely many $\xi$. When there is no risk of confusion, we will write $\RH^\gamma:=\RH^\gamma(\Xi)$.

\subsection{An algebra of operators}

Let $G$ be a group that acts from the left on $\Xi$, so that the action is free, i.e.~only the identity of $G$ has fixed points. We denote by $g\act \xi$ the action of $g\in G$ on $\xi\in\Xi$. Starting from the weight function $\langle\cdot\rangle$ on $\Xi$ we define one on $G$ by
\begin{equation}\label{eq:modulusg}
\ang g := 1 + \sup_{\xi\,\in\,\Xi} \left|\ang{g\act\xi} - \ang{\xi}\right|.
\end{equation}
We assume that $G$ has a \emph{bounded range of action}, which means that $\ang g$ is finite for all $g\in G$.

It will be useful for future convenience to observe the following properties of the weight function:
\begin{lemm}
For all $f,g \in G$, $\xi \in \Xi$ and $t \in\R$ the following relations hold:
\begin{enumerate}
\item
\begin{equation}
\ang g = \left\langle g^{-1}\right\rangle;\label{eq:anginv}
\end{equation}
\item Peetre-type inequalities:
\begin{equation}
\ang g^{-1} \ang \xi \le \ang{g \act \xi} \le \ang g \ang \xi\label{eq:peetre}
\end{equation}
and
\begin{equation}
\ang{fg}^t \le \min\left\{\ang f^t \ang g^{\abs t}, \ \ang f^{\abs t} \ang g^{t}\right\}.\label{eq:peetre2}
\end{equation}
\end{enumerate}
\end{lemm}

\begin{proof}
For all $g \in G$, $\xi \in \Xi$ the definition~\eqref{eq:modulusg} implies~\eqref{eq:anginv} and the estimates
\begin{align}\label{bounds on g act xi}
\max\big\{1, 1 +\ang{\xi} -\ang{g}\big\} \leq \ang{g\act \xi} \leq \ang\xi +\ang g -1.
\end{align}
Note the relations
\begin{equation}\label{upper1}
a +1 -b =\big((b -1)(a -b) +a\big)/b \geq a/b,\quad \text{for all } a\geq b\geq 1.
\end{equation}
and
\begin{equation}\label{lower1}
a +b -1 \leq a +b -1 +(a -1)(b -1) =ab,\quad \text{for all } a, b\geq 1.
\end{equation}
The first estimate in~\eqref{eq:peetre} follows from~\eqref{bounds on g act xi} and~\eqref{upper1}, the second from~\eqref{bounds on g act xi} and~\eqref{lower1}. Now by~\eqref{eq:modulusg} and~\eqref{upper1} for all $f,g \in G$ we obtain
\begin{equation}\label{eq:angprod}
\begin{aligned}
\ang{fg}&\leq 1 +\sup_{\xi\,\in\,\Xi} \left|\ang{fg\act\xi} - \ang{g\act \xi}\right| +\sup_{\xi\,\in\,\Xi}\left|\ang{g\act\xi} - \ang{\xi}\right| \\
&=\ang f +\ang g -1 \leq \ang f\ang g,
\end{aligned}
\end{equation}
which implies~\eqref{eq:peetre2} for $t >0$. Now~\eqref{eq:angprod} and~\eqref{eq:anginv} imply
\begin{equation}
\ang g =\left\langle f^{-1}fg\right\rangle \leq\ang f\ang{fg}\text{ and }\ang f =\left\langle fgg^{-1}\right\rangle \leq \ang{fg}\ang{g},
\end{equation}
which delivers~\eqref{eq:peetre2} for $t < 0$. The case $t = 0$ is trivial.
\end{proof}

\begin{defi}
We call a function $b : G \times \Xi \to \C$, $(g,\xi) \mapsto b_g(\xi)$ an almost periodic
\emph{symbol} if there exists a countable set $\Theta \subset G$, closed under inversion and containing the identity $\id_G$, such that for all $g \in G \setminus
\Theta$, $b_g(\xi) \equiv 0$. Whenever there is no risk of confusion, we will write $\id := \id_G$. We call $\Theta$ a \emph{frequency set} for $b$ and the functions $\{b_\theta(\,\cdot\,)\}_{\theta\,\in\,\Theta}$ the \emph{Fourier coefficients} of $b$. For every symbol $b$ and every $\gamma\in\R$, $l\geq 0$, we define the family of norms
\begin{equation}\label{eq:defsnorm}
\snorm{b}{\gamma}{l} := \sum_{\theta\,\in\, \Theta} \ang \theta^l
\sup_{\xi\,\in\,\Xi} \left(\ang{\xi}^{-\gamma}
\left|b_\theta(\xi)\right|\right).
\end{equation}
The class of \emph{symbols of order $\gamma$} is defined as
\begin{equation}\label{eq:sgamma}
\BS^{\gamma}:=\BS^{\gamma}(G,\Xi):=\set{ b:G\times\Xi\to\C :
\snorm{b}{\gamma}{l}<\infty \text{ for all } l\geq 0}.
\end{equation}
\end{defi}

The space of symbols is naturally a linear space. It is clear that if $\Theta$ is a frequency set for a symbol, then any $\Gamma \supset \Theta$ is also one. It is obvious from the definition that $\snorm{\cdot}{\gamma}{l}$ is a decreasing function of $\gamma$ and an increasing function of $l$, thus
\begin{equation}
\BS^{\gamma_1}\subseteq\BS^{\gamma_2},\quad \text{ for all }\gamma_1\leq \gamma_2.
\end{equation}
We introduce
\begin{equation}\label{eq:algebradef}
\BS^\infty:= \bigcup_{\gamma\, \in\,\R} \BS^\gamma \quad \text{ and }
\quad \BS^{-\infty}:= \bigcap_{\gamma\,\in \,\R} \BS^\gamma.
\end{equation}

\begin{lemm}\label{lem:complete}
For every $\gamma\in\R$, the space $\BS^\gamma$ equipped with the family of norms $\lbrace\snorm{\cdot}{\gamma}{l}\rbrace_{l\,\geq\,0}$ is a Fréchet space.
\end{lemm}


\begin{proof}
Consider a sequence
\begin{equation}
(b_n)_{n\,\geq\,1}\subset \BS^\gamma
\end{equation}
that is Cauchy with respect to $\snorm{\cdot}{\gamma}{l}$ for every $l \geq 0$, and denote by $\Theta(n)$ a frequency set for each $b_n$. Then, for all $\theta \in G$, we observe that $b_\theta(\xi):=\lim\limits_{n\,\to\,\infty}(b_n)_\theta(\xi)$ exists and vanishes outside the countable set $\Theta = \bigcup_n \Theta(n)$. It is a simple computation to see that $b\in\BS^\gamma$ with $\snorm{b_n-b}{\gamma}{l}\to 0$, as $n\to\infty$, for all $l\geq 0$. Hence, the claim follows.
\end{proof}

\begin{defi}
Let $b:G\times\Xi\to\C$ be a symbol with frequency set $\Theta\subseteq G$ and
\begin{equation}
\big(b_{\theta}(\xi)\big)_{\theta\,\in\,\Theta}\in \ell^2(\Theta), \text{ for all } \xi\in\Xi.
\end{equation}
Then the \emph{almost periodic linear operator associated to $b$} is
\begin{equation}
B:= \Op(b):\spann(\CE)\to \ell^2(\Xi)
\end{equation}
defined by
\begin{equation}\label{eq:defopB}
B \be_\xi := \sum_{\theta\,\in\,\Theta} b_\theta(\xi) \be_{\theta\,\act\,\xi}, \quad\text{for all }
\xi\in\Xi.
\end{equation}
\end{defi}

\begin{rema}
If $b\in\BS^\infty$, then, in view of~\eqref{eq:defsnorm} and~\eqref{eq:sgamma}, $(b_{\theta}(\xi))_{\theta\,\in\,\Theta}\in \ell^1(\Theta)\subseteq
\ell^2(\Theta)$ holds for all $\xi\in\Xi$. This means that we can associate an almost periodic operator to every symbol in $\BS^\infty$. On the other hand, since the group action of $G$ on $\Xi$ is free, $b$ can be recovered from $B$ via the identity
\begin{equation}\label{eq:symbrecovery}
b_g(\xi)=\left(\be_{g\,\act\,\xi},B\be_\xi\right)_{\ell^2(\Xi)}, \quad\text{for all }g\in G , \xi\in\Xi.
\end{equation}
Thus, there is a one-to-one correspondence between almost periodic symbols and almost periodic operators. This correspondence is in contrast to the case of classical pseudo-differential operators where this correspondence is only modulo smoothing operators. Hence, we allow ourselves to overload the notation and write $B = \Op(b) \in \BS^\gamma$ if $b \in \BS^\gamma$, $\gamma \in \R \cup \set{\pm \infty}$, and let $\snorm B \gamma l := \snorm b
\gamma l$ for all $l \ge 0$, $\gamma \in \R$. Note that this correspondence gets lost if one does not require the group action of $G$ on $\Xi$ to be free. Our construction can be generalised to such non-free group actions, but for simplicity of the exposition we do not do it in this paper.
\end{rema}

We call $B$ \emph{quasi-periodic} if $b$ admits a finite frequency set. A simple example of a quasi-periodic operator of class $\BS^\gamma$, $\gamma\in\R$, is $\Op(h)$ with
\[
h_g(\xi):=
\begin{cases}
\tilde h(\xi) & \text{if } g = \id, \\ 0
&\text{otherwise.}
\end{cases}
\]
Here, $\tilde h$ is a function on $\Xi$ satisfying $\big|\tilde h(\xi)\big| \leq\ang\xi^\gamma$ for all $\xi\in\Xi$.


\begin{rema}
Our terminology is justified by the following example. Suppose that $G$ is a locally compact abelian (LCA) group and $G_B$ is its Bohr compactification, see~\cite[\S 1]{Shubin1978}. Index by $\Xi$ the set of characters $\widetilde{\CE}:=\lbrace \tilde{\be}_\xi:\xi\in\Xi\rbrace$ of $G$ or, equivalently, $G_B$. On $\CAP(G)$, the continuous almost periodic functions on $G$, we can define an inner product $(f,g) = \CM(f\overline g)$, where $\CM(f)$ is the mean of $f$ with respect to the normalised Haar measure on $G_B$. The Besicovitch space $\RB^2(G)$ is defined as the closure of $\CAP(G)$ with respect to the norm induced by this inner product. By~\cite[Proposition~1.5]{Shubin1978}, the map
\begin{equation}
\CE\to\widetilde{\CE}\, \quad \be_\xi\mapsto\tilde{\be}_\xi,
\end{equation}
extends to an isometric isomorphism $\ell^2(\Xi)\to \sfB^2(G)$. In particular, for $G = (\R^d,+)$, one has $\widetilde{\CE}=\lbrace \bx\mapsto \exp(i \bx\cdot
\bxi),\ \bxi\in\R^d\rbrace$ and the operators in $\BS^\infty(\R^d,\R^d)$ correspond to almost periodic pseudo-differential operators in $\sfB^2(\R^d)$ or $\sfL^2(\R^d)$, see~\cite[\S 3--4]{Shubin1978} and~\cite[Equation (8.8)]{ParSht2012}. The present work can be applied to more general settings, for example, when the underlying group $G$ is non-abelian. Note that the Bohr compactification construction is inadequate in that situation, e.g., for $G = \mathrm{SL}(2;\R)$ we have $G_B = \set \id$, see~\cite[p.~4]{Shubin1978}.
\end{rema}


From Lemma~\ref{lem:complete} we obtain the following corollary.
\begin{coro}\label{cor:complete}
Let $(B_n)_{n\,\geq\,1} \subset \BS^\gamma$ be such that
\begin{equation}
\sum_{n\,\geq\,1}\snorm{B_n}{\gamma}{l}<\infty
\end{equation}
for all $l \geq 0$. Then the sum
\begin{equation}
B:=\sum_{n\,\geq\,1}B_n
\end{equation}
converges in $\BS^\gamma$ with
\begin{equation}
\snorm{B}{\gamma}{l}\leq\sum\limits_{n\,\geq\,1}\snorm{B_n}{\gamma}{l}.
\end{equation}
\end{coro}

Up until now, operators from $\BS^\infty$ were only defined on $\spann(\CE)$. We now show that they can be extended in a natural way.
\begin{lemm} \label{lem:domain}
For every $\beta,\gamma\in\R$ the operator $B \in \BS^\gamma$ can be uniquely extended to a bounded linear operator $B : \sfH^{\beta} \to \sfH^{\beta-\gamma}$. Moreover, we have the bound
\begin{equation}
\norm{B}_{\sfH^{\beta}\,\to\,\sfH^{\beta-\gamma}}
\le \snorm{B}{\gamma}{\abs{\beta-\gamma}}.
\end{equation}
\end{lemm}
\begin{proof}
Let $x,y \in \spann(\CE)$, i.e.~$x_\xi=y_\xi=0$ for all but finitely many $\xi$. Then, the Cauchy--Schwarz and Peetre inequalities~\eqref{eq:peetre} imply
\begin{equation}
\begin{aligned}
\left|(x,By)_{\sfH^{\beta-\gamma}} \right|& = \abs{\sum_{\theta\,\in\,\Theta} \sum_{\xi\,\in\,\Xi}
\ang{\theta\act\xi}^{2(\beta - \gamma)} x_{\theta\,\act\,\xi} \overline{ b_\theta(\xi)\,y_\xi}}
\\
&\le \sum_{\theta\,\in\,\Theta} \ang \theta^{\abs{\beta-\gamma}}
\sup_{\zeta\,\in\,\Xi} \left(\ang{\zeta}^{-\gamma}\abs{b_\theta(\zeta)}\right) \times \\
& \qquad \times \left(\sum_{\xi\,\in\,\Xi}\ang{\theta \act \xi}^{2(\beta-\gamma)}
\abs{x_{\theta\act\xi}}^2\right)^{1/2}
\left(\sum_{\xi\,\in\,\Xi} \ang \xi^{2\beta} \abs{ y_\xi}^2\right)^{1/2}\\
&\leq \snorm{B}{\gamma}{\abs{\beta-\gamma}} \|x\|_{\sfH^{\beta-\gamma} } \|y\|_{\sfH^{\beta}}.
\end{aligned}
\end{equation}
The claim follows by density of $\spann\CE$ in $\sfH^{\alpha}$ for all $\alpha\in\R$.
\end{proof}

We obtain the following immediate corollary.

\begin{coro} \label{cor:domain}
Every $B \in \BS^0$ extends to a bounded operator on $\ell^2 =\sfH^0$ such that
\begin{equation}
\|B\|_{\ell^2\,\to\,\ell^2} \le \snorm{B}{0}{0}.
\end{equation}
\end{coro}

\begin{defi}
For $b\in\BS^\infty$, we define
\begin{equation}\label{eq:adjoint symbol}
b^\dagger_\theta(\xi):=
\begin{cases}\overline{b_{\theta^{-1}}(\theta\act\xi)}&
\text{if } \theta\in \Theta,\\
0& \text{if }\theta\in G\setminus\Theta
\end{cases}
\end{equation}
for all $\xi\in\Xi$, where $\Theta$ is a frequency set for $b$.
\end{defi}

\begin{lemm} \label{lem:adjoint}
If $b\in\BS^\gamma$, then $b^\dagger\in\BS^\gamma$. Moreover, for all $x,y\in\sfH^\gamma$, one has
\begin{equation}\label{eq:B+onH}
(x,By)_{\ell^2(\Xi)}=\left(B^\dagger x,y\right)_{\ell^2(\Xi)}, \quad \text{i.e.}\; B^\dagger\subseteq B^*.
\end{equation}
In particular, $B$ is symmetric on $\sfH^\gamma$ if and only if $B=B^\dagger$.
\end{lemm}

\begin{proof}
Every frequency set $\Theta$ for $b\in\BS^\gamma$ is also one for $b^\dagger$. Moreover, since $\Theta=\Theta^{-1}$ holds by convention, \eqref{eq:anginv} and~\eqref{eq:peetre} imply that for all $l\geq 0$,
\begin{equation*}
\begin{aligned}
\snorm{b^\dagger}{\gamma}{l}&=\sum\limits_{\theta\,\in\,\Theta_B} \ang{\theta}^l\sup\limits_{\xi\,\in\,\Xi}\big[\ang{\xi}^{-\gamma}|b_{\theta^{-1}}(\theta\act\xi)|\big]\\
&=\sum\limits_{\theta\,\in\,\Theta_B}\ang{\theta}^l\sup\limits_{\xi\,\in\,\Xi}\big[\left\langle\theta^{-1}\act\xi\right\rangle^{-\gamma}|b_{\theta^{-1}}(\xi)|\big]\\
&\leq \sum\limits_{\theta\,\in\,\Theta_B}\ang{\theta}^{l+|\gamma|}\sup\limits_{\xi\,\in\,\Xi}\big[\ang{\xi}^{-\gamma}|b_{\theta^{-1}}(\xi)|\big]\\
&=\sum\limits_{\theta\,\in\,\Theta_B}\ang{\theta}^{l+|\gamma|}\sup\limits_{\xi\,\in\,\Xi}\big[\ang{\xi}^{-\gamma}|b_{\theta}(\xi)|\big]\\
&=\snorm{b}{\gamma}{l+|\gamma|},
\end{aligned}
\end{equation*}
thus $b^\dagger\in\BS^\gamma$ holds. Moreover, \eqref{eq:symbrecovery} and~\eqref{eq:adjoint symbol} yield
\begin{equation}\label{eq:B+onspanE}
(\be_\eta,B\be_\xi)=\left(B^\dagger\be_\eta,\be_\xi\right), \qquad \text{for all }\eta,\xi\in\Xi.
\end{equation}
In view of Lemma~\ref{lem:domain} and the density of $\spann(\CE)$ in $\sfH^\gamma$, \eqref{eq:B+onspanE} extends to~\eqref{eq:B+onH}. This finishes the proof of the lemma.
\end{proof}

\begin{defi}
Let $a, b \in \BS^\infty$ be symbols with frequency sets $\Theta_a$ and $\Theta_b$. The \emph{composed symbol} $a \circ b$ with frequency set
\begin{align}\label{eq:frequencysetcomposition}
\Theta_{a\circ b}:=\Theta_a \Theta_b := \set{\theta_a \theta_b :
\theta_a \in \Theta_a, \theta_b \in \Theta_b}
\end{align}
is defined as
\begin{equation}\label{eq:symbolcomposition}
(a \circ b)_\theta(\xi) := \sum_{\theta_a \theta_b = \theta} a_{\theta_a}(\theta_b \act\xi) b_{\theta_b}(\xi) \quad\text{ for all }\theta\in\Theta_{a\circ b}, \xi\in\Xi.
\end{equation}
\end{defi}


\begin{lemm}\label{lem:product}
For $\alpha, \beta\in \R$ let $A = \Op(a) \in \BS^{\alpha}$ and $B = \Op(b) \in \BS^{\beta}$. Then $AB\in\BS^{\alpha+\beta}$ and $AB =\Op(a\circ b)$. Moreover, for all $l \ge 0$ we have the bound
\begin{equation}\label{eq:product}
\snorm{AB}{\alpha+\beta}{l} \leq
\snorm{A}{\alpha}{l}\snorm{B}{\beta}{l + \abs{\alpha}}.
\end{equation}
\end{lemm}

\begin{proof}
The frequency set $\Theta_{a\circ b}$ is, clearly, a countable set. For any $l \ge 0$, we have
\begin{equation}
\begin{aligned}
\snorm{a\circ b}{\alpha + \beta}{l} &= \sum_{\theta\,\in\,\Theta_{a\circ b}}
\,\sum_{\theta_a\theta_b = \theta} \ang{\theta}^l
\sup_{\xi\,\in\,\Xi} \left(\ang{\xi}^{-\alpha - \beta}
\abs{a_{\theta_a}(\theta_b\act\xi)}\abs{b_{\theta_b}(\xi)}\right) \\
&\le \sum_{\theta_b\,\in\,\Theta_b} \ang{\theta_b}^{l + \abs{\alpha}}
\sup_{\xi\,\in\,\Xi}\left(\ang{\xi}^{-\beta}
\abs{b_{\theta_b}(\xi)}\right)\\
&\qquad \times \sum_{\theta_a\,\in\,\Theta_a}\ang{\theta_a}^{l}
\sup_{\zeta\,\in\,\Xi}\left(\ang{\theta_b \act \zeta}^{-\alpha}
\abs{a_{\theta_a}(\theta_b \act \zeta)}\right) \\
&\le \snorm{a}{\alpha}{l} \snorm{b}{\beta}{l + \abs{\alpha}}.
\end{aligned}
\end{equation}
Thus $a\circ b\in \BS^{\alpha + \beta}$ and~\eqref{eq:defopB} implies $AB =\Op(a\circ b)$.
\end{proof}

It is natural to consider operators from $\BS^\infty$ on the common domain $\sfH^\infty$. Then Lemmata~\ref{lem:domain}, \ref{lem:adjoint}, and~\ref{lem:product} yield the following corollary.
\begin{coro}
$\BS^\infty=\bigcup_{\gamma\,\in\,\R}\BS^\gamma$ is a $*$-algebra of operators on $\sfH^\infty$, filtered by $\R$, with involution $\dagger$. The subalgebra of regularising operators $\BS^{-\infty}$ forms a two-sided ideal of $\BS^\infty$.
\end{coro}

We also consider the adjoint actions $\ad(A;B): = \rmi(AB -BA)$ with the frequency set $\Theta_{\ad(a;b)} = \Theta_{a \circ b } \cup \Theta_{b \circ a}$. The Fourier coefficients of $\ad(A;B)$ are
\begin{equation}\label{eq:ad symbol}
\ad(a;b)_\theta(\xi) = \rmi\, \left(\sum_{\theta_a \theta_b = \theta} a_{\theta_a}(\theta_b\act\xi) b_{\theta_b}(\xi) -
\sum_{\theta_b \theta_a = \theta} b_{\theta_b}(\theta_a\act\xi)a_{\theta_a}(\xi)\right),
\end{equation}
for all $\theta \in \Theta_{\ad(a;b)}$. If $G$ is commutative, \eqref{eq:ad symbol} simplifies to
\begin{equation}\label{eq:adcom}
\ad(a;b)_\theta(\xi) = \rmi\, \sum_{\theta_a \theta_b = \theta}\big(a_{\theta_a}(\theta_b\act\xi) b_{\theta_b}(\xi) - b_{\theta_b}(\theta_a\act\xi)a_{\theta_a}(\xi)\big).
\end{equation}

For $k=1,2,3,\,\dots$ and $A,B,B_1,\,\dots B_k\in\BS^\infty$, we define recursively \label{page:defad}
\begin{equation}\label{eq:defadk}
\begin{aligned}
\ad(A;B_1,\,\dotsc,\,B_k) &:= \ad\left(\ad(A;B_1,\dotsc,B_{k-1});B_k\right),\\
\ad^0(A;B)&:= A,\\
\ad^k(A;B)&:= \ad\left(\ad^{k -1}(A;B);B\right).
\end{aligned}
\end{equation}

The following lemma is a direct consequence of Lemma~\ref{lem:product}.
\begin{lemm}\label{lem:weakcommest}
Let $k \in \N$ and assume that $A_j \in \BS^{\gamma_j}$ for $0 \le j \le k$. Put
\begin{equation}
\gamma = \sum_{j=0}^k \gamma_j, \qquad \hat\gamma =\sum_{j=0}^k |\gamma_j|.
\end{equation}
Then $\ad(A_0;A_1,\,\dotsc,\,A_k) \in \BS^{\gamma}$. Furthermore, if for all $0 \le j\le k$ we have $A_j = A_j^\dagger$, then $\ad(A_0;A_1,\,\dotsc,\,A_k)=\ad(A_0;A_1,\,\dotsc,\,A_k)^\dagger$. Moreover, for all $l
\ge 0$ we have
\begin{equation}
\begin{aligned}\label{eq:normestcomm}
\snorm{\ad(A_0;A_1,\,\dotsc,\,A_k)}{\gamma}{l}\leq 2^{k}
\prod_{j=0}^k \snorm{A_j}{\gamma_{j}}{l + \hat\gamma - |\gamma_j|}.
\end{aligned}
\end{equation}
In particular, for any $A\in\BS^{\alpha}$, $B\in\BS^{0}$ and $k\in\N$ we obtain the estimate
\begin{align}\label{eq:normestcommzeroorder}
\snorm{\ad^k(A;B)}{\alpha}{l}\leq 2^k \snorm{A}{\alpha}{l}
\left(\snorm{B}{0}{l + \abs{\alpha}}\right)^k.
\end{align}
\end{lemm}
For some $\Xi$ and $G$ it may be possible to improve this lemma and show that $\ad(A;B)\in\BS^\gamma$ holds with $\gamma<\alpha+\beta$ for all $A\in\BS^{\alpha}$, $B\in\BS^{\beta}$. This will be discussed in Section~\ref{sec:sgt}.

The following proposition provides bounds on norms of operators restricted to ``annuli'' in $\Xi$.

\begin{prop}\label{prop:normorder}
For $1 \le m \le M \le \infty$, let $\Upsilon \subseteq \set{\xi \in \Xi : m \le \ang \xi \le M}$ and denote by $P_{\Upsilon}$ the orthogonal projection in $\ell^2(\Xi)$ onto the closure of $\spann \{\be_\xi :$ \linebreak$\xi \in \Upsilon\}$. Then, for any $A \in \BS^\gamma$ with $\gamma \ge 0$, the norm inequality
\begin{equation}\label{eq:upslarge}
\norm{A P_\Upsilon}_{\ell^2\,\to\,\ell^2} \le M^\gamma\snorm{A}{\gamma}{0}
\end{equation}
holds. For any $A \in \BS^\gamma$ with $\gamma \le 0$, we get the inequality
\begin{equation}
\norm{A P_\Upsilon}_{\ell^2\,\to\,\ell^2} \le m^\gamma\snorm{A}{\gamma}{0}.
\end{equation}
\end{prop}
\begin{proof}
Observe that $P_\Upsilon$ is a quasi-periodic operator with a frequency set $\Theta=\{\id\}$ and the symbol $(p_\Upsilon)_{\id} =\bone_{\Upsilon}$ (the indicator function of $\Upsilon$). Thus, for all $\gamma\in\R$ and $l\geq0$,
\begin{equation}\label{eq:upscases}
\begin{aligned}
\snorm{P_\Upsilon}{-\gamma}{l} &= \sup_{\xi \,\in\,\Upsilon} \ang \xi^\gamma \leq
\begin{cases}
\hfill m^\gamma &\text{if } \gamma \le 0, \\
\hfill M^\gamma &\text{if } \gamma \ge 0.
\end{cases}
\end{aligned}
\end{equation}
If $M < \infty$ or $\gamma \le 0$, then Corollary~\ref{cor:domain} and Lemma~\ref{lem:product} imply the bound
\begin{equation}
\norm{A P_\Upsilon}_{\ell^2\,\to\,\ell^2} \le \snorm{A}{\gamma}{0} \snorm{P_\Upsilon}{-\gamma}{\abs{\gamma}},
\end{equation}
and the statement of the lemma follows from~\eqref{eq:upscases}. On the other hand, the inequality~\eqref{eq:upslarge} is trivial for $M = \infty$ and $\gamma > 0$.
\end{proof}


\section{Elliptic and diagonal operators}\label{sec:perturbation}

In this section, we introduce particular classes of operators from $\BS^\infty$ and study their properties. Some of these classes do depend on the specific choice of orthonormal basis $\CE$ for $\ell^2(\Xi)$. However, the class of operators on which our main theorems depend, that of elliptic operators, is invariant under change of basis.

\begin{defi}\label{def:dapo}
The subalgebra $\BD\BS^\infty \subset \BS^\infty$ of \emph{diagonal operators} is defined as
\begin{equation}
\BD\BS^\infty := \set{A = \Op(a) \in \BS^\infty:\set{\id} \text{ is a frequency set for } a}.\label{eq:dapo}
\end{equation}
For symbols of operators from $\BD\BS^\infty$ we can suppress the subscript $\id$, i.e. we let $a(\xi) :=a_{\id}(\xi)$ for all $A =\Op(a)\in\BD\BS^\infty$, $\xi\in\Xi$. For $\alpha \in \R\cup \set{-\infty}$ we define $\BD\BS^\alpha := \BD\BS^\infty \cap \BS^\alpha$. Introduce the map $\CD : \BS^\infty \to \BD\BS^\infty$, $A\mapsto A^\CD$, that projects $A=\Op(a)$ onto its diagonal part $A^\CD:= \Op(a^\CD)$ where
\begin{equation}
a^\CD(\xi) := a_{\id}(\xi),
\end{equation}
i.e.
\begin{equation}
A^\CD\be_{\xi} =\langle\be_{\xi}, A\be_{\xi}\rangle \be_{\xi}
\end{equation}
holds for all $\xi \in\Xi$. We also define the \emph{off-diagonal part} as $A^{\CO\CD} := \Op(a^{\CO\CD})$ with $a^{\CO\CD}:=a-a^\CD$.
\end{defi}

Note that for any $A\in\BS^\alpha$ with $\alpha\in\R$ and all $l\geq 0$,
\begin{equation}\label{eq:partition}
\snorm{A^{\CD}}{\alpha}{l}+\snorm{A^{\COD}}{\alpha}{l}=\snorm{A}{\alpha}{l}
\end{equation}
and
\begin{equation}\label{eq:diagonal l irrelevant}
\snorm{A^\CD}{\alpha}{l} =\snorm{A^\CD}{\alpha}{0}.
\end{equation}

\begin{defi}\label{def:DES}
The set $\BD\BE\BS^\alpha$ of \emph{diagonal elliptic operators of order $\alpha\in\R$} is defined as the set of operators $A = \Op(a) \in \BD\BS^\alpha$ for which there exist \emph{ellipticity parameters} $\kappa> 0$ and $\er\ge 1$ such that
\begin{equation}\label{eq:ellconst}
|a(\xi)|\geq \kappa\ang{\xi}^\alpha \qquad \text{ for all $\xi\in\Xi$ such that }\ang{\xi}\geq
\er.
\end{equation}
Let the set of ellipticity parameters $(\kappa, \er)$ of $A$ be denoted by $\FE(A)$. Note that $(\kappa, \er)\in \FE(A)$ implies $(\tilde \kappa,\widetilde
\er)\in \FE(A)$ for all $0 <\tilde \kappa\leq \kappa$, and $\widetilde \er \geq \!
\er$.
\end{defi}

\begin{defi}\label{def:ES}
The set $\BS\BE\BS^\alpha$ of \emph{strongly elliptic operators of order $\alpha\in\R$} consists of operators $A\in\BS^\alpha$ such that $A^\CD \in \BD\BE\BS^\alpha$ and $A^{\CO\CD} \in
\BS^\gamma$ for some $\gamma < \alpha$. For $(\kappa,\er)\in \FE(A^\CD)$ we define $P_\er$ as the diagonal operator with symbol $\bone_{\{\xi: \ang{\xi}\leq
\er\}}$. We also define $P_\er^c$ as $\Id - P_\er$, and
\begin{equation}\label{eq:tilde A}
\tilde A_{\kappa,\er} := A^\CD P_\er^c + \kappa \er^\alpha P_\er.
\end{equation}
\end{defi}

\begin{defi}\label{def:WES}
The set $\BE\BS^\alpha$ of \emph{elliptic operators of order $\alpha\in\R$} consists of operators $A\in\BS^\alpha$ for which there exists a unitary $U\in\BS^0$ with $U A U^\dagger\in\BS\BE\BS^\alpha$.
\end{defi}

Clearly, both $\BS\BE\BS^\alpha$ and $\BE\BS^\alpha$ are closed under addition of operators in $\BS^\beta$, $\beta < \alpha$. As we did with diagonal operators, we set
\begin{equation}
\BT^\infty:=\bigcup\limits_{\gamma\,\in\,\R}\BT^\gamma, \qquad \BT^{-\infty}:=\bigcap\limits_{\gamma\,\in\,\R}\BT^\gamma, \ \text{for} \ \BT\in \lbrace \BD\BE\BS, \BS\BE\BS, \BE\BS \rbrace.
\end{equation}

\begin{prop}\label{prop:globell}
Let $A\in\BS^\infty$ and $\alpha >0$ such that $A^{\CD}\in\BD\BE\BS^\alpha$. For any $(\kappa, \er)\in\FE(A^\CD)$ the operator $\tilde A_{\kappa,\er}$ is invertible with $\tilde A_{\kappa,\er}^{-1}\in\BD\BS^{-\alpha}$ and for all $l \ge 0$ we have
\begin{equation}\label{eq:globinvsnorm}
\snorm{\tilde A_{\kappa,\er}^{-1}}{\gamma}{l} =\snorm{\tilde A_{\kappa,\er}^{-1}}{\gamma}{0} \leq \kappa^{-1}
\begin{cases}
\er^{-\alpha} &\text{for } \gamma \geq 0,\\ \er^{-\alpha -\gamma} &\text{for }-\alpha \leq
\gamma < 0.
\end{cases}
\end{equation}
Moreover, the following estimates hold for all $\gamma \in\R$, $l\geq 0$:
\begin{align}
\snorm{A \tilde A_{\kappa,\er}^{-1} - \Id}{\gamma - \alpha}{l} &\le
\er^{\alpha -\gamma} + \frac 1 \kappa \left(
\er^{\alpha - \gamma} \snorm{A^\CD}{\alpha}{0} +
\snorm{A^{\CO\CD}}{\gamma}{l}
\right)\label{eq:AA_R -Id}
\\
\intertext{and}
\snorm{\tilde A_{\kappa,\er}^{-1} A - \Id}{\gamma - \alpha}{l} &\le
\er^{\alpha -\gamma} + \frac 1 \kappa \left(
\er^{\alpha - \gamma} \snorm{A^\CD}{\alpha}{0} +
\snorm{A^{\CO\CD}}{\gamma}{l+\alpha}
\right).\label{eq:A_RA -Id}
\end{align}
\end{prop}

\begin{proof}
We have that $\tilde A_{\kappa,\er} - A^{\CD} = (\kappa \er^\alpha - A^{\CD}) P_\er$, and since $P_\er \in \BS^{-\infty}$, which is an ideal of $\BS^\infty$, we observe that $\tilde A_{\kappa,\er} \equiv A^{\CD}
\mod{\BS^{-\infty}}$. By~\eqref{eq:ellconst} and~\eqref{eq:tilde A}, $\tilde A_{\kappa,\er}\in\BD\BS^\infty$ and its symbol satisfies
\begin{equation}
\abs{\tilde a_{\kappa, \er}(\xi)} = \abs{a_{\id}(\xi)}\bone_{\set{\ang \xi >\er}} +
\kappa \er^\alpha\bone_{\set{\ang \xi\le\er}} \ge \kappa \ang
\xi^\alpha\bone_{\set{\ang \xi>\er}} + \kappa \er^\alpha\bone_{\set{\ang \xi\le\er}}
\end{equation}
for all $\xi \in \Xi$. Hence $\tilde A_{\kappa,\er}^{-1} =\Op\big(\tilde a_{\kappa,\er}^{-1}\big)\in\BD\BS^{-\alpha}$ and~\eqref{eq:globinvsnorm} holds.

The estimates~\eqref{eq:AA_R -Id} and~\eqref{eq:A_RA -Id} follow by applying~\eqref{eq:product} term-wise to the right hand sides of the identities
\begin{equation*}
\begin{split}
A \tilde A_{\kappa,\er}^{-1} - \Id &= -P_\er +\kappa^{-1}\er^{-\alpha}A^{\CD}P_\er +A^{\CO\CD}\tilde A_{\kappa,\er}^{-1},\\
\tilde A_{\kappa,\er}^{-1}\,A - \Id &= -P_\er +\kappa^{-1}\er^{-\alpha}A^{\CD}P_\er +\tilde A_{\kappa,\er}^{-1}\,A^{\CO\CD}
\end{split}
\end{equation*}
and taking~\eqref{eq:diagonal l irrelevant} into account.
\end{proof}


Our current goal is to understand perturbations of strongly elliptic operators of positive order. In particular, we show that operators of lower order are relatively bounded with respect to them.

\begin{lemm}\label{lem:relative boundedness}
Let $\beta\in\R$, $\alpha>\max(\beta,0)$, $0 <\gamma < \alpha$, and assume that $A \in \BS\BE\BS^\alpha$ with $A^{\CO\CD}\in \BS^\gamma$ and $B \in \BS^\beta$. Then for $\beta\le 0$ the operator $B$ is bounded, and, in the case of $\alpha
>\beta >0$, for every $x\in\sfH^\alpha$ and
\[
(\kappa,
\er)\in\FE\left(A^\CD\right)\cap \left\{\er
\geq\left(\left\|A^{\CO\CD}\right\|_0^{(\gamma)}/\kappa\right)^{1/(\alpha
-\gamma)}\right\},
\]
we have
\begin{equation}\label{eq:infinitesimal estimate}
\|Bx\| \le \frac{\er^{\beta -\alpha}\|B\|_0^{(\beta)}}{\kappa -\er^{\gamma -\alpha}\|A^{\CO\CD}\|_0^{(\gamma)}}\Big(\|Ax\| +\kappa \er^{\alpha}\left(1 +\kappa^{-1}\left\|A^\CD\right\|_0^{(\alpha)}\right)\|x\|\Big).
\end{equation}
In particular, $B$ is infinitesimally $A$-bounded in $\ell^2(\Xi)$.
\end{lemm}

\begin{proof}
The only non-trivial case is $\alpha >\beta >0$. For every $x\in\sfH^\alpha$ we have
\begin{equation*}
\|Bx\| \le \left\|B\tilde A_{\kappa,\er}^{-1}\right\|\norm{A^\CD x} +\left\|B\left(\tilde A_{\kappa,\er}^{-1}A^\CD -\Id\right)\right\|\norm{x},
\end{equation*}
with $\tilde A_{\kappa,\er}$ defined as in~\eqref{eq:tilde A}. Corollary~\ref{cor:domain} and displays~\eqref{eq:product}, \eqref{eq:globinvsnorm} and~\eqref{eq:A_RA -Id} imply the estimates
\begin{align*}
\left\|B\left(\tilde A_{\kappa,\er}^{-1}A^\CD -\Id\right)\right\| &\le
\|B\|_0^{(\beta)}\left\|\tilde A_{\kappa,\er}^{-1}A^\CD -\Id\right\|_{|\beta|}^{(-\beta)}\\
&\le \er^\beta\|B\|_0^{(\beta)}\big(1 +\kappa^{-1}\left\|A^\CD\right\|_0^{(\alpha)}\big)
\\
\intertext{and}
\left\|B\tilde A_{\kappa,\er}^{-1}\right\| &\le \|B\|_0^{(\beta)}\left\|\tilde A_{\kappa,\er}^{-1}\right\|_0^{(-\beta)}\\ &\le \kappa^{-1}\|B\|_0^{(\beta)}\er^{\beta -\alpha},
\end{align*}
and we obtain
\begin{equation}\label{eq:prelim infinitesimal estimate}
\|Bx\| \le \kappa^{-1}\er^{\beta -\alpha}\|B\|_0^{(\beta)}\Big(\left\|A^\CD x\right\| +\kappa \er^\alpha\left(1 +\kappa^{-1}\left\|A^\CD\right\|_0^{(\alpha)}\right)\|x\|\Big),
\end{equation}
which is~\eqref{eq:infinitesimal estimate} with $A^\CD\in \BD\BE\BS^\alpha$ replacing $A$. Applying~\eqref{eq:prelim infinitesimal estimate} with $B=A^{\CO\CD}$, we arrive at
\begin{equation}
\begin{aligned}
\left\|A^{\CO\CD}x\right\| &\le \kappa^{-1}\er^{\gamma -\alpha}\left\|A^{\CO\CD}\right\|_0^{(\gamma)}\left\|A^\CD x\right\|  \\ &\qquad +\er^\gamma\left\|A^{\CO\CD}\right\|_0^{(\gamma)}\left(1 +\kappa^{-1}\left\|A^\CD\right\|_0^{(\alpha)}\right)\|x\|.
\end{aligned}
\end{equation}
Hence we have
\begin{align*}
\|Ax\| &\ge \left\|A^\CD x\right\| -\left\|A^{\CO\CD}x\right\|\\
&\ge \left(1 -\kappa^{-1}\er^{\gamma -\alpha}\left\|A^{\CO\CD}\right\|_0^{(\gamma)}\right)\left\|A^\CD x\right\| -\er^\gamma\left\|A^{\CO\CD}\right\|_0^{(\gamma)}\left(1 +\kappa^{-1}\left\|A^\CD\right\|_0^{(\alpha)}\right)\|x\|,
\end{align*}
which implies
\begin{multline}\label{eq:A^CD control}
\left\|A^\CD x\right\|\\
\le \left(1 -\kappa^{-1}\er^{\gamma -\alpha}\left\|A^{\CO\CD}\right\|_0^{(\gamma)}\right)^{-1} \Big(\|A x\| +\er^\gamma\left\|A^{\CO\CD}\right\|_0^{(\gamma)}\left(1 +\kappa^{-1}\left\|A^\CD\right\|_0^{(\alpha)}\right)\|x\|\Big).
\end{multline}
Substituting~\eqref{eq:A^CD control} into~\eqref{eq:prelim infinitesimal estimate} we obtain~\eqref{eq:infinitesimal estimate}.
\end{proof}


We conclude the section with the following proposition.

\begin{prop}\label{prop:closedness and self-adjointness}
For $\alpha \in\R$, every operator from $\BE\BS^\alpha$ is closed on $\sfH^{\max\{\alpha, 0\}}$ in the Hilbert space $\ell^2(\Xi)$. Every symmetric operator from $\BE\BS^\alpha$ defined on $\sfH^{\max\{\alpha, 0\}}$ is self-adjoint.
\end{prop}
\begin{proof}
For $\alpha\leq 0$ we have $\BE\BS^\alpha\subset \BS^0$, and the statements follow from Corollary~\ref{cor:domain}. Now assume $A\in \BS\BE\BS^\alpha$ with $\alpha >0$. By~\eqref{eq:inner product}, Definition~\ref{def:ES} and Lemma~\ref{lem:domain}, for any $(\kappa,
\er)\in\FE(A^\CD)$ we have the estimates
\begin{equation}
\kappa^2\|x\|_{\sfH^\alpha}^2 \le \left\|\tilde A_{\kappa,\er}x\right\|^2 \le
\left(\left\|\tilde A_{\kappa,\er}\right\|^{(\alpha)}_{\alpha}\right)^2\|x\|_{\sfH^\alpha}^2
\end{equation}
for all $x \in \RH^\alpha$. Hence the graph norm of $\tilde A_{\kappa,\er}$ is equivalent to the norm of $\sfH^\alpha$, and $\tilde A_{\kappa,\er}$ is closed on $\sfH^\alpha$. If $\tilde A_{\kappa,\er}$ is symmetric, then for every $x\in\dom(\tilde A_{\kappa,\er}^*)$ there exists $C_x\geq 0$ such that for all $y \in \RH^\alpha$
\begin{equation}
\left|\left(x, \tilde A_{\kappa,\er}y\right)\right|\leq C_x\|y\|_{\ell^2(\Xi)}.
\end{equation}
In particular, with $(y_{n})_\xi:=\overline{(\tilde a)_{\kappa, \er, \id}(\xi)}\bone_{\ang
\xi\,\le\,n}\ x_\xi$ for $n\geq \er$, $\xi\in \Xi$, we obtain by~\eqref{eq:ellconst} and~\eqref{eq:tilde A} that
\begin{equation}
\begin{split}
\sum_{\substack{\xi\,\in\,\Xi \\
\ang \xi\,\le\,n}}\ang{\xi}^{2\alpha}|x_\xi|^2 &\leq \kappa^{-1}\left(x, \tilde A_{\kappa,\er}y_{n}\right)\leq
\kappa^{-1}C_x\|y_n\|_{\ell^2(\Xi)}\\
&\leq \kappa^{-1}C_x\snorm{\tilde{A}_{\kappa,\er}}{\alpha}{0}\left(\sum_{\substack{\xi\,\in\,\Xi \\ \ang \xi\,\le\,n}}\ang{\xi}^{2\alpha}|x_\xi|^2\right)^{1/2}.
\end{split}
\end{equation}
Passing to the limit $n\to \infty$, it follows from~\eqref{eq:inner product} that $\|x\|_{\sfH^\alpha} \leq
\kappa^{-1}C_x\|\tilde{A}_{\kappa,\er}\|^{(\alpha)}_{0}$, i.e.
\[
\dom\left(\tilde
A_{\kappa,\er}^*\right)\subset \sfH^\alpha =\dom\left(\tilde A_{\kappa,\er}\right),
\]
hence
$\tilde A_{\kappa,\er}$ is self-adjoint. By Lemma~\ref{lem:relative boundedness} $A -\tilde A_{\kappa,\er}$ is infinitesimally $\tilde A_{\kappa,\er}$-bounded, so that $A$ is also self-adjoint (see, e.g., Theorems~3.4.2 and 4.1.9 in~\cite{BirmanSolomyak}).

For $A\in\BE\BS^\alpha$, by Definition~\ref{def:WES} there exist a unitary $U\in\BS^0$ and $H\in\BS\BE\BS^\alpha$ such that $A=UH U^\dagger$. Moreover, it follows from Lemma~\ref{lem:domain} that $U\sfH^\alpha=U^\dagger\sfH^\alpha=\sfH^\alpha$. Now, let $(x_n)_{n\,\in\,\N}\subset\sfH^\alpha$ with $x_n\to x$ and $U^\dagger HUx_n\to z$ in $\ell^2$, as $n\to\infty$. Since $U$ is bounded, $Ux_n\to Ux$ and $HUx_n\to Uz$, thus the closedness of $H$ implies that $Ux\in \sfH^\alpha$ and $HUx=Uz$, i.e.~$x\in\sfH^\alpha$ and $U H U^\dagger x=z$. Hence, $A$ is closed on $\sfH^\alpha$ and self-adjoint if symmetric.
\end{proof}

\section{The Density of States Measure and von Neumann Algebras} \label{sec:ids}


In this section, following~\cite{Shubin1979a}, we consider a representation of $\BS^\infty$ into another operator algebra, affiliated with an infinite factor (accounting for the almost periodicity), and define the density of states measure (DSM) for self-adjoint operators in $\BE\BS^\infty$ with respect to this representation. For a suitable representation, this DSM will coincide with the classically defined DSM on elliptic differential operators with almost periodic coefficients. We follow the construction and terminology of~\cite[\S 1]{Shubin1979a}, generalising Shubin's symbol classes to the ones defined in Section~\ref{sec:operators}.

\subsection{Representations of the operator algebra}

Before stating the abstract conditions we assume on the algebra of almost-periodic operators on $\Xi$, let us recall a few definitions, which can be found in~\cite[\S 34--38]{Naimark}.
\begin{defi}\label{def:vNa}
Let $\FH$ be a Hilbert space and $\CB(\FH)$ be the algebra of bounded linear operators on $\FH$. For a subalgebra $\FA \subset \CB(\FH)$, its
\emph{commutant} is defined as
\begin{equation}
\FA' := \set{B \in \CB(\FH) : AB = BA \text{ for all } A \in \FA}.
\end{equation}
We say that $\FA \subset \CB(\FH)$ is a \emph{factor} if $\FA \cap \FA' =
\spann(\Id_\FH)$. A family of densely defined, not necessarily bounded, operators $\FS$ with domains in $\FH$ is said to be \emph{affiliated to} $\FA$ if every $B \in \FS$ commutes with every unitary $U \in \FA'$, this relationship is denoted $\FS \eta \FA$. Similarly, a subspace $\FK \subset
\FH$ is affiliated to $\FA$, again denoted $\FK \eta \FA$ if it is invariant under the action of every unitary $U \in \FA'$. If $\FK$ is closed, this is readily seen to be equivalent to the projection $P_{\FK}$ being in $\FA$.
\end{defi}

Let us now set up some notation for the rest of this section. We set $\FH$ as some Hilbert space, and $\FA$ as a factor of either type $\rmI_\infty$ or $\rmII_\infty$ in $\CB(\FH)$, the algebra of bounded linear operators in $\FH$. The precise definition of factor type is not relevant to us, we only use the fact that they carry a well-defined notion of trace, see the beginning of Section~\ref{sec:dsm}.

Let $\widetilde{\sfH}^\infty$ be a dense subspace of $\FH$ and
\[
\widetilde{\BS}^\infty =\bigcup_{\gamma\,\in\,\R}\widetilde\BS^\gamma
\]
be a
$*$-algebra of unbounded linear operators in $\FH$ defined on $\widetilde{\sfH}^\infty$, filtered by $\R$. We assume that $\widetilde{\BS}^\infty\widetilde{\sfH}^\infty\subseteq\widetilde{\sfH}^\infty$, and that $\widetilde{\BS}^\infty$ is invariant under the involution $\tilde A\mapsto\tilde{A}^\dagger:=\big.\tilde{A}^*|_{\widetilde{\sfH}^\infty}$, where $\tilde{A}^*$ is the adjoint to $\tilde A\in\widetilde{\BS}^\infty$. We also suppose that $\widetilde{\BS}^\infty$ is affiliated with the factor $\FA$. Finally, we assume there is a representation $\rho: \BS^\infty \to \widetilde{\BS}^\infty$ having the following properties:
\begin{enumerate}\romanenumi
\item \label{item:rhohomo}
$\rho$ is a homomorphism of filtered $*$-algebras with $\rho(\BS^\gamma)\subseteq \widetilde\BS^\gamma$, for all $\gamma\in\R$.
\item \label{item:rhoS0bounded}
For every $A\in\BS^0$, $\rho(A)$ extends to a bounded linear operator on $\FH$ with
\begin{equation}\label{eq:normpreserving}
\norm{\rho(A)}_{\FH\,\to\,\FH} = \norm{A}_{\mathsf{\ell}^2(\Xi)\,\to\,\mathsf{\ell}^2(\Xi)}.
\end{equation}
\item For all $A\in\BS^\infty$, $\rho(A)$ is closable in $\FH$ with the closure $A^\sharp:=\overline{\rho(A)}$. For every $\alpha>0$ there exists a dense subspace $\widetilde{\sfH}^\alpha\supseteq\widetilde{\sfH}^\infty$ such that
\begin{enumerate}
\item $\widetilde{\sfH}^\alpha\subseteq\widetilde{\sfH}^\gamma$ if $0 <\gamma\leq \alpha$,
\item \label{item:domsharpDES}
$A\in\BD\BE\BS^\alpha$ implies $\dom(A^\sharp)=\widetilde{\sfH}^\alpha$,
\item \label{item:domsharpinv}
for all $B\in\BS^0$, $B^\sharp \widetilde{\sfH}^\alpha\subseteq\widetilde{\sfH}^\alpha$.
\end{enumerate}
\item \label{item:sharpAselfadjoint}
If $A\in\BD\BE\BS^\alpha$, $\alpha>0$, is self-adjoint on $\sfH^\alpha$, then $A^\sharp$ is self-adjoint.
\end{enumerate}

\begin{rema}
When $\FA$ is a $\rmI_\infty$ factor some of the statements in this section become rather trivial. However, we include this case for applications in Section~\ref{sec:bs}.
\end{rema}
\begin{rema}\label{rem:Shubin}
In~\cite{Shubin1979a}, Shubin considers $G = \R^d$ acting on itself by translation, with almost periodic operators acting both in Besicovitch space $\sfB^2(\R^d)\cong \ell^2(\R^d)$ and in $\sfL^2(\R^d)$ through the Fourier integral representation of pseudo-differential operators. The appropriate Hilbert space is then
\begin{equation}\label{eq:hilbert}
\FH = \sfB^2\left(\R^d\right) \otimes \sfL^2\left(\R^d\right),
\end{equation}
and the $\rmII_\infty$ factor $\FA$ is generated by the two families of operators
\begin{equation}
\set{\be_{\bxi} \otimes \be_{\bxi}: \bxi \in \R^d} \text{ and } \set{I \otimes T_{\bxi}: \bxi \in \R^d},\label{eq:iiinffact}
\end{equation}
where $\be_{\bxi}$ is multiplication by the character $\be_{\bxi}(\bx)=\rme^{\rmi\bxi\cdot \bx}$ and $T_{\bxi}$ is the translation operator $T_{\bxi} f(\bx) = f(\bx - \bxi) $. The representation $\rho$ is given on $A=\Op(a)\in \BS^\infty$ by the linear operator $\rho(A)=a(\bx+\by;D_\by)$ acting on
\begin{equation}
\widetilde{\sfH}^\infty:=\sfB^2\left(\R^d\right) \otimes \hat{\sfH}^\infty\left(\R^d\right).
\end{equation}
Here, $\bx$ is the variable of functions in $\sfB^2(\R^d)$, $\by$ is the variable of functions in $\RL^2(\R^d)$, $D_y=-\rmi\nabla_\by$, and $\hat{\sfH}^\infty(\R^d):= \{f\in \sfC^\infty(\R^d): \partial^\alpha f\in\sfL^2(\R^d)\text{ for all }\alpha\in\N_0^d\}$.
\end{rema} 
Properties~\eqref{item:rhohomo} and~\eqref{item:rhoS0bounded} of the representation $\rho$ imply the following lemma.
\begin{lemm}\label{lem:sharpS0}
If $A\in\BS^0$, then $A^\sharp$ is defined on $\FH$ and satisfies 
\[
\big(A^\sharp\big)^*=\big(A^\dagger\big)^\sharp = \big(A^*\big)^\sharp\quad\text{and}\quad\left\|A^\sharp\right\|_{\FH\,\to\,\FH}=\|A\|_{\ell^2(\Xi)\,\to\,\ell^2(\Xi)}.
\]
In particular, the map $\BS^0\to\FB(\FH),\ A\mapsto A^\sharp$ is an injective homomorphism of $*$-algebras. If $U\in\BS^0$ is unitary, then so is $U^\sharp$.
\end{lemm}
We will now carry over Lemma~\ref{lem:relative boundedness} to images under $\sharp$. This provides us with some information on the domains of operators from $(\BS^\infty)^\sharp$.
\begin{lemm}\label{lem:relboundsrho}
Let $\beta\in\R$, $B\in\BS^\beta$ and $A\in\BS\BE\BS^\alpha$ for some $\alpha>0$. Then
\begin{enumerate}
\item\label{approx_domain}
$\bigcup_{\zeta\,>\,\max\{\beta, 0\}}\widetilde{\sfH}^\zeta\subseteq\dom(B^\sharp)$,
\item\label{tildeHalpha}
$\dom(A^\sharp)=\widetilde{\sfH}^\alpha$.
\item\label{relbound}
Suppose $\beta < \alpha$ and $0 <\gamma < \alpha$ with $A^{\CO\CD}\in \BS^\gamma$. Then for $\beta\le 0$ the operator $B^\sharp$ is bounded, and, otherwise, for every $\phi\in\widetilde{\sfH}^\alpha$ and
\[
(\kappa, \er)\in\FE\left(A^\CD\right)\cap \Big\{\er
\geq\left(\left\|A^{\CO\CD}\right\|_0^{(\gamma)}/\kappa\right)^{1/(\alpha
-\gamma)}\Big\},
\]
we have
\begin{equation}\label{eq:relboundssharp}
\left\|B^\sharp\phi\right\|_{\FH} \le \frac{\er^{\beta -\alpha}\|B\|_0^{(\beta)}}{\kappa -\er^{\gamma -\alpha}\left\|A^{\CO\CD}\right\|_0^{(\gamma)}}\Big(\left\|A^\sharp\phi\right\|_{\FH} +\kappa \er^{\alpha}\left(1 +\kappa^{-1}\left\|A^\CD\right\|_0^{(\alpha)}\right)\|\phi\|_{\FH}\Big).
\end{equation}
In particular, $B^\sharp$ is infinitesimally $A^\sharp$-bounded in $\FH$.
\end{enumerate}
\end{lemm}

\begin{proof}
For $\beta\le 0$ the statements \emph{\eqref{approx_domain}} and~\eqref{relbound} follow from~\eqref{eq:normpreserving}. Let now $\beta>0$ and assume that $0 <\gamma < \alpha$ with $A^{\CO\CD}\in \BS^\gamma$. Following the proof of Lemma~\ref{lem:relative boundedness} and applying properties~\eqref{item:rhohomo} and~\eqref{item:rhoS0bounded} of the representation $\rho$ where necessary, we derive~\eqref{eq:relboundssharp} for $\phi\in\widetilde{\sfH}^\infty$. Consequently, the graph norm of $\rho(B)$ is dominated by the graph norm of $\rho(A)$, thus $\dom(B^\sharp)\supseteq \dom(A^\sharp)$. Applying~\eqref{eq:relboundssharp} for $A^{\CO\CD}$ instead of $B$ and $A^{\CD}$ instead of $A$, we conclude that the graph norms of $\rho(A)$ and $\rho(A^\CD)$ are equivalent, thus~\eqref{item:domsharpDES} implies $\dom(A^\sharp)=\dom((A^\CD)^\sharp)=\widetilde{\sfH}^\alpha$, which is~\eqref{tildeHalpha}. Now~\eqref{approx_domain} follows by varying $A\in \BS\BE\BS^\alpha$ with $\alpha >\beta$. Finally, we can extend~\eqref{eq:relboundssharp} from $\widetilde{\sfH}^\infty$ to $\widetilde{\sfH}^\alpha$ by density with respect to the graph norm of $A^\sharp$.
\end{proof}

Properties~\eqref{item:domsharpDES} and~\eqref{item:sharpAselfadjoint} of the map $\rho$ can also be extended to operators from the classes $\BE\BS^\alpha$, $\alpha>0$.
\begin{lemm}\label{lem:sharpWES}
Let $\alpha>0$ and $A\in\BE\BS^\alpha$. Then $\dom(A^\sharp)=\widetilde{\sfH}^\alpha$ and for all unitary $U\in\BS^0$
\begin{equation}\label{eq:Usharp*AsharpUsharp}
U^\sharp A^\sharp \big(U^\sharp\big)^* = \big(U A U^\dagger\big)^\sharp \quad \text{holds on} \ \widetilde{\sfH}^\alpha.
\end{equation}
Moreover, if $A$ is self-adjoint on $\sfH^\alpha$, then $A^\sharp$ is self-adjoint.
\end{lemm}

\begin{proof}
Assume first that $A\in\BS\BE\BS^\alpha$, so that $A^\CD\in\BD\BE\BS^\alpha$ and $A^{\CO\CD}\in\BS^\gamma$ for some $0<\gamma<\alpha$. According to Lemma~\ref{lem:relboundsrho}$(\ref{approx_domain}, \ref{tildeHalpha})$ we have that $\dom(A^\sharp)=\widetilde{\sfH}^\alpha\subseteq\dom((A^{\CO\CD})^\sharp)$. Moreover, if $A$ is self-adjoint, then $(A^\CD)^\sharp$ is self-adjoint on $\widetilde{\sfH}^\alpha$ and $\rho(A^{\CO\CD})$ is symmetric on $\widetilde{\sfH}^\infty$, as follows from properties~\eqref{item:sharpAselfadjoint} and~\eqref{item:rhohomo} of $\rho$, respectively. Since $\dom((A^{\CO\CD})^\sharp)\supseteq\widetilde{\sfH}^\alpha$ is the closure of $\widetilde{\sfH}^\infty$ with respect to the graph norm of $(A^{\CO\CD})^\sharp$, the operator $(A^{\CO\CD})^\sharp$ is also symmetric on $\widetilde{\sfH}^\alpha$. Moreover, by Lemma~\ref{lem:relboundsrho}\eqref{relbound} it is infinitesimally $A^\sharp$-bounded. Thus, \cite[Theorem~4.1.9]{BirmanSolomyak} implies that $A^\sharp=(A^\CD +A^{\CO\CD})^\sharp$ is self-adjoint on $\widetilde{\sfH}^\alpha$.

Let now $A\in\BE\BS^\alpha$. By definition, there exist $H\in\BS\BE\BS^\alpha$ and $V\in\BS^0$ unitary such that $A=V^\dagger HV$ on $\sfH^\infty$. Since $\rho$ is a $*$-homomorphism,
\begin{equation}\label{eq:rhoAprod}
\rho(A)=\rho(V)^\dagger \rho(H)\rho(V)
\end{equation}
holds on $\widetilde{\sfH}^\infty$. By Lemma~\ref{lem:sharpS0} the operator $V^\sharp$ is unitary, and property~\eqref{item:domsharpinv} implies that
\begin{equation}\label{eq:VsharptildeHalpha}
V^\sharp \widetilde{\sfH}^\alpha= \big(V^\sharp\big)^*\widetilde{\sfH}^\alpha=\widetilde{\sfH}^\alpha.
\end{equation}
We have already proved in Lemma~\ref{lem:relboundsrho}\eqref{tildeHalpha} that $\dom(H^\sharp)=\widetilde{\sfH}^\alpha$, thus the argument at the end of the proof of Proposition~\ref{prop:closedness and self-adjointness} implies that $(V^\sharp)^\ast H^\sharp V^\sharp$ is closed on $\widetilde{\sfH}^\alpha$. As by~\eqref{eq:rhoAprod} and Lemma~\ref{lem:sharpS0} this operator is an extension of $\rho(A)$, it follows that $\dom(A^\sharp)\subseteq\widetilde{\sfH}^\alpha$. Similarly, we have on $\widetilde{\sfH}^\infty$ that
\begin{equation}
\rho(H)=\rho(V)\rho(A)\rho(V)^\dagger
\end{equation}
and $V^\sharp A^\sharp (V^\sharp)^*$ is a closed operator on $V^\sharp\dom(A^\sharp)\subseteq \widetilde{\sfH}^\infty$. Thus,
\[
\widetilde{\sfH}^\alpha=\dom\left(H^\sharp\right)\subseteq
V^\sharp\dom\left(A^\sharp\right),
\]
and~\eqref{eq:VsharptildeHalpha} yields $\widetilde{\sfH}^\alpha\subseteq\dom(A^\sharp)$. Hence $\dom(A^\sharp)=\widetilde{\sfH}^\alpha$. Finally, let $U\in\BS^0$ be unitary. Then
\begin{equation}
\rho(U)\rho(A)\rho(U)^\dagger=\rho\left(U A U^\dagger\right)\subseteq \left(U A U^\dagger\right)^\sharp,
\end{equation}
so that on the domain $\widetilde{\sfH}^\alpha=\dom((U A U^\dagger)^\sharp)$ we have that $U^\sharp A^\sharp(U^\sharp)^*$ is a closed extension of $\rho(U) \rho(A)\rho(U)^\dagger$, in other words~\eqref{eq:Usharp*AsharpUsharp} holds. If $A$ is self-adjoint on $\sfH^\alpha$, then so is $H$, thus $H^\sharp$ by the first part of the proof. Hence, the self-adjointness of $A^\sharp$ follows from~\eqref{eq:Usharp*AsharpUsharp} with $U=V$ and $H$ instead of $A$.
\end{proof}

\subsection{The density of states measure}\label{sec:dsm}

Since $\FA$ is a factor of type $\rmI_\infty$ or $\rmII_\infty$, there exists, by definition, a semi-finite faithful normal trace $\FT$ on $\FA$, see~\cite[I.6 and I.8.4]{Dixmier}. Moreover, due to~\cite[I.6.4, Corollary]{Dixmier}, this trace is unique up to multiplication by a positive number. As in Definition~\ref{def:vNa}, we write $\FL\eta\FA$ to denote that $\FL\subset \FH$ is a closed linear subspace affiliated to $\FA$, i.e.~$P_{\FL}\in\FA$, where $P_\FL$ is the projection onto $\FL$. If $\FL\eta \FA$, the relative dimension of $\FL$ is defined by
\begin{align*}
\FD(\FL):=\FT(P_{\FL})\in [0,\infty].
\end{align*}
If $\FA$ is a $\rmi_\infty$-factor, the range of the relative dimension is $c\N_0\cup\set{\infty}$, for some $c>0$. It is $[0,\infty]$ if $\FA$ is a $_\infty$-factor.

\begin{defi}
Let $A\in\BS^0\cup\BE\BS^\infty$ be symmetric and $J\subseteq \R$ be a Borel measurable set. Denote by $E_J(A^\sharp)$ the spectral projection of $A^\sharp$ for $J$. We define the \emph{density of states measure} (DSM) of $A$ on $J$, relative to the representation $\rho$, by
\begin{equation}
N(J; A) := \FT\big(E_J(A^\sharp)\big)=\FD\big(E_J(A^\sharp)\FH\big).\label{eq:dfnids}
\end{equation}
\end{defi}


\begin{rema}\label{rem:iiinfinityfactor}
Usually, the dependence on the representation $\rho$ and the factor $\FA$ is unambiguous and is thus not reflected in the notation.
\end{rema}
The following corollary generalises~\cite[Lemma~4.4]{ParSht2012}. It follows directly from Lemma~\ref{lem:sharpWES} (or Lemma~\ref{lem:sharpS0} for $A\in\BS^0$) and the invariance of $\FT$ under unitary transformations in $\FA$. We remark at this point that, since $\widetilde{\BS}^\infty\eta\FA$, one has $U^\sharp\in\FA$ for every unitary $U\in\BS^0$, see Lemma~\ref{lem:sharpS0} and~\cite[\S 35.1]{Naimark}.
\begin{coro} \label{cor:unitary}
Let $U\in\BS^0$ be unitary and let $A\in\BS^0\cup\BE\BS^\infty$ be symmetric. Then one has $N(J;A) = N(J;U A U^\dagger)$ for any Borel measurable set $J\subseteq \R$.
\end{coro}

In the remainder of this section, we investigate the behaviour of the DSM for elliptic operators of positive order under perturbations. In~\cite{MorParSht2014,ParSht2012,ParSht2016} such an analysis was conducted for operators that are bounded from below and the particular case $J = (-\infty,
\lambda)$, $\lambda\in\R$.

Before continuing, let us introduce the following notation. For any interval $J =[s,t]\subset \R$, $s <t$ and $\eps \in\R$, we define
\[
J_{\eps}:=
\begin{cases}
\hfill\varnothing\hfill & \text{for } \eps < \dfrac{s-t}{2}, \\
[s-\eps, t+\eps] & \text{otherwise}.
\end{cases}
\]
The following lemma gives us a variational characterisation of the DSM (cf.~\cite[Lemma~4.1]{ParSht2012}).

\begin{lemm}\label{lem:variational}
Let $A\in \BS^0\cup\BE\BS^\infty$ be symmetric. Then, for any interval $J=[q-r,q+r]$ with $q\in\R$ and $r>0$, we have
\begin{multline}\label{eq:variational}
N(J;A) = \sup\big\{\FD(\FL) : \FL\subset\dom(A^\sharp),\ \FL \eta \FA,
\\\text{\emph{and} }
\left\|\left(A^\sharp-q\right)\phi\right\|_{\FH}\leq r\|\phi\|_\FH \ \forall \phi
\in \FL\big\}.
\end{multline}
The analogous statement holds for the open interval $J=(q-r,q+r)$ with strict inequality in~\eqref{eq:variational}.
\end{lemm}

\begin{rema}
Usually, variational characterisations such as~\eqref{eq:variational} are given in terms of quadratic forms rather than norms. The reason why we cannot do so is because we do not assume the operator $A^\sharp$ to be semi-bounded, $J$ a semi-infinite interval. One can interpret Lemma~\ref{lem:variational} in terms of quadratic forms as usual for the nonnegative operator $(A^\sharp - q)^2$.
\end{rema}

\begin{proof}
Choosing $\FL := E_J(A^\sharp) \FH$, we observe that $N(J;A)$ is at most the right hand side of~\eqref{eq:variational}. Suppose that there exists a subspace $\FL$ that satisfies the assumptions on the righthand side of~\eqref{eq:variational} and $\FD(\FL) > \FD(E_J(A^\sharp)\FH)$. Then~\cite[\S 37.1, Lemma]{Naimark} implies that $\FL$ contains an element $\phi$ orthogonal to $E_J(A^\sharp)\FH$, implying that $\|(A^\sharp-q)\phi\|_{\FH}^2>r^2\|\phi\|_{\FH}^2$, which is a contradiction.
\end{proof}

The following lemma generalises~\cite[Corollary~4.3]{ParSht2012} to operators that are not necessarily bounded below and unbounded perturbations.

\begin{lemm} \label{lem:monotonicity}
Let $A\in\BS\BE\BS^\alpha$, $\alpha>0$, and $B\in\BS^\beta$, $\beta<\alpha$, symmetric operators. Let $J:=[q-r,q+r]\subset \R$ be the interval of length $2r > 0$ centred at $q\in\R$. Then there exists a constant $C\geq 0$ depending only on $A$ and $\beta$ such that, for
\begin{equation}\label{eq:epsa}
\epsilon:=\epsilon_{J,A,B}:=
\begin{cases}
\hfill\norm{B}\hfill & \text{if }\beta \leq 0,\\
\dfrac{\snorm{B}{\beta}{0}}{2 +\snorm{B}{\beta}{0}}\Big(r + \abs q+ C\big(1+\snorm{B}{\beta}{0}\big)^{\frac{\alpha}{\alpha-\beta}}\Big)& \text{if } \beta>0,
\end{cases}
\end{equation}
the inequality
\begin{equation}\label{eq:bounds}
N(J_{-\epsilon}; A)\leq N(J; A+B)\leq N(J_{\epsilon}; A)
\end{equation}
holds.
\end{lemm}

\pagebreak
\begin{proof}
In view of Lemma~\ref{lem:relboundsrho}(\ref{approx_domain},\ref{relbound}) and property~\eqref{item:rhohomo} of the representation $\rho$, one has that $(A+B)^\sharp =A^\sharp+B^\sharp$ on $\dom(A^\sharp)\subseteq\dom(B^\sharp)$. Fix $\phi\in \FL := E_{J}((A+B)^\sharp)\FH\subseteq \dom((A+B)^\sharp)=\dom(A^\sharp)$, so that
\begin{align}
\norm{\left(A^\sharp + B^\sharp - q\right)\phi}_\FH &\le r \norm \phi_\FH.\label{A+B}
\intertext{We will show that}
\norm{\left(A^\sharp - q\right) \phi}_\FH &\le (r + \eps) \norm \phi_\FH\label{eq:A-Meps}
\end{align}
holds, which in view of~\eqref{eq:variational} implies the second inequality in~\eqref{eq:bounds}. Since
\begin{equation}\label{eq:A-Mxtriag}
\begin{split}
\norm{\left(A^\sharp-q\right)\phi}_\FH &\le \norm{\left(A^\sharp+B^\sharp - q\right) \phi}_\FH +
\norm{B^\sharp \phi}_\FH\\
&\le r \norm \phi_\FH + \norm{B^\sharp \phi}_\FH,
\end{split}
\end{equation}
it is sufficient to estimate $\|B^\sharp \phi\|_\FH$. For $\beta \le 0$, Lemma~\ref{lem:sharpS0} and Corollary~\ref{cor:domain} imply
\begin{equation}
\norm{B^\sharp \phi}_\FH \le
\norm{B^\sharp}\norm{\phi}_\FH=\norm{B}\norm{\phi}_\FH = \eps \norm \phi_\FH,
\end{equation}
and~\eqref{eq:A-Meps} follows from~\eqref{A+B} and~\eqref{eq:A-Mxtriag}.

From now on, we consider $\beta > 0$. By assumption we can choose $\gamma\in (\beta,\alpha)$ such that $A^{\CO\CD} \in
\BS^\gamma$. Let $(\kappa,\er)\in\FE(A^\CD)$ with
\begin{equation}
\er\geq\max\left\{\left(\frac{4\left(1+\snorm{B}{\beta}{0}\right)}{\kappa}\right)^{1/(\alpha-\beta)};\quad\left(\frac{2\snorm{A^{\CO\CD}}{\gamma}{0}}{\kappa}\right)^{1/(\alpha-\gamma)}\right\}.
\end{equation}
As $\phi\in \dom(A^\sharp)$, Lemma~\ref{lem:relboundsrho}\eqref{relbound} yields
\begin{equation}\label{eq:relativbdspecificR}
\begin{aligned}
\norm{B^\sharp \phi}_\FH&\leq 2\snorm{B}{\beta}{0}\left[\er^{\beta-\alpha}\kappa^{-1}
\norm{A^\sharp \phi}_\FH +
\er^\beta\left(1+\kappa^{-1}\snorm{A^\CD}{\alpha}{0}\right)\norm{\phi}_\FH\right]\\
&\leq \frac{\snorm{B}{\beta}{0}}{2}\left[\frac{\norm{\left(A^\sharp-q\right)\phi}_\FH +|q|\|\phi\|_\FH}{1+\snorm{B}{\beta}{0}}+ C\left(1+\snorm{B}{\beta}{0}\right)^{\frac{\beta}{\alpha-\beta}}\norm{\phi}_\FH\right],
\end{aligned}
\end{equation}
where $C$ is a constant only depending on $A$ and $\beta$. Combining~\eqref{eq:A-Mxtriag} and~\eqref{eq:relativbdspecificR}, we get
\begin{multline*}
\frac{2 +\snorm{B}{\beta}{0}}{2\big(1+\snorm{B}{\beta}{0}\big)}\left\|\left(A^\sharp-q\right)\phi\right\|_\FH
 \\
\leq
\left[r+\frac{\snorm{B}{\beta}{0}}{2\left(1+\snorm{B}{\beta}{0}\right)}\left(|q|+C\left(1+\snorm{B}{\beta}{0}\right)^{\frac{\alpha}{\alpha-\beta}}\right)\right]\norm{\phi}_\FH.
\end{multline*}
Hence, we arrive at~\eqref{eq:A-Meps} with $\epsilon$ as in~\eqref{eq:epsa}.

For the first inequality in~\eqref{eq:bounds} the only non-trivial case is $\epsilon\leq r$. For all $\phi\in E_{J_{-\eps}}(A^\sharp)\FH\subset\dom(A^\sharp)$ we have
\begin{align}
\norm{\left(A^\sharp-q\right)\phi}_\FH &\leq (r-\epsilon)\norm{\phi}_\FH.\label{eq:A-Mr-eps}
\intertext{This implies}
\norm{\left(A^\sharp+B^\sharp-q\right)\phi}_\FH &\leq (r-\epsilon)\norm{\phi}_\FH+\norm{B^\sharp \phi}_\FH,
\end{align}
where in view of~\eqref{eq:relativbdspecificR} and~\eqref{eq:epsa}
\begin{equation}
\norm{B^\sharp \phi}_\FH \leq \frac{\snorm{B}{\beta}{0}}{2 +2\snorm{B}{\beta}{0}}\left(r + \abs q+ C\left(1+\snorm{B}{\beta}{0}\right)^{\frac{\alpha}{\alpha-\beta}}\right)\norm{\phi}_\FH
\leq\epsilon \norm{\phi}_\FH.
\end{equation}
Thus, the first inequality in~\eqref{eq:bounds} follows and the Lemma~\ref{lem:monotonicity} is proved.
\end{proof}


The next lemma deals with perturbations that are ``spectrally far'' from a given interval. It is a generalisation of~\cite[Lemma~11.1]{MorParSht2014} for operators that are not necessarily bounded below.


\begin{lemm} \label{lem:spectralperturb}
For $\alpha>0$, $\beta < \alpha$ let $H_0 \in \BD\BE\BS^\alpha$, $B \in \BS^\beta$, and $A\in\BS^0$ be symmetric operators and set $H:=H_0+B\in\BS\BE\BS^\alpha$. Suppose that there exists a family of orthogonal projections $\set{P_l}_{l=0}^L$ with $P_l\in \BS^{-\alpha}$, $0\leq l\leq L-1$, and $P_L\in \BS^0$ that all commute with $H_0$ and satisfy
\begin{equation}\label{eq:projections}
\sum_{l=0}^L P_l = I, \quad \text{and} \quad A = A P_0, \quad B_{n,l} := P_n B P_l = 0, \quad \text{for} \ \abs{n-l}>1.
\end{equation}
Moreover, let $J = (q-r,q+r)$ be an interval such that
\begin{equation}\label{eq:distspectrum}
D_l:=\operatorname{dist}\big(J,\sigma\big((P_l H P_l)^\sharp\big)\big)>0,\quad \text{for all}\ \ 0\leq l<L.
\end{equation}
Finally, assume that
\begin{equation}\label{eq:tecassL}
3^Lr\geq d_L:=\min\limits_{1\,\leq\,l\,<\,L} D_l
\end{equation}
and
\begin{equation}\label{eq:projectionnorm}
\max_{0\,\le\,l\,<\,L} (\norm{B_{l,l-1}} + \norm{B_{l,l+1}})/D_l \le 1/4,
\end{equation}
where we use the convention $B_{0,-1}:=0$.

Then for
\begin{equation}\label{eq:epschoice}
\eps := 3^{2 - \frac{L}{2}} \left(\frac{r}{d_L}\right)^{1/2} \|A\|
\end{equation}
we have that
\begin{equation}
N(J_{-\eps};H) \le N(J;H + A) \le N\left(J_{\eps};H\right).\label{eq:approximationcontrol}
\end{equation}
\end{lemm}
\begin{proof}
We only prove the first inequality; the second inequality follows analogously. It suffices to show that for any $\phi \in E_{J_{-\eps}}(H^\sharp) \FH\subseteq\dom(H^\sharp)=\widetilde{\sfH}^\alpha$, one has
\[
\norm{\left(H^\sharp + A^\sharp - q\right) \phi}_\FH \leq r \norm \phi_\FH.
\]
For any $K \in \N$, we split the interval $J_{-\epsilon}$ into
$2K+1$ subintervals of equal width: for $- K \le k \le K - 1$, set
\begin{equation}\label{eq:intIk}
I_k := \left(q + (2k - 1)\frac{(r - \eps)}{2K+1}, q + (2k+1)\frac{(r -
\eps)}{2K+1}\right]
\end{equation}
and
\begin{equation}\label{eq:intIK}
I_K:= \left(q + (2K -1)\frac{(r-\eps)}{2K+1}, q + r - \eps\right).
\end{equation}
For $\phi \in E_{J_{-\eps}}(H^\sharp)\FH$ and $- K \le k \le K$ define $\phi^{k}:= E_{I_k} (H^\sharp)\phi\in\widetilde{\sfH}^\alpha$ and
\begin{equation}\label{eq:etak}
\eta^{k} := H^\sharp \phi^{k} - \left(q + 2k \frac{(r - \eps)}{2K+1}\right) \phi^k,
\end{equation}
so that
\begin{equation}\label{eq:normetak}
\norm{\eta^{ k}}_\FH \le \frac{r}{(2K + 1)} \norm{\phi^{ k}}_\FH
\end{equation}
holds. We also introduce
\[
\phi_l^k:=P_l^\sharp\phi_k \quad \text{ and } \quad
\eta_l^k:=P_l^\sharp\eta_k, \quad \text{
for } -K\leq k\leq K \, \text{ and } \, 0\leq l \leq L.
\]
For $0\leq l<L$, we clearly have
$P_l^\sharp H^\sharp=(P_lH)^\sharp$ on $\widetilde{\sfH}^\infty$ and, since $P_lH\in\BS^0$, this identity extends to $\widetilde{\sfH}^\alpha$. Moreover, $P_l$ commutes with $H_0$ so that~\eqref{eq:projections} implies that on $\widetilde \RH^\alpha$
\begin{align*}
P_l^\sharp H^\sharp &= (P_lH)^\sharp=(P_lHP_l)^\sharp+ B_{l,l-1}^\sharp+B_{l,l+1}^\sharp\\
&=(P_lHP_l)^\sharp P_l^\sharp+ B_{l,l-1}^\sharp P_{l-1}^\sharp+B_{l,l+1}^\sharp P_{l+1}^\sharp,
\end{align*}
where we use the convention $P_{-1}:=0$. Thus, applying $P_l^\sharp$ to~\eqref{eq:etak}, we arrive at
\begin{equation}\label{eq:etalk}
\eta_l^{ k} = B_{l,l-1}^\sharp \phi_{l-1}^{ k} + \left((P_lHP_l)^\sharp -
\left(q + 2 k\frac{(r-\eps)}{2K+1}
\right) \right)\phi_l^{ k} + B_{l,l+1}^\sharp \phi_{l+1}^{k},
\end{equation}
for $0 \le l <L$, and Lemma~\ref{lem:sharpS0} together with~\eqref{eq:distspectrum} and~\eqref{eq:normetak} gives for $0 \le l < L$,
\begin{equation}\label{eq:normphilk}
\begin{aligned}
\norm{\phi_l^{ k}}_\FH &\le D_l^{-1}\Big(\norm{\eta_l^{ k}}_\FH + \norm{B_{l,l-1}}
\norm{\phi_{l-1}^{ k}}_\FH + \norm{B_{l,l+1}} \norm{\phi_{l+1}^{ k}}_\FH \Big)\\
&\le \frac{r}{(2K+1) d_L} \norm{\phi^{k}}_\FH +
\frac{\norm{\phi_{l-1}^{k}}_\FH +
\norm{\phi_{l+1}^{k}}_\FH}{4}.
\end{aligned}
\end{equation}
Recursively for $0 \le l < L$ we deduce that
\begin{equation}
\norm{\phi^{k}_l}_\FH \le \frac{2r}{(2K+1) d_L}\norm{\phi^{k}}_\FH +
\frac{1}{3}\norm{\phi_{l+1}^k}_\FH.
\end{equation}
Hence, employing the trivial bound $\|\phi_L^k\|_\FH\le \|\phi^k\|_\FH$, we get that
\begin{equation}\label{eq:phi0phik}
\norm{\phi_0^k}_\FH \le \left(\frac{3r}{(2K+1) d_L} + 3^{-L}\right)\norm{\phi^k}_\FH.
\end{equation}
In view of Lemma~\ref{lem:sharpS0}, it follows that for all $- K \le k \le K$,
\begin{equation}
\begin{aligned}
\norm{A^\sharp \phi^{k}}_\FH &= \norm{(AP_0)^\sharp \phi^{k}}_\FH\\&=
\norm{A^\sharp \phi_0^k}_\FH \\&\le \left(\frac{3r}{(2K+1) d_L} + 3^{-L}\right) \norm{A}\norm{\phi^k}_\FH,
\end{aligned}
\end{equation}
whence the Cauchy--Schwarz inequality and the Pythagorean theorem yield
\begin{equation}
\begin{aligned}
\norm{A^\sharp\phi}_\FH &\le \sum_{-K\,\le\,k\,\le\,K} \norm{A^\sharp
\phi^k}_\FH \\
&\le
\left(\frac{3r}{(2K+1) d_L} + 3^{-L}\right)\sqrt{2K+1}\norm{A}
\norm{\phi}_\FH.
\end{aligned}
\end{equation}
We choose
\begin{equation}\label{eq:deltachoice}
K = \flo{\frac{3^{L+1}r}{2d_L} - \frac 1 2} + 1,
\end{equation}
so that $\frac{3^{L+1}r}{d_L}\leq 2K+1\leq \frac{3^{L+2}r}{d_L}$. Then, by~\eqref{eq:tecassL}, we have
\begin{equation}\label{eq:boundsecterm}
\left(\frac{3r}{(2K + 1)d_L} + 3^{-L}\right)\sqrt{2K+1} \le 3^{2 -
\frac{L}{2}} \Big(\frac{r}{d_L}\Big)^{1/2}.
\end{equation}
Consequently, we arrive at
\begin{equation}\label{eq:therightthing}
\begin{split}
\norm{\left(H^\sharp + A^\sharp - q\right) \phi}_\FH &\le \norm{\left(H^\sharp - q\right)
\phi}_\FH+ \norm{A^\sharp\phi}_\FH \\
&\le \left((r - \eps) + 3^{2 - \frac{L}{2}} \left(\frac{r}{d_L}\right)^{1/2}\norm{A}\right)
\norm{\phi}_\FH\\
&= r \norm \phi_\FH,
\end{split}
\end{equation}
where we used that $\phi\in E_{J_{-\eps}}(H^\sharp)$ and the value of $\eps$ given in~\eqref{eq:epschoice}.
\end{proof}

\section{Gauge Transform}\label{sec:gt}

Let $\alpha\in \R$ and $A = \Op(a) \in \BS\BE\BS^\alpha$ be symmetric, thus extends to a self-adjoint linear operator on $\sfH^{\alpha}$ by Proposition~\ref{prop:closedness and self-adjointness}.

\begin{defi}
For every symmetric $\Psi\in\BS^0$, the unitary transformation of $A$ into
\[
[A] := [A]_\Psi := \exp(- \rmi \Psi) A \exp(\rmi \Psi)
\]
is called a
\emph{gauge transform}.
\end{defi}

\noindent We remark here that, due to Lemma~\ref{lem:product} and Corollary~\ref{cor:domain}, the series
\begin{equation}\label{eq:exppsi}
\exp(\rmi\Psi)=\sum\limits_{k=0}^\infty \frac{\big(\rmi\Psi\big)^k}{k!}
\end{equation}
converges both in $\BS^0$ and in the operator norm. In particular, Lemma~\ref{lem:complete} implies that $\exp(\rmi\Psi)\in\BS^0$, whence $\exp(\rmi\Psi)$ is unitary and $[A]_\Psi\in\BE\BS^\alpha$ is symmetric. The following lemma provides an expansion of $[A]_\Psi$ into a series of multiple commutators of $A$ with $\Psi$, see~\eqref{eq:defadk} for the definition of $\ad^k$.
\begin{lemm}\label{lem:A'series}
We have
\begin{equation}\label{eq:H'formalseries}
[A]_\Psi = \sum_{k = 0}^\infty \frac{1}{k!}\ad^k(A;\Psi),
\end{equation}
where the series converges absolutely in $\BS^\alpha$.
\end{lemm}

\begin{proof}
Lemma~\ref{lem:product} yields the bounds
\begin{equation}
\snorm{\Psi^jA\Psi^m}{\alpha}{l}\leq \big(\snorm{\Psi}{0}{l}\big)^j\snorm{A}{\alpha}{l}\big(\snorm{\Psi}{0}{l+|\alpha|}\big)^m, \quad \text{for all}\ l\geq 0.
\end{equation}
Thus, the double series
\begin{equation}
[A]_\Psi=\sum_{j= 0}^\infty \frac{(-\rmi \Psi)^j}{j!} A \sum_{m =0}^\infty \frac{(\rmi \Psi)^m}{m!} = \sum_{j,m= 0}^\infty \frac{(-\rmi \Psi)^j}{j!} A \frac{(\rmi \Psi)^m}{m!}
\end{equation}
converges absolutely in $\BS^\alpha$. Recursively, we obtain
\begin{equation}
\ad^k(A;\Psi) = k! \sum_{j + m = k} \frac{(-\rmi \Psi)^j}{j!} A \frac{(\rmi \Psi)^m}{m!}, \quad\text{ for all $k \ge 0$}.\qedhere
\end{equation}
\end{proof}

In the remainder of this section, we look at gauge transforms that result in an operator $[A]_\Psi$ that is closer to a diagonal operator (i.e.~an operator in $\BD\BE\BS^\alpha$) than $A$. More precisely, we construct $\Psi$ in such a way that the gauge transform removes as much of the off-diagonal part $A^{\CO\CD}$ from $A$ as possible. Let $\beta<\alpha$ such that $A^\COD\in\BS^\beta$. Then we aim at
\begin{equation}\label{eq:H'goal}
[A]_\Psi=A^\CD + A^\CR + R,
\end{equation}
where $A^\CR\in\BS^\beta$ is an off-diagonal
\emph{resonant} part (which our transformation cannot eliminate) and $R\in\BS^\gamma$ for some $\gamma <\beta$. The exact form of the operators $A^\CR$ and $R$ depends on the choice of $\Psi$.

As a first step towards~\eqref{eq:H'goal}, let us rewrite the series~\eqref{eq:H'formalseries} as
\begin{equation}\label{eq:firstorder}
[A]_\Psi = A^\CD + A^{\CO\CD} + \ad\left(A^\CD;\Psi\right) + R,
\end{equation}
with
\begin{equation}\label{eq:defRconsecgauge}
R:=\ad(A^{\CO\CD};\Psi)+\sum\limits_{k= 2}^\infty\frac{1}{k!}\ad^k(A;\Psi).
\end{equation}
Suppose that $\Psi \in \BS^\zeta$ with $\zeta \in \R$ and let $\kappa \in \R$. In order to achieve that $R\in\BS^\gamma$ for some $\gamma<\beta$, we can use the following estimates:
\begin{enumerate}
\item If $\zeta <0$, then by Lemma~\ref{lem:weakcommest} we get $\ad(X;\Psi)\in\BS^{\varkappa +\zeta}$ for all $X \in \BS^\varkappa$. We call a gauge transform that only uses these trivial bounds on the commutator norms \emph{weak}.
\item Sometimes the structure of the commutators allows us to prove $\ad(X;\Psi)\in\BS^{\varkappa +\zeta -\varepsilon}$ for some $\varepsilon >0$ and appropriate $X \in \BS^\varkappa$. A gauge transform exploiting this improvement shall be called \emph{strong}.
\end{enumerate}

As we will see, the main issue with the strong gauge transform is that some conditions under which it can be used may not be formally invariant under the use of the gauge transform, which is in general an iterative scheme. Furthermore, due to combinatorial issues it may be harder to verify that those conditions are still satisfied as the number of steps increase. However, as we will see, in many situations it is sufficient to make one step of the strong gauge transform, and proceed from there with the weak one.



\subsection{The commutator equation}
We recall that after the gauge transform we would like to arrive at the operator $[A]_\Psi$ as in~\eqref{eq:H'goal}, in the best possible case with $A^{\CR}=0$. Comparing~\eqref{eq:H'goal} with~\eqref{eq:firstorder} we obtain that $A^{\CR}=0$ is equivalent to the commutator equation
\begin{equation}\label{eq:comm}
\ad\left(A^\CD;\Psi\right) + A^{\COD} = 0
\end{equation}
for $\Psi = \Op(\psi)$.

Let $A=\Op(a)$ and $\Theta$ be a frequency set for $a$. By~\eqref{eq:ad symbol}, equation~\eqref{eq:comm} is solved if $\Theta$ is a frequency set for $\psi$ and
\begin{equation}
a^\CD(\theta \act \xi) \psi_\theta(\xi) - \psi_\theta(\xi)a^\CD(\xi) = \rmi a^\COD_\theta(\xi)\quad
\end{equation}
holds for all $\theta\in\Theta':=\Theta\setminus \set{\id}$ and $\xi\in\Xi$. This leads to
\begin{equation}\label{eq:eq}
\psi_\theta(\xi) = \frac{\rmi a^\COD_\theta(\xi)}{a^{\CD}(\theta\act\xi) - a^\CD(\xi)},
\end{equation}
for $\theta\in\Theta'$ and $\xi\in\Xi$. However, the problem of small denominators $a^\CD(\theta\act\xi) - a^\CD(\xi)$ for some pairs $(\theta,\xi)$ generally prevents such choice of $\psi$. This motivates the following definition.

\begin{defi}\label{def:resonant}
For $\delta \in \R,\ \resc>0$, and $\theta \in G$, we call a set $\Lambda_\theta^{\delta,\resc}\subseteq \Xi$ a $\delta$\emph{-resonant region} generated by $\theta$ for $A^\CD$ if it satisfies
\begin{equation}\label{eq:resonantregion}
\Lambda_\theta^{\delta,\resc} \supseteq \set{\xi \in \Xi : \abs{a^\CD(\theta\act\xi) - a^\CD(\xi)} \leq \resc \ang \xi^\delta}.
\end{equation}
A corresponding \emph{resonance cut-off} is a function $\rco :=\rco^{\delta, s} : G\times\Xi \to \R$, mapping $(\theta,\xi)\mapsto \rco_\theta^{\delta, s}(\xi)$, such that for all $\theta\in G$, we have
\begin{equation}\label{eq:dennonres}
\begin{gathered}
0\leq \rco \leq 1,\\
\rco_{\theta}^{\delta,s}(\xi)=0,\text{ for all }\xi\in\Lambda_\theta^{\delta,s},\\
\overline{\rco_{\theta^{-1}}(\theta\act\xi)}=\rco_\theta(\xi),\text{ for all }\xi\in\Xi.
\end{gathered}
\end{equation}
For a fixed resonance cut-off, we define the \emph{resonant} part $B^{\CR}:=\Op(b^{\CR})$ and the \emph{non-resonant} part $B^{\CNR}:=\Op(b^{\CNR})$ of any operator $B=\Op(b)\in\BS^\infty$ via their symbols
\begin{equation}\label{eq:resnonres}
\begin{aligned}
b^\CR &{:=}\, b^{\CO\CD}\left(1 - \rco^{\delta,s}\right), \\
b^{\CNR} &{:=}\, b^{\CO\CD}\rco^{\delta,s}.
\end{aligned}
\end{equation}
\end{defi}
\begin{rema}\label{rem:res}\leavevmode
\begin{enumerate}\romanenumi
\item For any $\delta\in\R$ and $s>0$, the only $\delta$-resonant region generated by $\id$ is $\Lambda_{\id}^{\delta,s}=\Xi$. Hence, every resonance cut-off $\rco$ satisfies $\rco_{\id}\equiv 0$.
\item \label{item:remrescutoff}
If $\Lambda_\theta^{\delta,s}$ satisfies
\begin{equation}
\Lambda_{\theta^{-1}}^{\delta,s}=\theta\act\Lambda_{\theta}^{\delta,s},\quad \text{ for all } \ \theta\in G,
\end{equation}
then the resonance cut-off $\rco$ can be chosen as
\begin{equation}\label{eq:rescutoffstandard}
\rco_\theta(\xi):=\bone_{\Xi\setminus\Lambda_{\theta}^{\delta,s}}(\xi), \quad \text{ for all } (\theta,\xi)\in G\times\Xi.
\end{equation}
\item \label{item:remres}
If $B^{\CO\CD}\in\BS^\gamma$, $\gamma\in\R$, then
\begin{equation}
B^{\CO\CD}=B^{\CNR}+B^{\CR},
\end{equation}
and
\begin{equation}\label{eq:estpartbyfull}
\snorm{B^{\CNR}}{\gamma}{l}\leq\snorm{B^{\CO\CD}}{\gamma}{l},\qquad
\snorm{B^{\CR}}{\gamma}{l}\leq\snorm{B^{\CO\CD}}{\gamma}{l}
\end{equation}
hold for all $l\geq 0$. If $B$ is symmetric, then so are $B^{\CD}$, $B^{\CNR}$, and $B^{\CR}$.
\end{enumerate}
\end{rema}



With the help of Definition~\ref{def:resonant}, the problem of small denominators in~\eqref{eq:eq} can be circumvented. Let $\delta\in\R$, $s>0$, and fix a resonance cut-off $\rco$ corresponding to $\delta$-resonant regions $\Lambda_\theta^{\delta,\resc}, \ \theta\in G$, for $A^\CD$. Using~\eqref{eq:resnonres}, we define
\begin{equation}\label{eq:defpsiabstract}
\psi_\theta^{\delta,s}(\xi):=
\begin{cases}
\dfrac{ia^\CNR_\theta(\xi)}{a^\CD(\theta\act\xi)-a^\CD(\xi)}&
\text{if } \theta\in\Theta',\\[2ex] 0 & \text{otherwise}.
\end{cases}
\end{equation}
Recall that $A^{\COD}\in\BS^\beta$ so that, in view of Remark~\ref{rem:res}\eqref{item:remres}, $A^{\CNR}\in \BS^\gamma$ for some $\gamma\leq\beta$.
\begin{lemm}\label{lem:psi}
Let $\gamma\leq \beta$ with $A^{\CNR}\in \BS^\gamma$. Then~\eqref{eq:defpsiabstract} defines a symbol $\psi^{\delta,s} \in \BS^{\gamma-\delta}$. The operator $\Psi:=\Op(\psi^{\delta,s})$ is symmetric with
\begin{equation}\label{eq:estimatepsinorm}
\snorm{\Psi}{\gamma - \delta}{l} \le \frac 1 \resc \snorm{A^\CNR}{\gamma}{l}, \ \text{for all}\ l\geq 0.
\end{equation}
It satisfies
\begin{equation}\label{eq:modifiedcomm}
\ad\left(A^\CD;\Psi\right)+A^{\CNR}=0.
\end{equation}
\end{lemm}
\begin{proof}
The bounds~\eqref{eq:estimatepsinorm} follow directly from~\eqref{eq:defpsiabstract} and~\eqref{eq:resonantregion}--\eqref{eq:resnonres}. The equation~\eqref{eq:modifiedcomm} follows as in~\eqref{eq:comm}--\eqref{eq:eq} with $\CNR$ replacing $\CO\CD$.
\end{proof}


In view of~\eqref{eq:modifiedcomm}, \eqref{eq:firstorder} takes the form
\begin{equation}\label{eq:A'}
[A]_\Psi=A^\CD+A^{\CR}+R,
\end{equation}
with $R$ defined in~\eqref{eq:defRconsecgauge}.

\subsection{Weak gauge transform}\label{subsec:weakgauge}

Let $\gamma\leq\beta$ such that $A^{\CNR}\in\BS^\gamma$, $A^{\COD}\in\BS^\beta$. We choose $\delta >\gamma$, so that $\gamma -\delta< 0$ in Lemma~\ref{lem:psi}. Note that $\delta$ determines the size of the resonant regions and thus the efficiency of the gauge transform.
\begin{lemm}\label{lem:onestepweakgauge}
Let $\Psi=\Op(\psi^{\delta, s})$ be the operator defined in~\eqref{eq:defpsiabstract} and $R$ be as in~\eqref{eq:defRconsecgauge}. Then $\ad(A^\COD;\Psi)$, $R\in \BS^{\beta+\gamma - \delta}$ are symmetric and, for all $l\geq 0$,
\begin{equation}\label{eq:comaodpsi}
\snorm{\ad\left(A^\COD;\Psi\right)}{\beta+\gamma-\delta}{l}\le \frac{2}{\resc} \snorm{A^\COD}{\beta}{l + \abs{\gamma - \delta}}\snorm{A^\CNR}{\gamma}{l + \abs{\beta}},
\end{equation}
and
\begin{equation}\label{eq:Rnormboundsweak}
\snorm{R}{\beta+\gamma-\delta}l \le \frac{3}{\resc} \snorm{A^\COD}{\beta}{l + \abs{\gamma - \delta}}\snorm{A^\CNR}{\gamma}{l + \abs\beta+\abs{\gamma - \delta}}
\exp\left(\frac 2 \resc \snorm{A^\CNR}{\gamma}{l + \abs\beta+\abs{\gamma - \delta}}\right).
\end{equation}
\end{lemm}

\begin{proof}
The estimates~\eqref{eq:normestcomm} and~\eqref{eq:normestcommzeroorder} together with $\delta>\gamma$ and Lemma~\ref{lem:psi} imply that, for all $k\in\N$,
\begin{equation}\label{eq:estmultcomm}
\begin{split}
\snorm{\ad^k\left(A^\COD;\Psi\right)}{\beta+\gamma - \delta}{l} &\le 2
\snorm{\ad^{k-1}\left(A^\COD;\Psi\right)}{\beta}{l + \abs{\gamma - \delta}} \snorm{\Psi}{\gamma - \delta}{l + \abs \beta} \\
&\le \frac{2^{k}}{\resc} \left(\snorm{\Psi}{0}{l + \abs\beta+\abs {\gamma-\delta}}\right)^{k-1} \snorm{A^\COD}{\beta}{l + \abs{\gamma - \delta}}\snorm{A^\CNR}{\gamma}{l + \abs{\beta}}.
\end{split}
\end{equation}
Thus, \eqref{eq:comaodpsi} follows by choosing $k = 1$. Moreover, \eqref{eq:estimatepsinorm} implies that for all $k\in\N$,
\begin{equation}\label{eq:adkODweak}
\begin{split}
\snorm{\ad^k\left(A^\COD;\Psi\right)}{\beta+\gamma - \delta}{l} &\le \left(\frac{2}{\resc}\snorm{A^{\CNR}}{\gamma}{l + \abs\beta+\abs {\gamma-\delta}}\right)^{k} \snorm{A^\COD}{\beta}{l + \abs{\gamma - \delta}}.
\end{split}
\end{equation}
Similarly, we get from~\eqref{eq:modifiedcomm} that for all $k\geq 2$,
\begin{equation}\label{eq:adkDweak}
\begin{split}
\snorm{\ad^k\left(A^\CD;\Psi\right)}{\beta+\gamma - \delta}{l} &=\snorm{\ad^{k-1}\left(A^\CNR;\Psi\right)}{\beta+\gamma - \delta}{l}\\
&\le \left(\frac{2}{\resc}\snorm{A^{\CNR}}{\gamma}{l + \abs\beta+\abs {\gamma-\delta}}\right)^{k-1} \snorm{A^\CNR}{\beta}{l + \abs{\gamma - \delta}}.
\end{split}
\end{equation}
Hence, the bounds~\eqref{eq:Rnormboundsweak} follow from~\eqref{eq:defRconsecgauge}, \eqref{eq:adkODweak}, \eqref{eq:adkDweak}, and~\eqref{eq:estpartbyfull}.
\end{proof}

Lemmata~\ref{lem:psi} and~\ref{lem:onestepweakgauge} have the following immediate corollary, which follows by choosing $\gamma=\beta$ and applying~\eqref{eq:estpartbyfull}.
\begin{coro}\label{cor:onestepweakgauge}
Let $\Psi=\Op(\psi^{\delta, s})$ be the operator defined in~\eqref{eq:defpsiabstract} and $R$ be as in~\eqref{eq:defRconsecgauge}. Assume that $\delta>\beta$. Then $\Psi \in \BS^{\beta-\delta}$ is symmetric with
\begin{equation}\label{eq:estimatepsinormcor}
\snorm{\Psi}{\beta - \delta}{l} \le \frac 1 \resc \snorm{A^\COD}{\beta}{l}, \ \text{for all}\ l\geq 0.
\end{equation}
Moreover, $\ad(A^\COD;\Psi)$, $R\in \BS^{2\beta- \delta}$ are symmetric and for all $l\geq 0$,
\begin{align}
\snorm{\ad\left(A^\COD;\Psi\right)}{2\beta-\delta}{l}&\le \frac{2}{\resc} \left(\snorm{A^\COD}{\beta}{l + \abs\beta + \abs{\beta - \delta}}\right)^2,
\\
\intertext{and}
\snorm{R}{2\beta-\delta}l &\le \frac{3}{\resc} \left(\snorm{A^\COD}{\beta}{l + \abs\beta + \abs{\beta - \delta}}\right)^2 \exp\left(\frac 2 \resc \snorm{A^\COD}{\beta}{l + \abs\beta+\abs{\beta - \delta}}\right).\label{eq:Rnormboundsweakcor}
\end{align}
\end{coro}
As a consequence of Lemma~\ref{lem:onestepweakgauge} we have arrived at~\eqref{eq:H'goal} with $R\in\BS^{\beta+\gamma-\delta}$ and $\beta+\gamma-\delta<\beta$ as desired. One may now iterate the gauge transform to further reduce the order of the error term, starting from $[A]_\Psi$ in the next step. We call such an iterative scheme
\emph{serial gauge transform}. A few remarks on this iterative scheme are in order.
\begin{rema}\leavevmode
\begin{enumerate}\romanenumi
\item At each step of the serial gauge transform, the resonant regions can be chosen differently.
\item \label{item:consecsameres}
Let us consider a serial gauge transform consisting of $k$ steps, starting with the operator $A_0:=A$, and transforming into the operator
\begin{equation}
A_j:=\exp(-\rmi\Psi_j)A_{j-1}\exp(\rmi\Psi_j)=\big[\dots\big[[A]_{\Psi_1}\big]_{\Psi_2}\dots\big]_{\Psi_j}
\end{equation}
at step $j=1,2,\,\dots, k$. Moreover, suppose for simplicity that $\delta>\beta$ and that $\Lambda_{\theta}^{\delta,s}$, $\theta\in G$, are $\delta$-resonant regions for all $A_j^{\CD}$, $j=0,1,\,\dots k-1$, simultaneously (so that the resonant cut-off $\rco =\rco^{\delta, s}$ can be chosen at all steps as $\rco_\theta(\xi)=\bone_{\Xi\setminus\Lambda_{\theta}^{\delta,s}}(\xi), \ (\theta,\xi)\in G\times\Xi$). Then~\eqref{eq:A'} and Corollary~\ref{cor:onestepweakgauge} imply that
\begin{equation}
A_1:=[A]_{\Psi_1}=A_0^{\CD}+A_0^{\CR}+R_1
\end{equation}
with $R_1\in\BS^{2\beta-\delta}$. Repeating the procedure, we obtain
\begin{equation}
A_j=A_{j-1}^{\CD}+A_{j-1}^{\CR}+R_j,\quad j=1,2,\,\dots,\,k,
\end{equation}
with $R_j\in\BS^{\beta+j(\beta-\delta)}$.
\item A disadvantage of the serial gauge transform lies in the fact that already the operator $A_1=[A]_{\Psi_1}$ may have a frequency set as large as
\begin{equation}\label{eq:ZTheta}
Z(\Theta):= \bigcup_{n\,\in\,\N} \Theta^n,
\end{equation}
where
\begin{equation}
\Theta^n := \underbrace{\Theta \cdot \dotso \cdot \Theta}_{\text{product taken } n \text{ times}}.
\end{equation}
This set $Z(\Theta)$ is usually infinite, even when $\Theta$ is finite. Thus, the same holds for $A_1^{\CR}$, which arises after the second step of gauge transform. This might be inconvenient since one generally likes to keep the structure of the resonant operators as simple as possible. However, one can resolve this issue by excluding the terms belonging to $\BS^{\beta+j(\beta-\delta)}$ from $A_j^{\CNR}$ at the $j$th step (by moving them to the remainder). Then $\Theta^{j+1}$ will be the frequency set for $A_j^{\CR}$.
\end{enumerate}
\end{rema}
In the next sub-section, we describe a different iterative gauge transform scheme that we call the \emph{parallel gauge transform}. This is often more convenient to work with than the serial gauge transform.

\subsection{Parallel weak gauge transform} Here, we perform several steps of the gauge transform at the same time, i.e.
\begin{equation}
A^{(\tilde k)}=[A]_{\Psi^{(\tilde{k})}},
\end{equation}
where
\begin{equation}
\Psi^{(\tilde k)}=\sum\limits_{j=1}^{\tilde{k}}\Psi_j
\end{equation}
for some $\tilde{k}\in\N$. Fix again $\delta\in\R$, $s>0$ and a resonant cut-off $\rco^{\delta,s}$ satisfying~\eqref{eq:dennonres} corresponding to $\delta$-resonant regions $\Lambda_\theta^{\delta,s},\ \theta\in G$, for $A^{\CD}$, see Definition~\ref{def:resonant}. Following~\cite[Section~9]{ParSht2012}, the operators $\Psi_l$, $B_l$, and $T_l$ are recursively defined by
\begin{equation}\label{eq:defb1}
B_1:=A^\COD,
\end{equation}
\begin{equation}\label{eq:defbltl}
\begin{split}
B_l:&=\sum\limits_{j=1}^{l-1}\frac{1}{j!}\sum\limits_{k_1+k_2+\dots +k_j=l-1}\ad\left(A^\COD;\Psi_{k_1},\Psi_{k_2},\,\dots,\,\Psi_{k_j}\right),\ l\geq 2,\\
T_l:&=\sum\limits_{j=2}^l\frac{1}{j!}\sum\limits_{k_1+k_2+\dots +k_j=l}\ad\left(A^\CD;\Psi_{k_1},\Psi_{k_2},\,\dots,\,\Psi_{k_j}\right),\ l\geq 2,
\end{split}
\end{equation}
and the relations
\begin{equation}\label{eq:defpsijparalell}
\begin{split}
\ad\left(A^\CD;\Psi_1\right)+B_1^{\CNR}&=0,\\
\ad\left(A^\CD;\Psi_l\right)+B_l^{\CNR}+T_l^{\CNR}&=0, \ l\geq 2.
\end{split}
\end{equation}
More precisely, let $\Theta$ be a frequency set for $A$ and for all $l\geq 1$, let $b_l$ and $t_l$ be the symbols of $B_l$ and $T_l$, respectively. Analogously to~\eqref{eq:defpsiabstract}, we solve~\eqref{eq:defpsijparalell} by choosing $\Psi_l:=\Op(\psi_l)$ with
\begin{align}
(\psi_1)_\theta(\xi)&{:=}
\begin{cases}\dfrac{i\left(b_1^\CNR\right)_\theta(\xi)}{a^\CD(\theta\act\xi)-a^\CD(\xi)}&
\text{if } \theta\in\Theta',\\[2ex] 0 & \text{otherwise}
\end{cases}
\\
\intertext{and}
(\psi_l)_\theta(\xi) &{:=}
\begin{cases}\dfrac{\rmi\left(b_l^\CNR\right)_\theta(\xi)+\rmi\left(t_l^{\CNR}\right)_\theta(\xi)}{a^\CD(\theta\act\xi)-a^\CD(\xi)}&
\text{if } \theta\in(\Theta^l)',\\[2ex] 0 & \text{otherwise}
\end{cases}
\end{align}
for $l\geq 2$. Note that for all $l\geq 1$, $\Theta^l$ is a frequency set for $B_l$, $T_l$, and $\Psi_l$. Finally, put
\begin{align}
Y_{\tilde{k}}&{:=}\sum\limits_{l=1}^{\tilde{k}}B_l+\sum\limits_{l=2}^{\tilde{k}}T_l,\label{eq:defpsiytildek}
\\
\intertext{and}
R_{\tilde{k}+1} &{:=}B_{\tilde{k}+1}+R_{\tilde{k}+1}^{(1)}+R_{\tilde{k}+1}^{(2)},\label{eq:defR}
\end{align}
with
\begin{equation}
\begin{split}
R_{\tilde{k}+1}^{(1)}:&=\sum\limits_{j\,\geq\,\tilde{k}+1}\frac{1}{j!}\ad^j\left(A;\Psi^{(\tilde k)}\right),\\
R_{\tilde{k}+1}^{(2)}:&=\sum\limits_{j=1}^{\tilde{k}}\frac{1}{j!}\sum\limits_{k_1+k_2+\dots + k_j\,\geq\, \tilde{k}+1}\ad\left(A;\Psi_{k_1},\Psi_{k_2},\,\dots,\,\Psi_{k_j}\right).
\end{split}
\end{equation}
Then we arrive at
\begin{equation}\label{eq:afterparallelgauge}
A^{(\tilde k)}=A^\CD+Y_{\tilde{k}}^{\CD}+Y_{\tilde{k}}^{\CR}+R_{\tilde{k}+1},
\end{equation}
see Lemma~\ref{lem:A'series}, where $Y_{\tilde{k}}^{\CR}$ is an operator with frequency set $\Theta_{\tilde{k}}$. The following Proposition provides norm estimates for the operators after the parallel gauge transform. In particular, it shows that if $\delta>\beta$, then we can assure that the error term $R_{\tilde{k}}$ belongs to classes of arbitrarily small order by choosing $\tilde{k}$ sufficiently large.

\begin{prop}\label{prop:symbolestweakgauge}
Let $A^{\CO\CD}\in\BS^\beta$ with $\delta>\beta$. Then we have for all $l\geq 0$,
\begin{equation}
\begin{split}\label{eq:boundspsilbltlparallelweakgauge}
\snorm{\Psi_k}{k(\beta-\delta)}{l}&\ll \left(\snorm{A^\COD}{\beta}{l+n_k}\right)^k,\ k\geq 1\\
\snorm{B_k}{k(\beta-\delta)+\delta}{l}+\snorm{T_k}{k(\beta-\delta)+\delta}{l}&
\ll \left(\snorm{A^\COD}{\beta}{l+n_k}\right)^k,\ k\geq 2,
\end{split}
\end{equation}
where $n_k$ is an increasing function of $k$, depending on $k$, $\beta$ and $\delta$, and the implied constants depend only on $k$, $\beta$, $\delta$, and $s$ in~\eqref{eq:resonantregion}. Moreover, the operators $\Psi^{(\tilde k)}\in\BS^{\beta-\delta}$, $Y_{\tilde{k}}\in\BS^{\beta}$, $R_{\tilde{k}+1}\in \BS^{\tilde{k}(\beta-\delta)+\beta}$ are symmetric and satisfy the bounds
\begin{equation}\label{eq:estpsiafterparallelgauge}
\begin{aligned}
\left\|\Psi^{(\tilde k)}\right\|^{(\beta-\delta)}_{l}+\snorm{Y_{\tilde{k}}}{\beta}{l}&\ll\left(1+\snorm{A^\COD}{\beta}{l+n_{\tilde{k}}}\right)^{\tilde{k}},\\
\snorm{R_{\tilde{k}+1}}{\tilde{k}(\beta-\delta)+\beta}{l}&\leq C_{A,\tilde{k},\beta,\delta,l},
\end{aligned}
\end{equation}
for all $l\geq 0$, where the implied constants only depend on $\tilde{k}, \beta, \delta$, and $s$; and $C_{A,\tilde{k},\beta,\delta,l}$ is a bounded function of the symbol norms $\lbrace \|A^\COD\|^{\beta}_{l}\rbrace_{l\,\geq\,0}$, $\tilde{k}$, $\beta$, $\delta$, and $l$.
\end{prop}
\begin{proof}
The bounds~\eqref{eq:boundspsilbltlparallelweakgauge} are easily deduced from Corollary~\ref{cor:onestepweakgauge} by induction in $k$, estimating all involved commutators using~\eqref{eq:normestcomm}. The estimates on the symbol norms of $\Psi^{(\tilde k)}$ and $Y_{\tilde{k}}$ follow readily.

Let us prove the estimates on the norms of $R_{\tilde{k}+1}$. Starting with $R_{\tilde{k}+1}^{(1)}$ we note that, for $m\geq \tilde{k}+1$ and $\Psi:= \Psi^{(\tilde k)}$
\begin{multline}
\left\|\ad^m(A;\Psi)\right\|^{\left(\tilde{k}(\beta-\delta)+\beta\right)}_l\\
\begin{aligned}
&\leq\left\|\ad^m\left(A^\CD;\Psi\right)\right\|^{\left(\tilde{k}(\beta-\delta) +\beta\right)}_l+\left\|\ad^m\left(A^\COD;\Psi\right)\right\|^{\left(\tilde{k}(\beta-\delta)+\beta\right)}_l\\
&=\left\|\ad^{m-1}\left(Y_{\tilde{k}}^{\CNR};\Psi\right)\right\|^{\left(\tilde{k}(\beta-\delta)+\beta\right)}_l +\left\|\ad^m\left(A^\COD;\Psi\right)\right\|^{\left(\tilde{k}(\beta-\delta)+\beta\right)}_l\\
&\leq 2^{m-\tilde{k}-1}\left\|\ad^{\tilde{k}}\left(Y_{\tilde{k}}^{\CNR};\Psi\right)\right\|^{\left(\tilde{k}(\beta-\delta)+\beta\right)}_l\left(\snorm{\Psi}{0}{l+|\tilde{k}(\beta-\delta)+\beta|}\right)^{m-\tilde{k}-1} \\
& \quad +2^{m-\tilde{k}}\left\|\ad^{\tilde{k}}(A^\COD;\Psi)\right\|^{\left(\tilde{k}(\beta-\delta)+\right)}_l \left(\snorm{\Psi}{0}{l+\left|\tilde{k}(\beta-\delta)+\beta\right|}\right)^{m-\tilde{k}},
\end{aligned}
\end{multline}
where we apply~\eqref{eq:normestcommzeroorder} in the second inequality. Dividing by $m!$ and summing over $m\geq \tilde{k}+1$ we obtain a convergent sum, for which we use the estimates on the norms of $Y_{\tilde{k}}$ and $\Psi^{(\tilde k)}$. Estimating $\|R_{\tilde{k}+1\|^{(2)}}{\tilde{k}(\beta-\delta)+\beta}{l}$ is somewhat easier since there are no convergence issues. This finishes the proof of the Proposition~\ref{prop:symbolestweakgauge}.
\end{proof}

\subsection{Strong gauge transform} \label{sec:sgt}
The aim of any (iterative) gauge transform scheme is to force the error term $R$ after the gauge transform, see e.g.~\eqref{eq:H'goal} or~\eqref{eq:afterparallelgauge}, into a class of (relatively) small order. For instance, in~\eqref{eq:firstorder} and~\eqref{eq:defRconsecgauge}, we aim at $R \in \BS^\gamma$ for some $\gamma < \beta$. If $\Psi$ belongs to a class of negative order, as is the case if one can choose $\delta > \beta$ in Definition~\ref{def:resonant}, for example, it can be trivially satisfied for some $\eta<\beta$, as was seen to be the case with the weak gauge transform. In some cases, however, one can not guarantee more than $\Psi \in \BS^0$, whether it be by choosing $\delta = \beta$, or by introducing additional cut-offs in the definition of $\Psi$. This is notably the case for Schr\"odinger-type operators, whenever the perturbation is not in $\BS^\beta$ for $\beta < \alpha - 1$. In such a case, one can no longer rely on the trivial product estimates for $\ad(A^{\COD};\Psi)$ to get that $\ad(A^{\COD};\Psi)\in\BS^\eta$ for some $\eta<\beta$. In the next lemma, we give a sufficient condition that, nevertheless, yields the required improvement through commuting with $\Psi$.


\begin{lemm}\label{lem:sgt}
Suppose that $A \in \BE\BS^\alpha$ and that $\Psi \in \BS^0$ is defined as in~\eqref{eq:defpsiabstract}. If both $\ad(A^{\CO\CD};\Psi) \in \BS^{\gamma}$ and $\ad(A^{\CNR};\Psi) \in \BS^{\gamma}$, then, with
\begin{align*}
[A]_\Psi =&\, \exp(-i\Psi) A \exp(i\Psi),
\\
\intertext{we have that}
R {:=}&\,[A]_\Psi - A^{\CD} - A^{\CR} \in \BS^{\gamma}.
\end{align*}
\end{lemm}
\begin{rema}
If, in addition to the assumptions of this lemma, we have $A^{\CO\CD} \linebreak\not \in\BS^{\gamma}$, this would mean that commuting with $\Psi$ has improved order and we can therefore call this gauge transform strong. Note that we do not require improvements in order to happen at every iteration of the commutator, but only at the first step.
\end{rema}

\begin{proof}
It follows from Lemma~\ref{lem:A'series} and equation~\eqref{eq:modifiedcomm} that
\begin{equation}
\begin{aligned}
R &= \sum_{k\,\ge\,1} \frac{1}{k!} \ad^k\left(A^{\CO\CD};\Psi\right) + \sum_{k\,\ge\,2} \frac{1}{k!} \ad^{k}\left(A^{\CD},\Psi\right) \\
&= \sum_{k\,\ge\,1} \frac{1}{k!} \ad^{k-1}\left(\ad\left(A^{\CO\CD};\Psi\right);\Psi\right) + \sum_{k\,
\ge\,2} \frac{1}{k!} \ad^{k-1}\left(\ad\left(A^{\CNR};\Psi\right);\Psi\right).
\end{aligned}
\end{equation}
As in the proof of Lemma~\ref{lem:A'series}, both series converge absolutely in $\BS^{\gamma}$.
\end{proof}

\begin{rema}
The hypotheses of the previous lemma can be achieved in many ways. The most common one does not depend on the operator $A$, but only on the algebraic structure of $\BS^\infty$: it is when commutators naturally improve order. The principal example is pseudo-differential operators in $\RL^2(\R^d)$ that are almost periodic with respect to the translation group $\R^d$. To obtain the commutator estimates in this case one requires some limited smoothness in $\xi$. We refer to~\cite[Lemma~3.4]{Sobolev2005} for a proof, and~\cite{MorParSht2014,
ParSob2010,ParSht2012,
ParSht2016,Sobolev2006} for examples of further applications. In all of these cases, the smooth structure of functions on $\R^d$ was used, and the resonance cut-off functions were taken to be smooth approximations to indicator functions of the non-resonant regions rather the indicators themselves.

It is also possible that one cannot reach $\Psi \in \BS^0$ through only non-resonant cut-offs. In such a case, in order to achieve convergence of the series for $\exp(i\Psi)$ and $[A]_\Psi$ achieved in~\eqref{eq:exppsi} and Lemma~\ref{lem:A'series} we will need energy cut-offs, i.e. cutting off large $\xi$. See~\cite{MorParSht2014,ParSob2010} where this idea is being used.
\end{rema}

\begin{rema}
While the weak and the strong gauge transform are defined similarly, the heuristic explanation as to why they work is quite different. One can think of the weak gauge transform as a very sophisticated perturbation theory. Indeed, if $\Psi$ is of negative order, its norm is small and so $\exp(i\Psi)$ is a small perturbation of the identity. This means that perturbations are additive in the first order, and the gauge transform is a convenient way of doing the bookkeeping. On the other hand, the strong gauge transform works due to certain algebraic structure present in the problem. As a toy example, consider the operator $A=A_0+B$ acting in $\ell^2(\Z)$, where $A_0$ is diagonal with $(A_0)_{jj}=j$ for $j\in\Z$, and $B$ is self-adjoint and T\"oplitz (and, for simplicity, bounded). Define $\Psi$ to be the off-diagonal matrix defined by
\[
\Psi_{jk}=
\begin{cases}\frac{iB_{jk}}{k - j},& \ j \ne k\\
0,& \ j= k
\end{cases}.
\]

Then, $\Psi$ is also self-adjoint and T\"oplitz. Since T\"oplitz operators commute~\eqref{LP3} immediately implies that $\exp(i\Psi)A\exp(-i\Psi)=A_0$. Note that no smallness of $B$ (and thus $\Psi$) is assumed here.
\end{rema}

\section{Systems of Almost Periodic Operators} \label{sec:mapo}

In this section, we provide a construction suitable to describe almost periodic operators with matrix-valued symbols within the framework of Sections~\ref{sec:operators} --~\ref{sec:gt}.

\subsection{Symbol formalism for systems of almost periodic operators}
Let the index set $\Xi$ and the group $G$ be as in Section~\ref{sec:operators}. Let $m\in\N$ and $\bbb:G\times\Xi\to\CL(\C^m)$ be a function such that there exists a countable frequency set $\Theta=\Theta^{-1}\subseteq G$ with $\bbb_\theta(\xi)=0$ for all $\theta\in G\setminus \Theta$ and $\xi\in\Xi$. Furthermore, assume that
\begin{equation}
\sum\limits_{\theta\,\in\,\Theta}\left\|\bbb_\theta(\xi)\right\|^2 <\infty,\ \text{for all } \xi\in\Xi,
\end{equation}
where $\|\,\cdot\,\|$ is the operator norm on $\CL(\C^m)$. For every $\xi\in\Xi$, let $\set{v_j(\xi):\ j\in \Z/m\Z}$ be an orthonormal basis for $\C^m$ so that $\lbrace\be_\xi\otimes v_j(\xi)\rbrace_{\xi\,\in\,\Xi,\,j\,\in\,\Z/m\Z}$ is an orthonormal basis for $\ell^2(\Xi;\C^m)=\ell^2(\Xi)\otimes\C^m$. In analogy to~\eqref{eq:defopB}, an almost periodic operator $\BB$ in $\ell^2(\Xi;\C^m)$ with symbol $\bb$ is defined by
\begin{equation}\label{eq:Basmatrixop}
\BB\big(\be_\xi \otimes v_j(\xi)\big):=\sum_{\theta\,\in\,\Theta}\be_{\theta\act\xi} \otimes \big[\bbb_{\theta}(\xi)v_j(\xi)\big].
\end{equation}
We introduce the index set $\Xi_m := \Xi \times \Z/m\Z$ equipped with the weight function
\begin{equation}\label{eq:angxim}
\ang{(\xi,j)}_m:= \ang \xi
\end{equation}
and define the group $G_m:= G \times \Z/m\Z$, and its (free) action on $\Xi_m$ by
\begin{equation}\label{eq:actiongprime}
(g,k)\act(\xi,j):= (g\act \xi,k+j).
\end{equation}
Applying the unitary map $T_m:\ell^2(\Xi_m)\to\ell^2(\Xi;\C^m)$ defined by
\begin{equation}\label{eq:defTm}
T_m\big(\be_{(\xi,j)}\big):= \be_\xi\otimes v_j(\xi),\ (\xi,j)\in \Xi_m,
\end{equation}
we can relate the operator $\BB$ to the operator
\begin{equation}\label{eq:defBfromBB}
B:=T_m^\ast \BB T_m
\end{equation}
in $\ell^2(\Xi_m)$. For every $g\in G$ and $\xi\in\Xi$, let $[\bbb_g(\xi)]$ be the matrix representation of $\bbb_g(\xi):\C^m\to\C^m$ with respect to the pair of bases $\set{v_j(\xi)}_j$ and $\set{v_j(g\act \xi)}_j$ on the domain and codomain, respectively, i.e.
\begin{equation}\label{eq:bbbonbasis}
\bbb_g(\xi)v_j(\xi)=\sum\limits_{k\,\in\,\Z/m\Z}\big[\bb_g(\xi)\big]_{kj}v_k(g\act\xi), \quad \text{for all} \ j\in \Z/m\Z.
\end{equation}
Define the scalar symbol $b:G_m\times\Xi_m\to\C$ by
\begin{equation}\label{eq:symbolgprime}
b_{(g,k)}(\xi,j):=\left[\bbb_g(\xi)\right]_{(k+j)j}
\end{equation}
for $(g,k)\in G_m$ and $(\xi,j)\in\Xi_m$, and note that $\Theta\times \Z/m\Z$ is a frequency set for $b$. In view of~\eqref{eq:defopB}, \eqref{eq:actiongprime}, and~\eqref{eq:symbolgprime}, we have
\begin{equation}\label{eq:actionmapo}
\Op(b) \be_{(\xi,j)} = \sum_{(\theta,k)\,\in\,\Theta \times \Z/m\Z} \left[\bbb_\theta(\xi)\right]_{kj} \be_{(\theta\,\act\,\xi,k)}, \quad \text{for} \ (\xi,j)\in\Xi_m.
\end{equation}
Hence, \eqref{eq:Basmatrixop}, \eqref{eq:defTm}, \eqref{eq:defBfromBB}, \eqref{eq:bbbonbasis}, and~\eqref{eq:actionmapo} yield $T_m^* \BB T_m = B = \Op(b)$. This justifies calling the operators from $\BS^\infty(G_m,\Xi_m)$ \emph{systems of almost periodic operators}, also known as matrix-valued operators. We shall use the notation
\begin{equation}\label{eq:defsystemclasses}
\BT_m^\gamma:=\BT^\gamma(G_m,\Xi_m),\ \BT\in\{\BS, \BD\BS, \BD\BE\BS, \BS\BE\BS, \BE\BS \}, \ \gamma\in\R\cup\set{\pm\infty},
\end{equation}
and
\begin{equation}
\sfH_m^\gamma:=\sfH^\gamma(\Xi_m),\ \gamma\in\R\cup\set{\pm\infty}.
\end{equation}

Since the map of symbols $\bbb\mapsto b$ is one-to-one, for those $\bbb$ which are mapped to $b\in\BS_m^\infty$ we write $\Op(\bbb):=\Op(b)$. We use this identification to apply the results of Sections~\ref{sec:operators}--\ref{sec:gt} without always making explicit the conjugation by the operators $T_m$. Note that $\ang{(g,k)}_m=\ang{g}$, for all $(g,k)\in G_m$, see~\eqref{eq:modulusg} and~\eqref{eq:angxim}. Hence, for $b\in\BS_m^\beta$, we have the equivalence of norms,
\begin{equation}\label{eq:matrixseminorm}
\Fc_m^{-1}\snorm{b}{\beta}{l}\leq\sum_{\theta\,\in\,\Theta} \ang \theta^l \sup_{\xi\,\in\,\Xi} \ang\xi^{-\beta}\norm{\bbb_{\theta}(\xi)}\leq \Fc_m\snorm{b}{\beta}{l},
\end{equation}
where the constant $\Fc_m>0$ only depends on $m$. Two sub-algebras of $\BS_m^\infty$ will be of particular interest in the sequel: uncoupled operators and diagonal operators.

\begin{defi}\label{def:uncoupled}
The \emph{uncoupled operators} in $\BS_m^\beta$, $\beta\in\R\cup\set{\pm\infty}$, are defined by
\begin{multline}
\BU\BS_m^\beta\\
\begin{aligned}
:=&\,\big\{\BB\in\BS_m^\beta: \ &\text{the matrix } \big[\bbb_g(\xi)\big], \text{ see~\eqref{eq:bbbonbasis}, is diagonal for all }g\in G,\ \xi\in\Xi\big\}\\
=&\,\big\{\BB\in\BS_m^\beta: \ &\text{the frequency set for }b\text{ can be chosen as a subset of }G\times\{0\}\big\}.
\end{aligned}
\end{multline}
For any operator $\BA =\Op(\ba)\in \BS_m^\alpha$, $\alpha\in\R\cup\set{\pm\infty}$, we define the symbol
\begin{equation}\label{eq:uncoupled}
\big[\ba_g^\CU(\xi)\big]_{k j}:=
\begin{cases}
\big[\ba_g(\xi)\big]_{k j} & \text{if }k =j\\ 0 & \text{if } k\neq j
\end{cases},\quad\text{ for all }g\in G, \xi\in \Xi, k,j\in\Z/m\Z.
\end{equation}
We write $\BA^\CU := \Op(\ba^\CU)$ for the projection of $\BA$ onto $\BU\BS_m^\alpha$ which we call the uncoupled part. We also denote $\BA^{\CC} := \BA - \BA^{\CU}$ its coupled part. It can easily be seen that if $\BA \in
\BS_m^\gamma$, $\gamma \in \R$ then for all $l \ge 0$
\begin{equation}\label{eq:estpartuncbyfull}
\snorm{\BA^\CU}{\gamma}{l} \le \snorm{\BA}{\gamma}{l}, \quad
\snorm{\BA^\CC}{\gamma}{l} \le \snorm{\BA}{\gamma}{l}
\end{equation}
and that if $\BA$ is symmetric so are $\BA^\CU$ and $\BA^\CC$.
\end{defi}

The second sub-algebra is $\BD\BS_m^\infty$, see~\eqref{eq:defsystemclasses} and Definition~\ref{def:dapo}. Noting that $\id_{G_m}=(\id_G,0)$, we infer from~\eqref{eq:symbolgprime} that
\begin{equation}
\BD\BS_m^\infty=\big\{B\in\BU\BS_m^\infty:\ \bbb_g(\xi)=0\ \text{for all} \ g\in G\setminus \set{\id_G}\big\},
\end{equation}
so that $\BD\BS_m^\infty \subseteq \BU\BS_m^\infty \subseteq \BS_m^\infty$. As in the scalar case, for any operator $\BA =\Op(\ba)\in \BS_m^\alpha$, $\alpha\in\R\cup\set{\pm\infty}$, we denote by
\begin{equation}\label{eq:diagonalmatrix}
\begin{split}
\BA^\CD &:=\Op(\ba^\CD)\quad\\ \big[\ba_g^\CU(\xi)\big]_{k j}&:=
\begin{cases}
\big[\ba_g(\xi)\big]_{k j} & \text{if }k =j
\text{ and } g = \id\\ 0 &
\text{otherwise}.
\end{cases}
\end{split}
\end{equation}
and $\BA^{\CO\CD} = \BA - \BA^\CD$. This definition makes it so that $\BA^\CD = T_m A^\CD T_m^*$. Similarly, if resonant and non-resonant regions are defined in terms of $\Xi_m$, we set $\BA^\CR$ = $T_mA^\CR T_m^*$ and $\BA^{\CNR} = T_m A^{\CNR} T_m^*$. We can also combine notions of coupling and resonance; we set for instance $\BA^{\CR,\CU} = (\BA^{\CR})^\CU$, and proceed similarly for other combinations of the indices.



The following lemma is useful when changing the reference orthonormal basis of $\C^m$. Nevertheless, for the rest of this section the reference basis of $\ell^2(\Xi;\C^m)$ will remain fixed as $\{\be_\xi\otimes v_j(\xi)\}_{(\xi,j)\,\in\,\Xi_m}$.
\begin{lemm}
Assume that, for every $\xi\in\Xi$, the set $\set{u_j(\xi):\ j\in \Z/m\Z}$ is an orthonormal basis for $\C^m$. Then the unitary operator
\begin{equation}\label{eq:defBU}
\BU:\ell^2\left(\Xi;\C^m\right) \to \ell^2\left(\Xi;\C^m\right), \quad \be_\xi\otimes v_j(\xi)\mapsto \be_\xi\otimes u_j(\xi)
\end{equation}
satisfies $T_m^*\BU T_m\in\BS_m^0$.
\end{lemm}
\begin{proof}
It is clear from~\eqref{eq:defBU} that $T_m^*\BU T_m=\Op(\bu)$ where $\bu_g(\xi)\in \operatorname{U}(m)$ is unitary for all $g\in G$, $\xi\in\Xi$, and $\set{\id_G}$ is a frequency set for $\bu$. Thus the equivalence of norms~\eqref{eq:matrixseminorm} implies that $T_m^*\BU T_m\in\BS_m^0$.
\end{proof}

The previous lemma has the following corollary, justifying our terminology of uncoupled operators.
\begin{coro}
Let $\set{v_j:\ j\in\Z/m\Z}$ be a fixed basis for $\C^m$. Then any operator $A \in \BU\BS_m^\infty$ is unitarily equivalent to an orthogonal sum $\bigoplus_{j\,\in\,\Z/m\Z}A_j$ where for every $j\in\Z/m\Z$, $A_j$ acts in $\ell^2(\Xi) \otimes \spann\set{v_j}$.
\end{coro}


\subsection{Gauge transform in \texorpdfstring{$\BS_m^\infty$}{Sm}: the reduction to uncoupled operators} \label{sec:reducscalar}

We would like to apply a weak gauge transform to an operator in the class $\BS\BE\BS^\infty_m$\linebreak --- cf.~\eqref{eq:defsystemclasses} --- in order to obtain an operator of the same order with an uncoupled principal symbol. In this section, we give two sufficient conditions that allow us to do this. The first one is more restrictive on the off-diagonal part and gives a non-trivial remainder, but allows for a principal symbol with multiple eigenvalue. The second one requires the principal symbol to have only simple eigenvalues, in which case the procedure is more efficient and the restrictions on the off-diagonal symbol are much milder.

\begin{theo}\label{thm:systemonestep}
Let $\BA = \Op(\ba) \in \BS\BE\BS_m^\alpha$ be symmetric and let $\beta <
\alpha$ be such that $\BA^{\CC}:= \Op(\ba^{\CC}) \in \BS^\beta_m$. Assume that
\begin{equation}
\big[\ba_{\id}\big]\left(\xi \right) = \ang \xi^\alpha
\diag(a_1(\xi),\,\dotsc,\,a_m(\xi))
\end{equation}
and that $\Theta$ is a frequency set for $\ba^\COD$. Here, for $j \in \Z/m\Z$, $a_j : \Xi \to \R$ are bounded functions such that for all $\theta \in Z(\Theta) = \bigcup_{k=1}^\infty \Theta^k$,
\begin{equation}\label{eq:shftinfty}
\lim_{\ang \xi\,\to\,\infty} \frac{a_j(\theta \act \xi)}{a_j(\xi)} = 1.
\end{equation}
Suppose finally that there exists $C, c> 0$ such that for every $j \in \Z/m\Z$ and $k \in \Z / m \Z
\setminus \set{0}$, either
\begin{equation}\label{eq:ajboundfrombelow}
\inf_{\ang \xi\,>\,C} \abs{a_j(\xi) - a_{j+k}(\xi)} \ge c > 0,
\end{equation}
or
\begin{equation}\label{eq:hypresonant}
\left[\ba^{\COD}\right]_{j,j+k} \in \BS^{2\beta - \alpha}.
\end{equation}
Then, for all $\eps > 0$ and $N \in \N$ there exists a symmetric operator $\BPsi \in \BS_m^{\beta - \alpha}$ such that
\begin{equation}\label{eq:desiredafteronestep}
[\BA]_{\BPsi} = \exp(-i\BPsi) \BA \exp(i\BPsi) = \BA^\CD + \BY + \BR_1 + \BR_2
\end{equation}
where $\BY \in \BU\BS_m^\beta$, $\BR_1 \in \BS^{2 \beta - \alpha}$, $\norm{\BR_2}_{\sfH^\beta_m \to \sfH_m^0} < \eps$ and $\BY, \BR_1,
\BR_2$ are symmetric. If $\BA$ is quasi-periodic, one can choose $\BR_2 = 0$.
\end{theo}

\begin{rema}
The conditions~\eqref{eq:ajboundfrombelow} and~\eqref{eq:shftinfty} are satisfied in the simple case of constant functions $a_j(\xi)=a_j\in\R\setminus
\{0\}$.
\end{rema}
\begin{proof}
Fix $\epsilon' > 0$. We first eliminate the long-range coupling. Since $\|\BA^{\COD}\|^{\beta}_{0} < \infty$, there exists a finite subset $\tilde{\Theta} \subseteq \Theta$, closed under inversion and containing the identity, such that
\begin{equation}\label{eq:smallnorm}
\sum_{\theta\,\in\,\Theta \setminus \tilde \Theta}\,
\sup_{\xi\,\in\,\Xi} \ang \xi^{-\beta}\norm{\ba_\theta^{\COD}(\xi)} < \epsilon'.
\end{equation}
Let $\BB:= \Op(\bbb)$ with the symbol
\begin{equation}\label{eq:btilde}
\bbb_\theta(\xi) :=
\begin{cases}
\ba_\theta^{\COD}(\xi) & \text{if } \theta \in \tilde \Theta, \\
0 & \text{otherwise.}
\end{cases}
\end{equation}
For $\tilde \BR: = \BA^{\COD} - \BB$, \eqref{eq:matrixseminorm} implies
\begin{equation}\label{eq:estRtilde}
\big\|\widetilde{\BR}\big\|_{0}^{(\beta)} < \Fc_m\epsilon',
\end{equation}
and we write $\widetilde{\BA}:= \BA^\CD + \BB$ so that $\BA=\widetilde{\BA}+\widetilde{\BR}$ and $\tilde \BA^\CD = \BA^\CD$. For every $j \in \Z/m\Z$, define the set
\begin{equation}\label{eq:specificconstants}
I_j := \set{k\,\in\,\Z/m\Z :~\eqref{eq:ajboundfrombelow} \text{ holds}}.
\end{equation}
Finiteness of $\tilde \Theta$ and bounded range of action imply that
\begin{equation}
\lim_{\ang \xi\,\to\,\infty} \sup_{\theta\,\in\,\tilde \Theta} \abs{\frac{\ang{\theta
\act \xi}}{\ang \xi} -1} = 0.
\end{equation}
Combining this with~\eqref{eq:shftinfty} and~\eqref{eq:ajboundfrombelow}, as well as boundedness of the functions $a_j$ implies the existence of $s'$ depending on $\eps'$ and the constants $c, C$ in~\eqref{eq:ajboundfrombelow} such that
\begin{equation}
\inf_{j\,\in\,\Z / m \Z} \; \inf_{k\,\in\,I_j} \; \inf_{\theta\,\in
\,\tilde \Theta} \;
\inf_{\ang \xi\,>\,s'} \Big|a_{j+k}(\theta \act \xi) \ang{\theta \act \xi}^\alpha - a_j(\xi) \ang \xi^\alpha\Big| > \frac c 2 \ang \xi^\alpha.
\end{equation}
Thus, for $(g,k) \in G_m$, the sets
\begin{equation}\label{eq:resregionsystem}
\Lambda_{(g,k)}^{\alpha,c/2} :=
\begin{cases}
\big\{(\xi,j) \in \Xi_m : \min\set{\ang \xi, \ang{g \act \xi}} \le s'\big\}, &
\text{if } g \in \tilde \Theta \text{ and } k \in I_j, \\
\Xi_m, & \text{otherwise}
\end{cases}
\end{equation}
are $\alpha$-resonant regions for the operator $\tilde A = T_m^* \BA T_m$, cf.~\eqref{eq:resonantregion}, and we choose the corresponding (scalar) resonance cut-off function
\begin{equation}
\chi_{(g,k)}(\xi,j):=\bone_{\Xi_m\setminus\Lambda_{(g,k)}^{\alpha,c/2}(\xi,j)},
\end{equation}
see Remark~\ref{rem:res}$\MK$\eqref{item:remrescutoff}. Thus, taking $\Psi$ as in Lemma~\ref{lem:psi}, we have that $\BPsi = T_m \Psi T_m^* \in \BS_m^{\beta - \alpha}$. In view of~\eqref{eq:A'}, we deduce that
\begin{equation}
[\tilde \BA]_{\BPsi} = \exp(-i \BPsi) \tilde \BA \exp(i \BPsi) = \BA^\CD +
\tilde \BA^\CR + \BR,
\end{equation}
where Corollary~\ref{cor:onestepweakgauge} and conjugation by $T_m$ give $\BR \in \BS_m^{2 \beta -
\alpha}$. We turn our attention to $\tilde \BA^\CR$. We decompose it as $\tilde \BA^\CR =
\tilde \BA^{\CR,\CU} + \tilde \BA^{\CR,\CC}$. By definition of the resonant region $\Lambda^{\alpha,c/2}_{(g,k)}$, we have that
\begin{equation}
\left[\tilde \ba^{\CR,\CC}\right]_{j,j+k} = \left[\tilde \ba^{\CO\CD,\CC}\right]_{j,j+k} \quad \text{if } k \not \in I_j
\end{equation}
and
\begin{equation}
\supp\left(\left[\ba^{\CR,\CC}\right]_{j,j+k}\right) \subset \set{\xi \in \Xi : \min_{\theta\,\in\,\tilde \Theta} \ang{\theta \act \xi} \le s'} \quad \text{if } k \in I_j.
\end{equation}
By~\eqref{eq:hypresonant}, for every $k \not \in I_j$ we have $[\tilde
\ba^{\CR,\CC}]_{j,j+k} \in \BS^{2\beta - \alpha}$. Finiteness of $\tilde
\Theta$ and bounded range of action imply that the support of $[\ba^{\CR,\CC}]_{j,j+k}$ is bounded for $k \in I_j$. Together, along with Proposition~\ref{prop:normorder}, this gives $\tilde \BA^{\CR,\CC} \in \BS_m^{2 \beta - \alpha}$.


All of this implies
\begin{equation}\label{eq:decompo}
[\BA]_{\BPsi} = \BA^\CD + \tilde \BA^{\CR,\CU} + \tilde \BA^{\CR,\CC} + \BR +
\exp(-i\BPsi) \tilde
\BR \exp(i \BPsi).
\end{equation}
We claim that this has the desired form~\eqref{eq:desiredafteronestep} with
\begin{equation}\label{eq:decompodetails}
\BY =
\tilde \BA^{\CR, \CU}, \qquad \BR_1 = \tilde \BA^{\CR,\CC} + \BR, \qquad
\text{and} \qquad \BR_2 =
\exp(-i \BPsi)
\tilde \BR \exp(i\BPsi).
\end{equation}
Indeed, it follows from~\eqref{eq:estpartbyfull} and~\eqref{eq:estpartuncbyfull} that $\tilde \BA^{\CR,\CU} \in \BU\BS^\beta_m$. We have that $\tilde \BA^{\CR,\CC}, \BR \in \BS_m^{2\beta - \alpha}$ so that their sum $\BR_1$ also is.

Finally, recall that $\exp(i\BPsi)\in\BS^0_m$ by Corollary~\ref{cor:complete}. In particular, we have
\begin{equation}
\snorm{\exp(i\BPsi)}{0}{|\beta|}\leq \sum\limits_{k=0}^\infty
\frac{\big(\snorm{\BPsi}{0}{|\beta|}\big)^k}{k!},
\end{equation}
where by Corollary~\ref{cor:onestepweakgauge} and conjugation with $T_m$ we have
\begin{equation}
\begin{aligned}
\snorm{\BPsi}{0}{|\beta|}\leq \snorm{\BPsi}{\beta-\alpha}{|\beta|}\le
\frac{4}{c}
\big\|\BB\big\|_{|\beta|}^{(\beta)} \leq \frac 4 c
\snorm{\BA^{\COD}}{\beta}{|\beta|}
\end{aligned},
\end{equation}
with $c$ the constant in~\eqref{eq:ajboundfrombelow}. Consequently, $\snorm{\exp(i\BPsi)}{0}{|\beta|}$ is bounded uniformly in $\epsilon'\searrow 0$. Hence, for any $\epsilon>0$, choosing
\begin{equation}
0< \eps' < \frac{\eps}{\left(\snorm{\exp(\rmi\BPsi)}{0}{|\beta|}\right)^2
\Fc_m},
\end{equation}
we obtain by Lemma~\ref{lem:domain} and~\eqref{eq:estRtilde} that
\begin{equation}\label{eq:normremainder}
\begin{aligned}
\norm{\BR_2}_{\RH_m^{\beta}\,\to\,\RH_m^0}&\leq \norm{\exp(i\BPsi)}_{\RH_m^0\,\to\,\RH_m^0}\big\|\widetilde{\BR}\big\|_{\RH_m^{\beta}\,\to\,\RH_m^0}\norm{\exp(i\BPsi)}_{\RH_m^{\beta}
\,\to\,\RH_m^{\beta}}\\
&\leq \big\|\widetilde{\BR}\big\|_{0}^{(\beta)}\left(\snorm{\exp(i\BPsi)}{0}{|\beta|}\right)^2<\epsilon.
\end{aligned}
\end{equation}
This finishes the proof.
\end{proof}

\begin{theo}\label{thm:gtsystem}
Let $\BA = \Op(\ba) \in \BS\BE\BS_{m}^\alpha$ be symmetric and let $\beta<\alpha$ such that $\BA^{\CO\CD}:=\Op(\ba^{\COD}) \in
\BS^\beta_m$ with frequency set $\Theta \subset G$. Assume that
\begin{equation}
\big[\ba_{\id}(\xi)\big]=\ang{\xi}^\alpha\diag(a_1(\xi),\,\dotsc,\,a_m(\xi))
\end{equation}
for some bounded functions $a_j:\Xi\to\R$. Moreover, suppose that there exist $C, c>0$ such that
\begin{equation}\label{eq:ajboundfrombelowbis}
\inf\limits_{\ang \xi >C}\min\limits_{j\neq k} |a_j(\xi)-a_k(\xi)|\geq c>0,
\end{equation}
and that, for all $j=1,2,\,\dots,\,m$, and $\theta\in Z(\Theta)=\bigcup\limits_{k=1}^\infty \Theta^k$,
\begin{equation}\label{eq:shiftsatinfinity}
\lim\limits_{\ang{\xi}\,\to\,\infty} \frac{a_j(\theta\act\xi)}{a_j(\xi)}=1.
\end{equation}
Then, for all $\epsilon>0$ and $N \in \N$ there exists a symmetric operator $\BPsi\in\BS^{\beta-\alpha}_m$ such that
\begin{equation}\label{eq:matgt}
[\BA]_{\BPsi} = \exp(i \BPsi) \BA \exp(i\BPsi) = \BA^{\CD} + \BY^{\CU} +
\BR_1 + \BR_2
\end{equation}
with $\BY\in\BS^\beta_m$, $\BR_1 \in \BS_m^{-N}$, $\norm{\BR_2}_{\sfH_m^{\beta}\,\to\,\sfH_m^{0}} < \epsilon$, and $\BY$, $\BR_1$, $\BR_2$ symmetric. If $\BA^{\COD}$ is quasi-periodic, then one can choose $\BR_2=0$.
\end{theo}
\begin{proof}
The proof essentially follows the scheme of the proof of Theorem~\ref{thm:systemonestep}. We first eliminate long-range coupling and find $\BB \in \BS_m^\beta$ and $\widetilde \BR$ such that $\BA = \BA^\CD +
\BB +
\widetilde \BR$ and
\[
\snorm{\widetilde \BR}{\beta}{0} < \eps'.
\]
Assumption~\eqref{eq:ajboundfrombelowbis} leads this time to $\alpha$-resonant regions
\begin{equation}\label{eq:resmultstep}
\Lambda^{\alpha,c/2}_{(g,k)} =
\begin{cases}
\set{(\xi,j) \in \Xi_m : \min\set{\ang \xi, \ang{g \act \xi} } \le s'}, &
\text{if } k \ne 0, \\
\Xi_m, & \text{if } k = 0
\end{cases}
\end{equation}
for some $s'$ depending on $\eps'$. Put
\begin{equation}
K := \frac{N + \beta}{\alpha - \beta}.
\end{equation}
We apply a parallel weak gauge transform according to~\eqref{eq:afterparallelgauge}. We have from Proposition~\ref{prop:symbolestweakgauge} and conjugating by $T_m$ that there exists symmetric operators $\BPsi \in \BS_m^{\beta - \alpha}$, $\BY \in \BS_m^\beta$ and $\BR
\in \BS_m^{-N}$ such that
\begin{equation}
[\tilde \BA]_{\BPsi} = \exp(-i \BPsi) \tilde A \exp(i \BPsi) = \BA^\CD + \BY^\CD + \BY^\CR + \BR
\end{equation}
and
\begin{equation}\label{eq:boundsnormpsi}
\snorm{\BPsi}{\beta - \alpha}{\abs{\beta}} \ll \left(1 +
\snorm{\BA^{\COD}}{\beta}{n_K}\right)^K.
\end{equation}
where the implicit constant depends only on $c$ and $K$. Inequality~\eqref{eq:boundsnormpsi} implies that $\|\exp(i\BPsi)\|_{\abs \beta}^{(0)}$ is uniformly bounded as $\eps' \searrow0$.

With the resonant region as in~\eqref{eq:resmultstep}, for every $\BY
\in \BS_m^\gamma$, $\gamma < \alpha$ we have that $
\BY^\CD+ \BY^\CR = \BY^\CU + \BR_\BY$, where the symbol of $\BR_\BY$ has bounded support, implying $\BR_\BY \in \BS_m^{-\infty}$. We put $\BR_1 = \BR + \BR_\BY \in
\BS^{-N}_m$ and $\BR_2 = \exp(-i\BPsi) \tilde \BR \exp(i \BPsi)$ gives us
\begin{equation}
\norm{\BR_2}_{\RH_m^\beta\,\to\,\RH_m^0} \le \eps'
\left(\snorm{\exp(i\BPsi)}{\abs
\beta}{0}\right)^2.
\end{equation}
Therefore for any $\eps > 0$ by choosing $ 0 < \eps' < \eps
\left(\snorm{\exp(i\BPsi)}{\abs
\beta}{0}\right)^{-2}$ we obtain
\begin{equation}
[\BA]_{\BPsi} = \BA^\CD + \BY^\CU + \BR_1 + \BR_2
\end{equation}
with the claimed properties.
\end{proof}



\section*{\textsc{Part II : Applications to asymptotic properties of systems}}


In this second part, we consider some specific examples where the methods and results developed in the first half are applicable. As was mentioned earlier, these methods work very well for operators
\[
H = H_0 + B
\]
of Schr\"odinger type acting on $\RL^2(\R^d)$. Here, $H_0 = (-\Delta)^{\alpha/2}$ and $B$ is a pseudo-differential perturbation of order $\beta < \alpha$. In particulars, the gauge transform method allows us to solve the following two types of problems, see~\cite{barbpar,MorParSht2014,parnovski,
ParSht2016,ParSht2012,ParSob2010,
Sobolev2006}:
\begin{itemize}
\item obtain a complete asymptotic expansion for the integrated density of states of almost periodic operators, as the spectral parameter goes to infinity;
\item prove that some elliptic periodic operators have the Bethe--Sommerfeld property, which asserts that the spectrum of such operators contains a half-line $[\lambda;\infty)$ for some $\lambda \in \R$.
\end{itemize}

We now consider these questions in the setting of elliptic systems of operators. We establish answers to both of these problems in the case where symbols are periodic, for the Bethe--Sommerfeld property, and almost periodic, for the integrated density of states. We will do so by using the tools developed in Part I of this paper to reduce these operators to uncoupled operators. We will show that such a reduction cannot change the integrated density of states too much, and we will show that it cannot open infinitely many gaps in the spectrum. Since elliptic systems of operators do not have to be semi-bounded, we will obtain these results as the spectral parameter goes to $\pm
\infty$. In order to do this, we will establish quantitative estimates based upon the results of Sections~\ref{sec:ids}, \ref{sec:gt} and~\ref{sec:mapo} under generic assumptions about the perturbations.

In Section~\ref{sec:besicovitch}, we describe the Besicovitch space of almost periodic functions and the operators acting on it. In particular, in Section~\ref{sec:prelim}, we describe the generic conditions required to prove the existence of complete asymptotics for the integrated density of states (IDS).
 
In Section~\ref{sec:asyexp}, we state and prove the main Theorems~\ref{thm:aexpconcretecut} and~\ref{thm:aexpconcrete} describing the asymptotic behaviour of the IDS for elliptic
systems of operators. Recall that they are more general versions of Theorems~\ref{thm:dirac2dids} and~\ref{thm:dirac3dids}. These theorems are proved by performing various reductions; in turn to a finite interval of the spectral parameter, to quasiperiodic operators and then to uncoupled operators, using the gauge transform.
 
From Sections~\ref{sec:periodic} to~\ref{sec:cg}, we change perspective and we study periodic operators. In Section~\ref{sec:periodic} we describe the structure of such operators, interpreting the Bloch--Floquet decomposition through the lens of almost periodic functions. We also introduce an auxiliary tool useful in the study of the Bethe—Sommerfeld property, the spectral overlap function.
 
In Section~\ref{sec:bs}, we give conditions under which elliptic systems of periodic operators enjoy the Bethe–Sommerfeld property. We then use the reduction to uncoupled operators and bounds for the density of states obtained in Section~\ref{sec:asyexp} to show that it is sufficient to prove that the spectral overlap function is sufficiently bounded away from 0 for uncoupled operators. This will be done by reusing the results of Section~\ref{sec:ids}, but interpreting fibrewise eigenvalue counting functions as instances of the IDS.
 
We prove those lower bounds in Section~\ref{sec:cg} by refining arguments based on combinatorial geometry that were previously used in proving the Bethe–Sommerfeld conjecture for Schr\"odinger-type operators.


Finally in Section~\ref{sec:dirac}, we spend a few words to show that periodic and almost periodic perturbations of the Dirac operator fit in the framework that we have described in this part.


\section{Besicovitch space and systems of operators} \label{sec:besicovitch}

In this section, we turn back to the space $\RB^2(\R^d;\C^m)$ of almost periodic vector-valued functions, corresponding to the case where $G = \Xi = \R^d$ and
\begin{equation}
\bg_1\bg_2:=\bg_1+\bg_2, \quad \bg\act\bxi:=\bg+\bxi,
\end{equation}
for all $\bg_1,\bg_2,\bxi\in\R^d$. The weight function is $\ang \bxi = 1 + \abs \bxi$. From~\eqref{eq:modulusg} we also get that $\ang \bg = 1 + \abs
\bg$ and that $G$ has bounded range of action. The case $m = 1$ corresponds to the usual Besicovitch space. We now offer a concrete description of this space, along with a few results relating the properties of operators acting on $\RL^2$ and $\RB^2$. These results can be found in~\cite{CoburnMoyerSinger,Shubin1978,Shubin1979}.

Let $\set{v_1,\,\dotsc,\,v_m}$ be an orthonormal basis for $\C^m$ and for $1 \le j
\le m$ let
\begin{equation}
\be_{\bxi,j}(\bx):= \exp(i \bxi \cdot \bx) \otimes v_j.
\end{equation}
The space $\RB^2(\R^d;\C^m)$ is the closure of
\begin{equation}
\spann \set{\be_{\bxi,j}:\ \bxi\in\R^d,\ j=1,\,\dots,\,m},
\end{equation}
taken with respect to the inner product
\begin{equation}
(f,g)_{\RB^2} = \lim_{L\,\to\,\infty} \frac{1}{(2L)^d}\int_{[-L,L]^d} f \cdot \bar g \de \bx.
\end{equation}

For the remainder of this article, we will use $\BS_m^\infty$, $\BD\BS_m^\infty$, etc. to refer to the spaces of almost periodic operators acting on $\RB^2(\R^d;\C^m)$. Let $\BA$ be an operator in $\BS_m^\alpha$ with symbol $\ba(\bx,\bxi)$. The action of $\BA$ in $\RL^2(\R^d;\C^m)$ as an operator in the H\"ormander class $\Psi^\alpha(\R^d;\C^m)$ with almost periodic symbol is defined through the usual Fourier integral representation of pseudo-differential operators (see e.g.~\cite{HormanderIII}) as
\[
\widetilde {\operatorname{Op}}(\ba)f(\bx) = \frac{1}{(2\pi)^d}\iint_{\R^d \times \R^d} \exp(i \bxi \cdot(\bx-\by))
\ba(\bx,\bxi) f(\by) \de \by \de \bxi.
\]
The following proposition links its properties as an operator in $\RL^2$ and $\RB^2$, respectively.

\begin{prop}\label{prop:normspec}
If $\BA \in \BS_m^\infty$ is bounded or elliptic, then
\begin{equation}
\spec_{\RB^2}(\BA) = \spec_{\RL^2}(\BA)
\end{equation}
as a set. In particular, if $\BA$ is bounded, its norm in $\RL^2$ and $\RB^2$ coincide.
\end{prop}

The proof of this proposition is exactly the same as the one in~\cite{Shubin1978} for the case $m = 1$. Indeed, it relies on some facts about function approximation proven in~\cite[Lemmata~4.1 and 4.2]{Shubin1978} which remain true as $m > 1$ since they apply coordinatewise. Boundedness or ellipticity then implies Proposition~\ref{prop:normspec}. When we refer to the norm of an operator, we will not distinguish whether that operator is acting in $\RL^2(\R^d;\C^m)$ or $\RB^2(\R^d;\C^m)$ since those norms are the same.

As mentioned in Remark~\ref{rem:Shubin}, there is a faithful, norm-preserving $*$-representation $\BA \mapsto \BA^\sharp$ of almost periodic operators given by $\BA^\sharp := \ba(\bx + \by,D_\by)$
acting in
\[
\FH_m := \RB^2\left(\R^d\right) \otimes \RL^2\left(\R^d\right) \otimes \C^m.
\]
Here $\bx$ is a variable of functions in $\RB^2(\R^d;\C^m)$ and $\by$ is a variable of functions in $\RL^2(\R^d;\C^m)$. The operator $\BA^\sharp$ is interpreted as a direct integral over $\bx$ of operators acting in $\RL^2(\R^d;\C^m)$. We denote by $e_J(x,y)$ the Schwartz kernel of the spectral projection~$E_J(\BA)$. Note that in view of Proposition~\ref{prop:normspec} and~\cite{CoburnMoyerSinger}, if $\BA \le \BB$ as operators, then $\BA^\sharp \le \BB^\sharp$ and $\norm \BA = \norm{\BA^\sharp}$.

Finally, the operator $\BA^\sharp$ is affiliated to the $\rmII_\infty$ factor $\FA$ generated by the two families of operators
\[
\set{\be_{\bxi} \otimes \be_{\bxi}\otimes M: \bxi \in \R^d,\ M\in \CM_m} \quad \text{and} \quad \set{I \otimes T_{\bxi}\otimes M: \bxi \in \R^d,\ M\in \CM_m},
\]
where $\be_{\bxi}$ is the operator of multiplication by $e^{i \bxi \cdot \bx}$, $T_{{\bxi}}$ is the operator of translation $T_{\bxi} f(\bx) = f(\bx - \bxi)$ and $\CM_m$ is the algebra of $m\times m$ matrices with complex entries. This means that the results of Section~\ref{sec:ids}--\ref{sec:mapo} on the density of states measure (DSM) also called the integrated density of states (IDS) apply to this algebra of operators and this representation.

In the classical setting, the IDS is defined for differential operators using the large box limit and for pseudo-differential operators as the trace of the Schwartz kernel
\begin{equation}
N(J;\BA) = M_\bx(\operatorname{tr} e_J(\bx,\bx)),
\end{equation}
where $M$ is the almost periodic mean. Note that this kernel is actually a smooth integral kernel whenever $J$ is a bounded interval, see~\cite{ParSht2016}.

Our terminology for the IDS is justified in~\cite[Remark~3.1]{Shubin1979}, where it is shown that the IDS as defined in Section~\ref{sec:ids} is the same as the one obtained from the classical definition for either differential or pseudo-differential operators.

\subsection{Concrete systems of operators}

From now on, we turn our attention to almost periodic pseudo-differential operators whose principal symbol is diagonal and nondegenerate.
\begin{defi}
A \emph{\pd} operator is an operator $\BA \in \BE\BS_m^\alpha$ for which there exists an unitary operator $\BU \in \BS_m^0$ so that $\BU^* \BA \BU = \BA_0 + \BB \in
\BS\BE\BS_m^\alpha$ has the following properties.
\begin{itemize}
\item The \emph{principal part} $\BA_0 \in \BD \BE \BS_m^\alpha$, with symbol
\begin{equation}\label{eq:defh0}
\ba_0(\bxi) = \diag\left(a_1 \abs \bxi^\alpha,\,\dotsc,\,a_m \abs
\bxi^\alpha\right),
\end{equation}
with $a_j \ne 0$ and without loss of generality $a_1 \ge \dotso \ge a_m$. We set $m^+ = \max\set{j : a_j > 0}$, where by convention $m_+ = 0$ if $a_1 < 0$.
\item The \emph{subprincipal part} $\BB \in \BS_m^\beta$ for $\beta < \alpha$ and has frequency set $\Theta$. We also suppose that $\BB$ is formally self-adjoint, i.e. that its symbol satisfies
\begin{equation}\label{eq:selfadjb}
\bbb_{\btheta}(\bxi) = \bbb_{-\btheta}(\bxi + \btheta)^*,
\end{equation}
for all $\bxi \in \RD$ and $\btheta \in \Theta$, where for a matrix $\ba$, $\ba^*$ is its conjugate transpose.
\end{itemize}
If $a_j \ne a_k$ for $j \ne k$, we say that $\BA$ is a \emph{\mpu} operator.
\end{defi}

\begin{rema}
Since we are interested only in spectral properties of elliptic operators, for the remainder of this paper we can always assume that the operators are already in $\BS\BE\BS_m^\alpha$.
\end{rema}

Without loss of generality, we assume that the frequency set $\Theta$ spans $\R^d$, contains $\mathbf 0$, and is symmetric about $\mathbf 0$. Recall from~\eqref{eq:ZTheta} that, using sum rather than product notations for the group of shifts in $\R^d$, that we also put
\begin{equation}\label{eq:thetasum}
\Theta^k = \Theta + \dotso + \Theta,
\end{equation}
where the sum is taken $k$ times, and
\begin{equation}\label{eq:ZTh}
Z(\Theta) = \bigcup_{k\,\in\,\N}
\Theta^k
\end{equation}
The set $Z(\Theta)$ is countable and non-discrete, unless $\Theta$ generates a lattice.


\subsection{Conditions on the perturbation and its frequency set}\label{sec:prelim}

In this section, we state the exact conditions under which we can obtain asymptotics for the integrated density of states for a system of operators acting in $\RB^2(\R^d;\C^m)$. We also show how we can reduce the problem to computing the IDS solely on some intervals contained in a large enough range of energies.


We are interested in the asymptotics for the positive energy and negative energy integrated densities of states for an \pd{} operator $\BA = \BA_0 + \BB$, defined as
\begin{align}
N^+(\lambda) &{:=} N^+(\lambda;\BA) := N([0,\lambda);\BA),\label{eq:npositive}
\\
\intertext{and}
N^-(\lambda) &{:=} N^-(\lambda;\BA) := N((-\lambda,0];\BA),\label{eq:nnegative}
\end{align}
as $\lambda \to \infty$. For this, we will need some conditions on the frequency set of the perturbation $\BB$. In Section~\ref{sec:asyexp}, we reduce the operator $\BA$ to a direct sum of operators of the type appearing in~\cite{MorParSht2014}. In that paper, the perturbations are required to satisfy some conditions, which we describe for completeness. Conditions~\ref{condI} and~\ref{condIV} correspond to~\cite[Conditions~A and~C]{MorParSht2014} and we do not use them explicitly. Condition~\ref{condII} addresses~\cite[equation 2.4]{MorParSht2014}, while Condition~\ref{condIII} addresses~\cite[Condition~B]{MorParSht2014}. We refer the reader to~\cite{MorParSht2014}, as well as~\cite{ParSht2012} for a more detailed discussion around these conditions and their genericity.

We first need the following generic condition on the set $Z(\Theta)$ defined in~\eqref{eq:ZTh}.

\begin{enonce}
{Condition}\label{condI}
Suppose that $\btheta_1,\,\dotsc,\,\btheta_d \in Z(\Theta)$. Then, $Z(\set{\btheta_1,\,\dotsc,\,\btheta_d})$ is discrete.
\end{enonce}

This condition is clearly satisfied for periodic $\BB$, but for quasi-periodic or almost periodic $\BB$ it is meaningful. The next two conditions describe how well $\BB$ is approximated by finite sums of homogeneous functions of $\bxi$, and by quasi-periodic operators.

\begin{enonce}
{Condition}\label{condII}
There exists a constant $C_0 > 1$ and a discrete subset $J \subset (-\infty,\beta]$ such that for all $\btheta \in \R^d$ and $\abs \bxi \ge C_0$,
\begin{equation}
(1 - \bone_{C_0}(\bxi)) \bbb_{\btheta}(\bxi) = \sum_{\iota\,\in\,J}
\abs{\bxi}^\iota
\bbb_{\btheta}^{(\iota)}\left(\frac{\bxi}{\abs \bxi} \right),
\end{equation}
where $\bbb_{\btheta}^{(\iota)} \in \BS_m^0$ is positively homogeneous of degree~$0$. We also suppose that for $\boeta \in \BBS^{d-1}$, $\bbb_{\btheta}^{(\iota)}(\boeta)$ has a series representation (written in multi-index notation)
\begin{equation}
\bbb_{\btheta}^{(\iota)}(\boeta) = \sum_{\bn\,\in\,\N_0^d}
\bbb_{\btheta}^{(\iota,\bn)} \boeta^\bn,
\end{equation}
which converges absolutely in a ball of radius greater than one of $\R^d$.
\end{enonce}

If $\BB$ is quasi-periodic and $J_0$ is finite, these are the only conditions that we need. Otherwise, we need to find a quasi-periodic approximation of $\BB$. In view of~\eqref{eq:smallnorm}, such an approximation will always exist, but we need a quantitative version of it.

\begin{enonce}
{Condition}\label{condIII}
For every $k \in \N$, there exists $C_k > C_0$ such that for each $\rho > C_k$, there exists a finite symmetric $\tilde \Theta \subset (\Theta \cap \CB(\rho^{1/k}))$ and a finite subset $J_k
\subset (-\infty,\beta]$ with
\begin{equation}\label{eq:cardtildej}
\# J_k \le \rho^{1/k}
\end{equation}
such that the symbol
\begin{equation}\label{eq:symbolrb2}
\br_{\btheta}^{(k)}(\bxi) :=
\begin{cases}
\bbb_{\btheta(\bxi)} & \text{if } \theta \not \in \tilde \Theta, \\
\bbb_{\btheta(\bxi)} - \sum_{\iota\,\in\,J_k} \abs{\bxi}^{\iota}
\bbb_{\btheta}^{(\iota)}\left(\frac{\bxi}{\abs \bxi} \right) & \text{if } \theta \in \tilde \Theta,
\end{cases}
\end{equation}
satisfies, for all $\ell \in \N$,
\begin{equation}\label{eq:smallseminorm}
\snorm{\br^{(k)}}{\beta}{\ell} \le c_{\ell,k} \rho^{-k},
\end{equation}
for some $c_{\ell,k} > 0$.
\end{enonce}

Finally, we need a Diophantine condition on the frequencies of $B$, for which we need some definitions. Fix $\tilde k \in \N$ (which will depend on the order of the remainder in the asymptotic expansion, but not on $k$ as in Condition~\ref{condIII}). We say that $\FV$ is a quasi-lattice subspace of dimension $q$ if there are linearly independent $\btheta_1,\,\dotsc,\,\btheta_q \in \tilde \Theta^{\tilde k}$ such that $\FV = \spann(\btheta_1,\,\dotsc,\,\btheta_q)$. We denote by $\CV$ the collection of all quasi-lattice subspaces.

We need a notion of non-trivial angle between two subspaces which are strongly distinct, i.e. so that neither is a subspace of the other. For this, we use the
\emph{largest principal angle} (which we abbreviate simply as ``the angle'') between subspaces, which is defined variationally recursively in the following way. Let $\FU, \FV \subset \R^d$ be subspaces with $\min(\dim(\FU),\dim(\FV)) = \ell$. The first angle $\phi_1(\FU,\FV) \in [0,\pi/2]$ between them is defined via
\begin{equation}
\cos(\phi_1(\FU,\FV)) := \max \set{\frac{\abs{\bx \cdot \by}}{\abs \bx \abs \by} : \bx \in \FU, \by \in \FV},
\end{equation}
and we fix unit vectors $\bx_1, \by_1$ realising this first angle. Then, for $1
\le q \le \ell$, the $q^{\rm th}$ angle is defined as
\begin{multline}
\cos(\phi_q(\FU,\FV)) :=\\
\max \set{\frac{\abs{\bx \cdot \by}}{\abs \bx \abs \by} : \bx \in \FU, \by \in \FV, \bx \cdot \bx_p = 0, \by \cdot \by_p = 0
\text{ for all } 1 \le p < q}.
\end{multline}
We then set the angle between $\FU$ and $\FV$ as $\phi(\FU,\FV) =
\phi_\ell(\FU,\FV)$. This angle is non-zero if and only if $\FU$ and $\FV$ are strongly distinct, and is $\pi/2$ as soon as there is a vector in one subspace orthogonal to the other.

Recalling that for any $k$ the choice of $\tilde \Theta$ depends on $\rho$, we put
\begin{equation}\label{eq:defnrR}
R(\rho) = \sup_{\btheta\,\in\,\tilde \Theta^{\tilde k}} \abs \btheta, \qquad r(\rho) = \inf_{\btheta\,\in\,(\tilde \Theta^{\tilde k})'} \abs \btheta,
\end{equation}
as well as
\begin{equation}
s := s(\rho) = s(\tilde \Theta^{\tilde k}) := \inf \sin(\phi(\FU,\FV)),
\end{equation}
where the infimum is over all strongly distinct pairs of subspaces in $\CV$. It is clear that
\begin{equation}\label{eq:boundonR}
R(\rho) = \bigo{\rho^{1/k}},
\end{equation}
where the implicit constant might depend on $k$ and $\tilde k$; however, we need the following condition for $r$ and $s$.

\begin{enonce}{Condition}\label{condIV}
For each fixed $k$ and $\tilde k$, the sets $\tilde \Theta$ can be chosen in such way that for sufficiently large $\rho$, depending on $k$ and $\tilde k$, the number of elements of $\tilde \Theta^{\tilde k}$ satisfies $\#\tilde \Theta^{\tilde k} \le \rho^{1/k}$ and we have that
\begin{equation}\label{eq:boundons}
s(\rho) \ge \rho^{-1/k}
\end{equation}
and
\begin{equation}\label{eq:boundonr}
r(\rho) \ge \rho^{-1/k}.
\end{equation}
\end{enonce}

\begin{rema}
Condition~\ref{condIV} is automatically satisfied for quasi-periodic and smooth periodic $\BB$. See~\cite{ParSht2012} for further discussion of this condition.
\end{rema}


\section{Asymptotic expansions for the IDS} \label{sec:asyexp}
We now suppose that the perturbation $\BB$ satisfies Conditions~\ref{condI}--\ref{condIV} and we set $\rho = \lambda^{1/\alpha}$, where $\alpha$ is the order of $\BA_0$. We prove the two following theorems, depending on whether all the $a_j$ in~\eqref{eq:defh0} are distinct or not. Recall that $m_+ = \max \set{j : a_j > 0}$, with $m_+ = 0$ if $a_j < 0$ for all $j$.
\begin{theo} \label{thm:aexpconcretecut}
Let $\BA$ be a \pd{} operator with subprincipal part $\BB \in \BS^{\beta}_m$, $\beta \le \alpha/2$ satisfying Conditions~\ref{condI}--\ref{condIV}. Suppose that there exists $\gamma \le 0$ such that whenever $a_j = a_k$ for some $1 \le j \ne k \le m$, then $[\BB]_{j,k} \in \BS^{\gamma}$ and put $\gamma^* = \max(2 \beta - \alpha, \gamma)$. Then, there exists a discrete set $L \subset (0,1 - \gamma^*)$ and constants $C_0^\pm$ and $C_{q,j}^\pm$, $0 \le q \le d-1$, $j \in L$ such that
\begin{equation}\label{eq:asymptoticsconcretecut}
\begin{aligned}
N^\pm\left(\BA; \rho^\alpha \right) = C_0^\pm \rho^d + \sum_{j\,\in\,L}\sum_{q =0}^{d-1} C_{j,q}^\pm \rho^{d-j} \log^q(\rho) + \bigo{\rho^{d - 1 + \gamma^*}},
\end{aligned}
\end{equation}
as $\rho \to \infty$. If $m^+ = m$ (resp. if $m^+ = 0$), then $C_0^- = C^-_{j,q} = 0$ (resp. $C_0^+ = C^+_{j,q}$ = 0) except for $(j,q) = (d,0)$.
\end{theo}

\begin{theo} \label{thm:aexpconcrete}
Let $\BA$ be a \mpu{} operator satisfying Conditions~\ref{condI}--\ref{condIV}. Then, for every $K \in \R$ there exists a discrete set $L \subset (0,d+K)$ and constants $C_0^\pm, C_{q,j}^\pm$, $0 \le q \le d-1$, $j \in L$, such that
\begin{equation}
\begin{aligned}
N^\pm\left(\BA; \rho^\alpha \right) = C_0^\pm \rho^d +
\sum_{j\,\in\,L}\sum_{q =0}^{d-1} C_{j,q}^\pm \rho^{d-j} \log^q(\rho) + \bigo{\rho^{-K}},
\end{aligned}\label{eq:asymptoticsconcrete}
\end{equation}
as $\rho \to \infty$. If $m^+ = m$ (resp. if $m^+ = 0$), then $C_0^- = C^-_{j,q} = 0$ (resp. $C_0^+ = C^+_{j,q}$ = 0) except for $(j,q) = (d,0)$.
\end{theo}

\begin{rema}
Note that the statement for $m^+ \in \set{0,m}$ follows from the operator being semi-bounded either above or below, respectively.

Note as well that if $J \subset \Z$, i.e. if the symbol of $\BA$ is a classical symbol, see~\cite[Chapter~7]{TaylorII}, then $L=
\set{0,\,\dotsc,\,K + d - 1}$.

The set $L$ of allowable exponents can be made explicit, depending on $J$ and $K$, see~\cite[Remark~2.7]{MorParSht2014}.
\end{rema}

The proof of Theorems~\ref{thm:aexpconcretecut} and~\ref{thm:aexpconcrete} are obtained after many reductions to simpler cases. Recall that they are the general versions of Theorems~\ref{thm:dirac2dids} and~\ref{thm:dirac3dids} in the introduction.

\subsection{IDS for uncoupled operators} \label{sec:uncoupled}

In this subsection, we prove that the conclusion of Theorem~\ref{thm:aexpconcrete} holds in the special case where $\BA \in \BU\BS_m^\alpha$, regardless of whether an operator is \pd{} or \mpu{}. This means that in addition of satisfying the conditions of Section~\ref{sec:prelim}, its symbol is given by
\begin{equation}\label{eq:modeldmapo}
\ba(\bx,\bxi) = \ba_0(\bxi) + \bbb(\bx,\bxi),
\end{equation}
where $\bbb(\bx,\bxi)$ is a diagonal matrix.

\begin{prop}\label{thm:aexpconcreteu}
Let $\BA \in \BU\BS^{\alpha}_m$ be an \pd{} operator satisfying Conditions~\ref{condI}--\ref{condIV}. Then, for every $K \in \R$ there exists a discrete set $L \subset (0,d+K)$ and constants $C_0^\pm, C_{q,j}^\pm$, $0 \le q \le d-1$, $j \in L$, such that
\begin{equation}\label{eq:asymptoticsconcreteunc}
\begin{aligned}
N^\pm\left(\BA; \rho^\alpha \right) = C_0^\pm \rho^d +
\sum_{j\in L}\sum_{q =0}^{d-1} C_{j,q}^\pm \rho^{d-j} \log^q(\rho) + \bigo{\rho^{-K}},
\end{aligned}
\end{equation}
as $\rho \to \infty$. If $m^+ = m$ (resp. if $m^+ = 0$), then $C^-_0 = C^-_{j,q} = 0$ (resp. $C_0^+ = C^+_{j,q} = 0$) except for $(j,q) = (d,0)$.
\end{prop}
\begin{proof}
Since $\BA \in \BU\BS_m^\alpha$, it can be split as a direct sum of operators $A_1 \oplus \dotso \oplus A_m$ acting in the mutually orthogonal subspaces $\RB^2(\R^d) \otimes v_j$. As such, we have that on any interval $J$,
\begin{equation}\label{eq:sumIDS}
N(J;\BA) = \sum_{j=1}^m N(J;A_j).
\end{equation}
This means that, for $j\le m^+$, $A_j$ is semi-bounded below and acts invariantly on $\RB^2(\R^d) \otimes v_j$ as the operator considered in~\cite{MorParSht2014}. For $j > m^+$, it is the operator $-\BA_j$ that acts in such a way. From~\cite[Theorem~2.5]{MorParSht2014}, this means that $N((-
\infty,\lambda);\BA_j)$ (resp. $N((\lambda,\infty); \BA_j)$) enjoys an asymptotic expansion of the form~\eqref{eq:asymptoticsconcrete} for $1 \le j \le m^+$ (resp. $m^+ < j \le m$). Observe that we have
\begin{multline}\label{eq:sumIDSsplit}
N^+\left(\rho^\alpha;\BA\right)=\\ \sum_{j=1}^{m^+} N((-\infty,\rho^\alpha);A_j) - \sum_{j=1}^{m^+} N((-\infty,0];A_j)
+ \sum_{j=m^++1}^m N((0,\rho^\alpha);A_j).
\end{multline}
The terms in the first sum have the required asymptotic expansion. The terms in the second sum do not depend on $\rho$, hence they might only change the constant term in~\eqref{eq:asymptoticsconcrete}. Finally, the operators in the third sum are semi-bounded above, hence for $\rho$ large enough the terms are constant and once again only affect the constant term. This proves the existence of the asymptotic expansion~\eqref{eq:asymptoticsconcrete} for $N^+$. The proof for $N^-$ is the same, interchanging the role of the semi-bounded below and above operators.
\end{proof}


\subsection{Reduction to a finite interval}

The strategy in this subsection is an adaptation of the one found in~\cite{MorParSht2014,ParSht2012}. It consists in showing that an asymptotic expansion holds in overlapping dyadic intervals $I_n$.

For $K > -d$, we choose $\rho_0$ sufficiently large, to be fixed later. For every $n \in \N$, we put $\rho_n := 2 \rho_{n-1} = 2^n \rho_0$. We also define the intervals $I_n := [\rho_{n-1},\rho_{n+1}]$. We prove the following theorem, which implies Theorems~\ref{thm:aexpconcretecut} and~\ref{thm:aexpconcrete} as a corollary.
\begin{theo} \label{thm:ainint}
Let $\BA$ be an operator satisfying the conditions of either Theorem~\ref{thm:aexpconcretecut} or~\ref{thm:aexpconcrete}. Then, for either $K = -d + 1 - \gamma^*$ in the former case or any $K \in \R$ in the latter, there exists $\rho_0$ large enough, a discrete set $L \subset (0,d+K)$ and constants $C_0^\pm, C^\pm_{j,q}$ for every $j \in L$ and $0
\le q \le d-1$ such that for every $n \in \N$ and every $0 < \mu < \nu$ with $\mu, \nu \in I_n$,
\begin{multline}\label{eq:ainint}
N((\mu^\alpha,\nu^\alpha);\BA) =\\
C_0^+(\nu^d - \mu^d) +
\sum_{j\in L}\sum_{q =0}^{d-1} C_{j,q}^+
\left(\nu^{d-j} \log^q(\nu) - \mu^{d-j} \log^q(\mu)\right) + \bigo{\rho_n^{-K}},
\end{multline}
where the implicit constants might depend on $K$, but not on $n$. Similarly,
\begin{multline}\label{eq:ainintneg}
N((-\nu^\alpha,-\mu^\alpha);\BA) =\\
 C_0^- (\nu^d - \mu^d) +
\sum_{j\in L}\sum_{q =0}^{d-1} C_{j,q}^-
\left(\nu^{d-j} \log^q(\nu) - \mu^{d-j} \log^q(\mu)\right) +
\bigo{\rho_n^{-K}}.
\end{multline}
\end{theo}

\begin{rema}
The reader familiar with previous works on the integrated density of states for almost periodic operators can notice that the roles of the dyadic decomposition in intervals $I_n$ is slightly different here. In previous work, this decomposition was necessary because the resonant zones were significantly different for different spectral intervals. This yielded coefficients $C^\pm$ depending possibly on $n$. It was however shown that the asymptotics had to match if the coefficients didn't grow too fast.

In our case, we need this decomposition in order to apply Theorem~\ref{thm:ctou} when the perturbation is unbounded. Indeed, it relies on Lemma~\ref{lem:spectralperturb} which can only be applied for some interval with control on how far away the endpoints can be. We will therefore obtain asymptotics when both endpoints belong to a specific dyadic interval, then glue the intervals together. We end up comparing the density of states with the one obtained in~\cite{MorParSht2014} for operators acting on scalar functions, i.e. the case $m = 1$. In such a case, the dependence on $n$ of the coefficients has already been removed.
\end{rema}

\begin{proof}[Proof of Theorems~\ref{thm:aexpconcretecut} and~\ref{thm:aexpconcrete} assuming Theorem~\ref{thm:ainint}] We prove the theorem for $N^+$, the proof for $N^-$ is the same. For $K \in \R$, suppose that $\rho_0$ is large enough for Theorem~\ref{thm:ainint} to hold. Suppose without loss of generality that for all $n$, $\rho_n$ is a point of continuity of $N^+$. For $\rho \in I_n$, we have that
\begin{equation}\label{eq:foldingsum}
\begin{aligned}
N^+(\rho^\alpha) &= N^+(\rho_0^\alpha) + \sum_{j=1}^{n-1} N((\rho_{j-1}^\alpha, \rho_j^\alpha);\BA) + N((\rho_{n-1}^\alpha,\rho^\alpha);\BA) \\
&= N^+(\rho_0^\alpha) +
\sum_{j\in L} \sum_{q =0}^{d-1} C_{j,q}^+
\left(\rho^{d-j} \log^q(\rho) - \rho_0^{d-j} \log^q(\rho_0)\right) +
\sum_{j=1}^n S_j,
\end{aligned}
\end{equation}
where $S_j = \bigo{\rho_j^{-K}}$. This implies that
\begin{equation}\label{eq:glued}
\sum_{j=1}^n S_j \ll \rho_0^{-K} \sum_{j = 1}^n 2^{-Kj} \ll \rho_n^{-K} \ll \rho^{-K}
\end{equation}
since $\rho \in I_n$. One can see that the term depending on $\rho_0$ is $\bigo{1}$, so that it can be included either in the error term $(\rho^{-K})$ when $K \le 0$ or in the constant term in~\eqref{eq:asymptoticsconcrete} and~\eqref{eq:asymptoticsconcretecut} otherwise.
\end{proof}

\subsection{Reduction to a quasiperiodic operator}

We now show in the following lemma that it is sufficient to prove Theorem~\ref{thm:ainint} for quasiperiodic operators.
\begin{lemm}\label{lem:qpenough}
Let $\BA \in \BS\BE\BS_m^\alpha$ be an \pd{} operator with subprincipal part $\BB
\in \BS^\beta_m$ satisfying Condition~\ref{condIII} and $k \ge 2$. Then, there exists $\rho_0 > 0$ and $0 < c_0 < 1$ so that for every $n \in \N$ there exists a quasi-periodic \pd{} operator $\BA' \in \BS\BE \BS_m^\alpha$ with frequency set $\tilde \Theta \subset
\CB(\rho_n^{1/k})$ such that
\begin{itemize}
\setlength
\itemsep{1em}
\item $\BA - \BA' \in \BS_m^\beta$;
\item $\supp(\ba'{}^{\CO\CD}) \subset \set{\abs{\bxi} > c_0 \rho_n}$;
\item there is $\eps \ll \rho_n^{\alpha-k}$ such that for all $J \subset I_n^\alpha$,
\begin{equation}\label{eq:qpapprox}
N((\pm J)_{-\eps};\BA') \le N(\pm J;\BA) \le N((\pm J)_\eps;\BA').
\end{equation}
\end{itemize}
\end{lemm}

\begin{proof}
For $k \in \N$ let $\tilde \Theta \subset \Theta \cap \CB(\rho_n^{1/k})$ be the frequency set given by Condition~\ref{condIII} with $\rho = \rho_n$, and $\BR \in \BS^\beta_m$ be the operator with symbol given in~\eqref{eq:symbolrb2}, which by~\eqref{eq:smallseminorm} satisfies $\snorm{\BR}{\beta}{0}
\ll \rho_n^{-k}$. Setting $\BA'' = \BA - \BR$ we have that $\BA - \BA'' \in
\BS_m^{\beta}$, and that $\tilde \Theta$ is a frequency set for $\BA''$ and that as long as $\rho_0$ is large enough,
\begin{equation}\label{eq:snormcomp}
\snorm{\BA''}{\gamma}{0} \le 2 \snorm{\BA}{\gamma}{0}
\end{equation}
for all $\beta \le \gamma \le \alpha$.

Writing any interval $J \subset \pm I_n^\alpha$ in the form $[M - r,M+r]$, it is easy to see that $\abs M + r \le (2 \rho_n)^\alpha$. Put $\beta_0 = \max\set{\beta,0}$. By Lemma~\ref{lem:monotonicity}, estimate~\eqref{eq:qpapprox} holds with $\BA''$ instead of $\BA'$ and
\begin{equation}
\eps_1 = \frac{\snorm{\BR}{\beta_0}{0}}{2 + \snorm{\BR}{\beta_0}{0}}\left(\abs M + r + C\left(1 +
\snorm{\BR}{\beta_0}{0}\right)^\frac{\alpha}{\alpha - \beta_0}\right) \ll \rho_n^{\alpha - k}
\end{equation}
instead of $\eps$. Let us now define
\begin{equation}
\BA' = \BA'' - \BA''{}^{\CO\CD}P_{c_0\rho_n} = (\BA'')^\CD +
\BA''{}^{\CO\CD}(1 - P_{c_0 \rho_n}),
\end{equation}
where $0 < c_0 < 1$ is to be determined later. By~\eqref{eq:snormcomp} and~\eqref{eq:partition}
\begin{equation}
\snorm{\BA'}{\gamma}{0} \le 4 \snorm{\BA}{\gamma}{0}
\end{equation}
for all $\beta \le \gamma \le \alpha$. We apply Lemma~\ref{lem:spectralperturb} with
\begin{equation}
H_0 = (\BA'')^{\CD}, \quad B = \BA''{}^{\CO\CD}(1 - P_{c_0\rho_n}), \quad A =
\BA''{}^{\CO\CD}P_{c_0\rho_n}, \quad H = \BA'.
\end{equation}
By Proposition~\ref{prop:normorder},
\begin{equation}
\norm{\BA''{}^{\CO\CD}P_{c_0\rho_n}} \le (c_0 \rho_n)^{\beta_0}
\snorm{\BA''{}^{\CO\CD}}{\beta_0}{0}.
\end{equation}
Set $X = \lfloor (2 - \alpha + k + \beta_0) \log_3 \rho_n \rfloor$, and let
\begin{equation}\label{eq:increasingproj}
Z_l := c_0 \rho_n + l \rho_n^{2/3}, \quad 0 \le l \le X - 1,
\end{equation}
so that if $\rho_0$ is large enough, $Z_{X - 1} \le 2 c_0 \rho_n$. For $0 \le l \le X$ introduce the family of projections
\begin{equation}
P_l :=
\begin{cases}
P_{Z_0} & \text{for } l = 0, \\
P_{Z_l} - P_{Z_{l -1}} & \text{for } 0 < l < X, \\
1 - P_{Z_{X - 1}} & \text{for } l = X.
\end{cases}
\end{equation}
We now verify that the conditions of Lemma~\ref{lem:spectralperturb} are satisfied. It is clear that $\BB P_{c_0\rho_n} P_{Z_0} = \BB P_{c_0\rho_n}$, and relations~\eqref{eq:projections} follow from~\eqref{eq:boundonR} and~\eqref{eq:increasingproj} as long as $k
\ge 2$ and $\rho_0$ is large enough. By Proposition~\ref{prop:normorder}, for $0 \le l < X$,
\begin{equation}\label{eq:dproj}
\norm{P_l \BA' P_l} \le Z_{X-1}^\alpha \snorm{\BA'}{\alpha}{0}
\le 4(2 c_0
\rho_n)^\alpha \snorm{\BA}{\alpha}{0}.
\end{equation}
We also have that
\begin{equation}\label{eq:odproj}
\norm{P_l \BA''{}^{\CO\CD} P_{l-1}} +
\norm{P_l \BA''{}^{\CO\CD} P_{l+1}} \le 2 Z_{X-1}^{\beta}
\snorm{\BA''}{\beta_0}{0} \le 4 (2c_0\rho_n)^{\beta_0}
\snorm{\BA}{\beta_0}{0}
\end{equation}
For $0 \le l < X$, set
\begin{equation}
D_l = \dist(J,\spec(P_l \BA' P_l)^\sharp).
\end{equation}
By~\eqref{eq:dproj} and Lemma~\ref{lem:sharpS0} for $l \le X - 1$
\begin{equation}
\spec((P_l \BA' P_l)^\sharp) \subset \left[- 4(2 c_0
\rho_n)^\alpha\snorm{\BA}{\alpha}{0}, 4(2 c_0 \rho_n)^\alpha\snorm{\BA}{\alpha}{0}
\right],
\end{equation}
so that setting $ c_0^{-\alpha} = 2^{\alpha+ 3} \snorm{\BA}{\alpha}{0}$ gives
\begin{equation}
D_l \ge \frac{\rho_n^\alpha}{2},
\end{equation}
in particular~\eqref{eq:distspectrum} holds. Combining with~\eqref{eq:odproj} we have that
\begin{equation}
\max_{0\,\le\,l\,<\,X} \set{\frac{
\norm{P_l \BA''{}^{\CO\CD} P_{l-1}}+
\norm{P_l \BA''{}^{\CO\CD} P_{l+1}}}{D_l}} \le 2 (2c_0)^\beta \rho_n^{\beta -
\alpha} \snorm{\BA}{\beta_0}{0},
\end{equation}
so that for $\rho_0$ large enough, \eqref{eq:projectionnorm} is satisfied.

Since the conditions of Lemma~\ref{lem:spectralperturb} are satisfied, for
\begin{equation}
\eps_2 = 3^{2 - X} \norm{\BA''{}^{\CO\CD}P_{c_0\rho_n}} \le 2\rho_n^{\alpha - k} \snorm{\BA}{\beta_0}{0},
\end{equation}
we have that
\begin{equation}
N\left(I_{- \eps_1 - \eps_2}; \BA' \right) \le N(I_{-\eps_1};\BA'') \le N(I;\BA);
\end{equation}
and
\begin{equation}
N(I_{+ \eps_1 + \eps_2};\BA') \ge N(I_{+\eps_1};\BA'') \ge N(I;\BA).
\end{equation}
Our claim therefore holds with $\eps = \eps_1 + \eps_2$.
\end{proof}



\subsection{Reduction to uncoupled operators} \label{sec:redtouncoupled}
In this subsection, we show that in order to prove Theorem~\ref{thm:ainint} for quasi-periodic operators, it is sufficient to do so for uncoupled operators. At the end of the section, we finally prove Theorem~\ref{thm:ainint} after all those reductions, which completes the proof of Theorems~\ref{thm:aexpconcretecut} and~\ref{thm:aexpconcrete}.

\begin{theo}\label{thm:ctoucut}
Let $\BA$ be an operator satisfying the conditions of Theorem~\ref{thm:aexpconcretecut}. Then, for every $n \in \N$ there is an operator $A' \in \BU\BS^\alpha_m$ and $\eps \ll
\rho_n^{\gamma^*}$ such that for all $\mu,\nu \in I_n$ and $I = (\mu^\alpha,
\nu^\alpha)$,
\begin{equation}\label{eq:errorestinterval}
N(\pm I_{- \eps}; \BA') \le N(\pm I;\BA) \le N(\pm I_\eps;\BA').
\end{equation}
In particular,
\begin{equation}\label{eq:errorest}
N(\pm I;\BA) = N(\pm I;\BA') + \bigo{\rho_n^{d - 1 + \gamma^*}}.
\end{equation}
\end{theo}

\begin{proof}
We only prove this theorem for $I$, the proof for $-I$ follows from the same argument, making the relevant sign changes. By Proposition~\ref{thm:aexpconcreteu}, we see that~\eqref{eq:errorest} follows from~\eqref{eq:errorestinterval}. We therefore only prove the latter.

We keep a quantitative track of the estimates found in Section~\ref{sec:reducscalar}. By Lemma~\ref{lem:qpenough}, we can without loss of generality assume for some fixed $C_0 > 0$ that $\BA$ is a quasi-periodic operator whose frequency set $\Theta$ lies in the ball $B(\rho_n^{1/k})$ for some $k \in \N$ and such that the support of $\ba^{\CO\CD}$ lies in $\set{\abs \xi > C_0 \rho_n}$. In particular, we can assume that there is $s > 0$ such that for all $\xi \in \supp(\ba^{\CO\CD})$, all $\theta \in \Theta$ and all $j,k$ such that $a_j \ne a_k$,
\begin{equation}
\Big|a_j \abs{\btheta + \bxi}^\alpha - a_k \abs{\bxi}^\alpha\Big| > s \abs
\bxi^\alpha.
\end{equation}
By Theorem~\ref{thm:systemonestep}, since $\BA$ is quasi-periodic there are symmetric operators $\BY \in
\BU\BS^{\beta}_m$, $\BR \in \BS_m^{\gamma^*}$ and $\BPsi \in \BS_m^{\beta - \alpha}$ such that $\BA$ is unitarily equivalent through conjugation with $\exp(i\BPsi)$ to
\begin{equation}
\BA' + \BR = \BA^{\CD} + \BY + \BR.
\end{equation}
Here, the symbol of $\BPsi$ is given by
\begin{equation}
[\boldsymbol \psi_{\btheta}(\bxi)]_{j,k} =
\frac{i[{\bbb^\CC}_{\btheta}(\bxi)]_{j,k} \chi_{j,k}}{a_j \abs{\bxi +
\btheta}^\alpha - a_k \abs \bxi^\alpha},
\end{equation}
where $\chi_{j,k} = 1$ if $a_j \ne a_k$ and $0$ otherwise.
Using the fact that
\[
\ad(\BA;\BPsi) = \ad\left(\BA^{\CO\CD};\BPsi\right) +
\ad\left(\BA^\CD;\BPsi\right) = \ad\left(\BA^{\CO\CD};\BPsi\right) - \BA^{\CNR},
\]
the operator $\BR$ is obtained from equations~\eqref{eq:decompo}, \eqref{eq:decompodetails} with $\widetilde \BR = 0$, and~\eqref{eq:defRconsecgauge} by
\begin{equation}\label{eq:defRonestepconcrete}
\BR = \BB^{\CR,\CC} + \ad\left(\BA^{\CO\CD};\BPsi\right) + \sum_{k=2}^\infty
\frac{1}{k!}\ad^k\left(\BA^{\CO\CD};\BPsi\right) - \sum_{k=2}^\infty \frac{1}{k!}
\ad^{k-1}\left(\BA^{\CNR};\BPsi\right).
\end{equation}
By Lemma~\ref{lem:monotonicity},
\begin{equation}
N(I_{-\eps};\BA') \le N(I;\BA' + \BR) \le N(I_\eps;\BA')
\end{equation}
for $\eps = \norm{\BR}$. Since $\BPsi$ has order $\beta - \alpha$ and is supported on $\set{\abs \xi > c \rho}$, by Corollary~\ref{cor:domain} and Lemma~\ref{lem:product} we have as in Proposition~\ref{prop:normorder} that
\begin{equation}
\norm{\ad\left(\BA^{\CO\CD};\BPsi\right)} \ll \rho_n^{2 \beta - \alpha} \snorm{\BA^{\CO\CD}}{\beta}{0}
\snorm{\BPsi}{-\beta}{\abs{\beta}},
\end{equation}
so this gives the contribution from the second term in~\eqref{eq:defRonestepconcrete}. The third and fourth terms uses the same estimate and the fact that this sum is absolutely convergent. Finally, for the first term we supposed that $\BB^{\CR,\CC} \in \BS_m^{\gamma}$, and it is also supported on $\set{\abs \xi > c \rho}$ so that by Proposition~\ref{prop:normorder},
\begin{equation}
\norm{\BB^{\CR,\CC}} \ll \rho_n^{\gamma}.
\end{equation}
Together, this completes the proof of Theorem~\ref{thm:ctoucut}.
\end{proof}

When $a_j \ne a_k$ whenever $j \ne k$, we get the following stronger statement.
\begin{theo} \label{thm:ctou}
Let $\BA$ be an operator satisfying the hypotheses of Theorem~\ref{thm:aexpconcrete}. There is a decreasing sequence $\set{\gamma_K}_{K\,\in\,\N}$, $\gamma_K \to - \infty$ such that for all $K, n \in \N$ there is an operator $\BA_K \in \BU\BS_m^\alpha$ and some $\eps \ll \rho_n^{-\alpha - d - K}$ such that for all $\mu,\nu \in I_n$, and $I = (\mu^\alpha,\nu^\alpha)$,
\begin{equation}\label{eq:desiredeps}
N(I_{-\eps};\BA_K) \le N(I;\BA) \le N(I_\eps;\BA_K)
\end{equation}
and such that if $K_1 < K_2$, then
\begin{equation}
\BA_{K_1} \equiv \BA_{K_2} \mod \BS^{\gamma_{K_1}}_m.
\end{equation}
In particular,
\begin{equation}\label{eq:approxbydiagonal}
N(\pm I;\BA) = N(\pm I; \BA_K) + \bigo{\rho_n^{-K}}.
\end{equation}
\end{theo}

\begin{proof}
This statement is proven in the same way as the previous one, replacing the use of Theorem~\ref{thm:systemonestep} with the parallel gauge transform Theorem~\ref{thm:gtsystem}, with a number of steps depending on $K$. This is possible because the condition on the terms coupling $a_j=a_k$ for $j
\ne k$ is vacuously verified, so that it is assuredly preserved after each step of gauge transform. This yields a remainder $\BR \in \BS_m^{-N}$ for any $N$, allowing for the arbitrary precision in the approximation for the density of states.
\end{proof}

\begin{rema}
Note that after this reduction, Conditions~\ref{condII} and~\ref{condIII}, corresponding to~\cite[Equation 2.4 and Condition B]{MorParSht2014}, do not hold anymore. However, the reason why these conditions are needed is to have a specific form for the functions $\bbb_{\btheta_1}(\bxi + \btheta_2)$, where $\btheta_1,\btheta_2 \in \tilde
\Theta$, see~\cite[Equation 10.5]{MorParSht2014}. This expansion still holds if these conditions are imposed on the symbol prior to reduction to uncoupled operators.
\end{rema}



\begin{proof}[Proof of Theorem~\ref{thm:ainint}]
By Theorems~\ref{thm:ctoucut} and~\ref{thm:ctou}, there is $\rho_0$ large enough so that for any $n$ there is an operator $\BA_K \in \BU\BS_m^\alpha$ such that for $\mu,\nu \in I_n$,
\begin{equation}
N((\mu^\alpha,\nu^\alpha);\BA) = N((\mu^\alpha,\nu^\alpha);\BA_K) +
\bigo{\rho_n^{-K}},
\end{equation}
where $K = 1- d - \gamma^*$ if $\BA$ satisfies the hypotheses of Theorem~\ref{thm:aexpconcretecut} and $K \in \R$ if $\BA$ satisfies the hypotheses of Theorem~\ref{thm:aexpconcrete}. Equation~\eqref{eq:ainint} (with coefficients $C^\pm_{j,q}$ depending on $n$) then follows by Proposition~\ref{thm:aexpconcreteu} for $\BA_K$ and the fact that
\begin{equation}\label{eq:munu}
N((\mu^\alpha,\nu^\alpha);\BA_K) = N^+(\nu^\alpha;\BA_K) - N^+(\mu^\alpha;\BA_K).
\end{equation}

In order to remove the dependence on $n$ of the coefficients, it is sufficient for every $K \in \N$ to prove that they must agree for all $n$ large enough. Since the coefficients obtained in Proposition~\ref{thm:aexpconcrete} do not depend on $n$, this means that as soon as $\mu, \nu \in I_{n} \cap I_{n+1}$, \eqref{eq:munu} gives the same coefficients for the asymptotic expansion up to terms of order $\rho_n^{-K}$, which means that the coefficients need to agree for all $n$ large enough.
\end{proof}



\section{The structure of periodic operators}\label{sec:periodic}

We now turn our attention to periodic operators. In this section, we describe the structure of operators that are periodic with respect to some lattice $\Lambda$, and we give a quantitative approach to the study of the Bethe--Sommerfeld property. We realise the usual Bloch--Floquet decomposition through Besicovitch spaces.


\subsection{Description of periodic operators}
For periodic operators we assume that $G$ is not $\R^d$ but rather the dual lattice $\Theta := \Lambda^\dagger \subset \R^d$. We note that in this case $Z(\Theta) = \Theta$.

Invariance of $\BA$ under the action of $\Lambda$ means that for all $\bk \in \R^d$, the subspace
\begin{equation}\label{eq:invsubspace}
\ell^2_\bk(\Theta;\C^m) := \overline{\spann \set{\be_{\bxi,j} : 1 \le j \le m, \bxi \in
\Theta + \bk} } \subset \RB^2\left(\R^d;\C^m\right)
\end{equation}
is an invariant subspace for $\BA$, and we denote by $\BA(\bk)$ the restriction of $\BA$ to this subspace. It is clear from the definition that we can restrict ourselves to $\bk \in \CO^\dagger = \R^d/\Lambda^\dagger = \R^d/\Theta$, and we call $\bk$ a
\emph{quasimomentum}. For any $\bxi \in \R^d$, its fractional part $\set \bxi \in
\CO^\dagger$ is the image of $\bxi$ under the quotient map. The spectrum of $\BA$ can be obtained as
\begin{equation}
\spec(\BA) = \bigcup_{\bk\,\in\,\CO^\dagger} \spec(\BA(\bk)),
\end{equation}
see~\cite[Theorem~4.5.1]{ Kuchment}. For every $\bk \in \CO^\dagger$ the spectrum of $\BA(\bk)$ in $\ell^2_\bk(\Theta;\C^m)$ is discrete.

The usual approach to studying the $\RL^2(\R^d;\C^m)$ theory of periodic operators is through the Bloch--Floquet decomposition, see e.g.~\cite{Kuchment}, where we represent $\BA$ as a direct integral over $\CO^\dagger$ of the fibre operators $\BA(\bk)$. This would require us to introduce a considerable amount of machinery. However, the Bethe-Sommerfeld property is strictly about the spectrum as a set, and to every elliptic periodic operator acting in $\RL^2(\R^d;\C^m)$ there corresponds an elliptic operator acting in $\RB^2(\R^d;\C^m)$ with the same spectrum. We therefore consider periodic operators as operators on Besicovitch space where we decompose them according to the invariant subspaces~\eqref{eq:invsubspace}. This makes our statements and proofs more direct.
\begin{rema}\label{rem:Iinfinitybs}
The subspaces $\ell^2_\bk(\Theta;\C^m)$ can be realised as
\[
\RH^0(\Theta + \bk;\C^m) = \overline{\spann\left(\left\{\be_{\btheta + \bk,j},
\btheta \in \Theta, 1 \le j \le m\right\}\right)}.
\]
The group $G$ is $\Theta$ acting on
$\Theta + \bk$ by translation. In this case, consider the $\rmI_\infty$ factor $\FA$ generated by $
\set{\be_{\btheta} \otimes M: \btheta \in \Theta, M \in \CM_m}. $

It is clear that the restriction of the subalgebra of periodic operators to $\ell^2_\bk(\Theta;\C^m)$ is affiliated to $\FA$, and that it respects the conditions described at the beginning of Section~\ref{sec:ids}. The associated trace of the spectral projection over an interval $J$ is simply $N(J;\BA(\bk)) := \#\set{j : \lambda_j(\BA(\bk)) \in J}$, the number of eigenvalues of $\BA(\bk)$ in that interval.

We also make the observation that for a bounded, periodic self-adjoint operator~$\BPsi$, the restriction to $\ell^2_\bk(\Theta;\C^m)$ of the unitary operator $\exp(i \BPsi)$ is still unitary, since $\ell^2_\bk(\Theta;\C^m)$ is an invariant subspace. This means that we can simultaneously use the gauge transform on each of the fibre operators and that the estimates from Section~\ref{sec:ids} hold uniformly for the counting function of the fibre operators.
\end{rema}

Let us now describe the structure of the spectrum of $\BA$ in terms of the spectra of the fibre operators $\BA(\bk)$. Since $\BA$ is self-adjoint, it has Fredholm index $0$, this implies that the Bloch variety
\begin{equation}
\set{(\bk,\lambda) \subset \CO^\dagger \times \R : \lambda \in \spec(\BA(\bk))}
\end{equation}
is a principal analytic set~\cite[Corollary~3.1.6 and Section~3.4.C]{Kuchment}. As such, if $\BA$ is semi-bounded below, we can naturally label the eigenvalues of $\BA(\bk)$ in non-decreasing order, counting multiplicity. Then, the functions $\lambda_j(\bk):= \lambda_j(\BA(\bk))$ are piecewise analytic functions of $\bk$. If $\BA$ is not semi-bounded, we can label the eigenvalues in non-decreasing order by $j \in \Z$ and it is possible to choose the labelling so that the functions $\lambda_j(\bk)$ are piecewise analytic. This requirement determines the labelling uniquely up to a uniform shift of the indices. Note that continuity in $\bk$ of the functions $\lambda_j$ and discreteness of the spectrum imply that labelling the eigenvalues at one quasimomentum $\bk$ induces a labelling everywhere in $\CO^\dagger$. The interval
\begin{equation}
\iota_j:= \iota_j(\BA) := \bigcup_{\bk\,\in \,\CO^\dagger} \lambda_j(\BA(\bk))
\end{equation}
is called the $j^{\rm th}$ spectral band of $\BA$.

\subsection{The overlap function}
In order to prove that an operator has the Bethe--Sommerfeld property we study the band overlap, characterized by the \emph{overlap function} $\zeta(\lambda;\BA)$, $\lambda \in \R$, introduced by M. Skriganov~\cite{Skriganov}. The overlap function is defined as the maximal number $t$ such that the symmetric interval $[\lambda-t,\lambda+t]$ is entirely contained in one band, i.e.
\begin{equation}\label{eq:defnoverlap}
\zeta(\lambda;\BA) :=
\begin{cases}
\max_j \max\left\{t \ge 0 : [\lambda-t,\lambda+t] \subset \iota_j\right\} & \text{if } \lambda \in
\spec(\BA)\\
0 & \text{if }\lambda \not \in \spec(\BA).
\end{cases}
\end{equation}
It is not hard to see that $\zeta$ is a continuous function of $\lambda$. In order to use our machinery we will relate the overlap function to the eigenvalue counting functions of the operators $\BA(\bk)$. This type of idea has been used in the past but crucially relied on the fact that $\BA$ was semi-bounded below. In the following proposition we find an equivalent formulation that is robust under perturbations yet works for operators that are not semi-bounded. Recall that for an interval $I = [s,t] \subset \R$ and $\eps
\in \R$, we define
\begin{equation}
I_\eps :=
\begin{cases}
\varnothing & \text{for } \eps < \frac{s-t}{2}, \\
[s - \eps,t + \eps] & \text{otherwise.}
\end{cases}
\end{equation}

\begin{lemm} \label{lem:overlapidscont}
Suppose that $\BA_1$, $\BA_2$ are self-adjoint periodic operators. Suppose that for all $p \in \set{1,2}$, $\bk \in \CO$, $\BA_p(\bk)$ has discrete spectrum. For $\lambda \in \R$ and $t > 0$, let
\begin{equation}\label{eq:smalldist}
\delta := \min_{\bk\,\in\,\CO^\dagger} \max\left\{\dist(\mu;\spec(\BA_1(\bk))) :
\mu \in [\lambda - t, \lambda + t]\right\}.
\end{equation}

Suppose that there is $0 \le \eps \le \delta/4$ such that for all $\bk \in
\CO^\dagger$ and any interval $I \subset [\lambda - t, \lambda + t]$
\begin{equation}\label{eq:condids}
N(I; \BA_2(\bk)) \le N(I_\eps;\BA_1(\bk))
\quad \text{and} \quad N(I;\BA_1(\bk)) \le N(I_\eps;\BA_2(\bk)).
\end{equation}

Then, for $p \in \set{1,2}$ there exist sets of consecutive integers $J_p \subset \Z$ and surjective maps
\[
\begin{aligned}
\lambda_\bullet(\BA_p(\bk)) : J_p &\to \spec(\BA_p(\bk)) \\
j &\mapsto \lambda_j(\BA_p(\bk))
\end{aligned}
\]
such that for all $j \in J_p$, $\lambda_j(\BA_p(\bk))$ are continuous in $\bk$ and such that
\begin{equation}\label{eq:evclose}
\left|\lambda_j(\BA_1(\bk)) - \lambda_j(\BA_2(\bk))\right| \le \eps
\end{equation}
for all $\bk$ and $j$ such that $\lambda_j(\BA_p(\bk)) \in [\lambda-t,\lambda+t]$.
\end{lemm}

\begin{rema}
We do not ask in the previous lemma that both operators share the properties of being either bounded, semi-bounded above or below, or unbounded in both directions.
\end{rema}

\begin{proof}
For any $\bk \in \CO^\dagger$, $0 < \eta \le \delta$, we say that $\mu \in [\lambda - t,
\lambda + t]$ is $\eta$-distant (from the spectrum of $\BA_1(\bk)$) at $\bk$ if
\begin{equation}\label{eq:choicemu}
\dist(\mu,\spec(\BA_1(\bk))) \ge \eta.
\end{equation}
By~\eqref{eq:smalldist}, for every $\bk \in \CO^\dagger$ there exists $\mu \in [\lambda - t, \lambda + t]$ which is $\delta$-distant at $\bk$. By the second inequality in~\eqref{eq:condids}, if $\mu$ is $\eta$-distant at $\bk$ for some $\eta > \delta/4$, then for $p \in \set{1,2}$
\begin{equation}\label{eq:muintempty}
(\mu - \eps,\mu + \eps) \cap \spec(\BA_p(\bk)) = \varnothing.
\end{equation}
Choose $\bk_0 \in \CO^\dagger$ and $\mu_0$ a point $\delta$-distant at $\bk_0$. Maps $j \mapsto \lambda_j(\BA_p(\bk))$ can be uniquely defined from the properties that they are nondecreasing, mapping to continuous functions in $\bk$, and that $\lambda_0(\BA_p(\bk_0))$ is the smallest eigenvalue larger than $\mu_0$. Note that the sets $J_1$ and $J_2$ are both defined uniquely from these properties, in particular if $\BA_p$ is unbounded both above and below then $J_p = \Z$.

We now prove that, for all $\bk \in \CO^\dagger$, if $\mu$ is $\delta/2$-distant at $\bk$, then for all $j \in J_1 \cap J_2$, then
\begin{equation}\label{eq:sameside}
\big(\lambda_j(\BA_1(\bk)) - \mu\big) \big(\lambda_j(\BA_2(\bk)) - \mu\big) > 0,
\end{equation}
in other words, for $p \in \set{1,2}$, $\lambda_j(\BA_p(\bk))$ are both on the same side of $\mu$. The functions $\lambda_j(\BA_p(\bk))$ were constructed specifically so that~\eqref{eq:sameside} holds at $\bk_0$ and $\mu_0$, our goal is to show that this property propagates to other $\mu$ and $\bk$.

We first prove that if~\eqref{eq:sameside} holds for some $\mu$ $\delta/2$-distant at $\bk$, then it holds for all other $\nu$ $\delta/2$-distant at $\bk$. This is a direct consequence of~\eqref{eq:condids} and~\eqref{eq:muintempty}, which imply that
\begin{equation}
N\big([\mu,\nu];\BA_1(\bk)\big) = N\big([\mu,\nu];\BA_2(\bk)\big).
\end{equation}

By continuity, for every $\bk \in \CO^\dagger$ there is $s_\bk > 0$ so that whenever $\mu$ is $\delta$-distant at $\bk$, $\mu$ is also $\delta/2$-distant at every $\bk' \in \CB(\bk,s_\bk)$. This also implies that if~\eqref{eq:sameside} holds at $\bk$ for one of those $\mu$, it also holds for that $\mu$ at every $\bk' \in \CB(\bk,s_\bk)$, and therefore at every $\nu$ $\delta/2$-distant at $\bk'$.

By compactness of $\CO^\dagger$, there are $\bk_1,\,\dotsc,\,\bk_\ell$ such that $\CO^\dagger$ is covered by the balls $U_j = \CB(\bk_j,s_{\bk_j}/2)$, with $0 \le j \le \ell$. If $U_{j} \cap U_{j'} \ne
\emptyset$, we have that $\bk_{j'} \in \CB(\bk_j,s_{\bk_j})$, so that if~\eqref{eq:sameside} holds for some $\mu$ $\delta$-distant at $\bk_j$ then it also holds for all $\nu$ $\delta/2$-distant at $\bk_{j'}$, and therefore also at any $\bk' \in U_{j'}$. By connectedness of $\CO^\dagger$, this means that~\eqref{eq:sameside} only needs to be verified for some $\bk_j$, $0 \le j \le \ell$ and one $\mu$ which is $\delta/2$-distant at $\bk_j$. Choosing $\bk_0$ and $\mu_0$, this means that~\eqref{eq:sameside} holds everywhere.



Suppose now that for some $p \in \set{1,2}$ there is some $j \in J_p$ and $\bk \in \CO^\dagger$ such that $\lambda_j(\BA_p(\bk)) \in [\lambda-t,\lambda+t]$ and
\begin{equation}
\lambda_j(\BA_1(\bk)) - \lambda_j(\BA_2(\bk)) > \eps.
\end{equation}
Let $\mu$ be $\delta$-distant at $\bk$. Without loss of generality assume that $[\mu,\infty) \cap \spec(\BA_p) \ne \varnothing$ and let
\begin{equation}
j' = \min \set{\ell : \lambda_\ell(\BA_p)(\bk) > \mu}.
\end{equation}
Supposing that $j \ge j'$, and using~\eqref{eq:muintempty} we obtain
\begin{equation}
N\left(\left[\mu,\lambda_j(\BA_2(\bk))\right];\BA_2(\bk)\right) \ge j + 1 - j' > N\left(\left[\mu,\lambda_j(\BA_2(\bk))\right]_{\eps};\BA_1(\bk)\right),
\end{equation}
which contradicts the second inequality in~\eqref{eq:condids}. Similarly, supposing that $\lambda_j(\BA_2(\bk))\linebreak - \lambda_j(\BA_1(\bk)) > \eps$ contradicts the first inequality in~\eqref{eq:condids}. The case $j < j'$ is treated analogously. We can therefore deduce that~\eqref{eq:evclose} holds at every $\bk \in \CO^\dagger$.
\end{proof}


The previous Lemma~\ref{lem:overlapidscont} admits the following corollary in the situation where the difference $\BA_1 - \BA_2$ is a bounded operator. The reader interested solely in this case might notice that the proofs of the statement could have been more direct on its own.

\begin{coro}
Let $\BA = \BA_0 + \BB$ be a self-adjoint unbounded periodic operator such that $\BA(\bk)$ has discrete spectrum for all $\bk \in \CO^\dagger$ and such that $\BB$ is bounded. Then, there exist labelings $\lambda_j(\BA_0(\bk))$ and $\lambda_j(\BA(\bk))$ of the eigenvalues of the fibre operators such that the functions $\lambda_j(\BA_0(\cdot))$ and $\lambda_j(\BA(\cdot))$ are both continuous on $\CO^\dagger$ and such that for every $\bk \in \CO^\dagger$
\begin{equation}
\left|\lambda_j(\BA_0(\bk) - \lambda_j(\BA(\bk))\right| \le \norm{\BB}.
\end{equation}
\end{coro}


\begin{proof}
It suffices to observe that $\norm{\BB(\bk)} \le \norm{\BB}$ for all $\bk \in
\CO^\dagger$. Defining the continuous family of operators $\BA_t = \BA_0 + t \BB$, it is easy to see that $\BA_1 = \BA$ and $\norm{\BA_t - \BA_s} = \abs{t-s}\norm \BB$. From Lemma~\ref{lem:monotonicity}, we know that for all $I \subset \R$, \eqref{eq:condids} holds for $\BA_s,\BA_t$ with $\eps = \abs{t-s}\norm \BB$. It is also clear that $\delta$ defined in~\eqref{eq:smalldist} is continuous in the parameter $t$. Setting
\begin{equation}
N = \ceil{\frac{\norm \BB}{\min_{0\,\le\,t \,\le\,1} \eps_t}},
\end{equation}
and applying recursively Lemma~\ref{lem:overlapidscont} to the operators $\BA_{j/N}$ and $\BA_{(j+1)/N}$, with $0 \le j < N$ yields the result we seek.
\end{proof}

The previous lemma and corollary provide us with an explicit way to compare the overlap function. This is made precise in the following proposition.

\begin{prop}\label{prop:overlapcontev}
Suppose that $\BA_1,\BA_2$ are self-adjoint, periodic operators such that for all $p \in \set{1,2}$ and $\bk \in \CO^\dagger$ the operator $\BA_p(\bk)$ has discrete spectrum. Suppose that for $\eps > 0$ there is a non-decreasing labelling of their eigenvalues so that whenever $\lambda_j(\BA_p(\bk)) \in [\lambda - 4\zeta(\lambda;\BA_1), \lambda + 4\zeta(\lambda;\BA_1)]$ we have
\begin{equation}\label{eq:evlabelclose}
\left|\lambda_j(\BA_1(\bk)) - \lambda_j(\BA_2(\bk))\right| \le \eps.
\end{equation}
Then,
\begin{equation}
\zeta(\lambda;\BA_2) \ge \zeta(\lambda;\BA_1) - 2\eps.
\end{equation}
\end{prop}

\begin{proof}
If $2\eps > \zeta(\lambda;\BA_1)$, the result follows trivially from nonnegativity of the overlap function. Otherwise, choose $j \in \Z$ such that
\[
\left[\lambda -
\zeta(\lambda;\BA_1),\lambda + \zeta(\lambda;\BA_1)\right] \subset \iota_j(\BA_1).
\]
Then~\eqref{eq:evlabelclose} implies
\begin{equation}
[\lambda - \zeta(\lambda;\BA_1) + \eps,\lambda + \zeta(\lambda;\BA_1) - \eps]
\subset \iota_j(\BA_1)_{-\eps} \subset \iota_j(\BA_2).
\end{equation}
The claim then follows from inspection of the definition~\eqref{eq:defnoverlap} of the overlap function.
\end{proof}

\section{Systems of periodic operators -- The Bethe--Sommerfeld property}\label{sec:bs}


In this section we prove that certain systems of periodic operators enjoy the Bethe--Sommerfeld property in a quantitative way. This will imply that the spectrum of an elliptic periodic operator $\BA$ in some classes contains a half-line. Our proof is again based on a reduction of the problem to uncoupled operators.

It is clear that if we show that the overlap function~\eqref{eq:defnoverlap} is bounded away from $0$ at sufficiently large $\lambda$ for some operator $\BA$, then $\BA$ has the Bethe--Sommerfeld property. This is the strategy employed in~\cite{ParSob2010}, where the self-adjoint operators of the form
\begin{equation}
A = (-\Delta)^\alpha + B,
\end{equation}
with $B\in \BS^\beta$, $\beta < 2\alpha$, and $B$ is $\Lambda$-periodic are studied. It is shown in~\cite{ParSob2010} that there are $S, c$ and $\lambda_0$, depending only on $\Theta$ and the symbol norms of $B$, such that for all $\lambda \ge \lambda_0$, $\zeta(\lambda;H) \ge c
\lambda^S$.

It is clear, since the spectrum of a finite direct sum of operators is the union of their individual spectra, that a direct sum of operators enjoying the Bethe--Sommerfeld property also enjoys the Bethe--Sommerfeld property. Nevertheless, the passage from scalar operators to uncoupled operators is not as easy as in Proposition~\ref{thm:aexpconcreteu} where the density of states of a direct sum of operators is readily seen to be the sum of the density of states of the summands. Indeed, when we try to establish the Bethe--Sommerfeld property we need the direct sum not only to have half-rays in its spectrum but also to preserve good lower bounds on the overlap function, otherwise the reduction to uncoupled operators could be able to open gaps. While it is possible a priori (see Example~\ref{ex:overlapcrash}) that for some direct sum of operators there are no lower bounds on the overlap function in terms of the overlap functions of the summands, our aim is now to show that for our class of operators this does not happen.

\subsection{The Bethe--Sommerfeld property} Our main theorem concerning systems of periodic operators is the following.

\begin{theo} \label{thm:bs}
Suppose that $\BA \in \BE\BS^\alpha_m$, $\alpha > 0$ is a periodic, self-adjoint
\mpu{} operator with $\BA^\CD$ of the form~\eqref{eq:defh0} and $a_j \ne a_k$ whenever $j \ne k$. Then, there exist positive $\tilde \lambda, S,c$ such that
\begin{enumerate}
\item \label{it:bspos}
if $\BA$ is unbounded above, $[\tilde \lambda,\infty) \subset
\spec(\BA)$ and for every $\lambda \ge \tilde \lambda$, $\zeta(\lambda;\BA)\linebreak \ge c
\lambda^{-S}$;
\item \label{it:bsneg}
if $\BA$ is unbounded below, $(- \infty,-\tilde \lambda] \subset
\spec(\BA)$ and for every $\lambda \ge \tilde \lambda$, $\zeta(-\lambda;\BA) \linebreak\ge c
\lambda^{-S}$.
\end{enumerate}
The overlap exponent $S$ depends only on $\alpha$ and the dimension $d$. The parameters $\tilde \lambda$ and $c$ can be chosen uniformly in the symbol norms of $\BA$ and $\BA^{\CO\CD}$.
\end{theo}


\begin{rema}
Saying that the parameters are chosen uniformly in the symbol norms means that if $\BA$ and $\BA'$ are operators satisfying the conditions of Theorem~\ref{thm:bs} and for all $s,\ell$, $\|\BA'\|^{s}_{\ell} \le \|\BA\|^{s}_{\ell}$ and $\|\BA'^{\CO\CD}\|^{s}_{\ell} \le \|\BA^{\CO\CD}\|^{s}_{\ell}$ then the parameters obtained for $\BA$ also work for $\BA'$.
\end{rema}


The proof of Theorem~\ref{thm:bs} hinges on three observations.
\begin{itemize}
\item For uncoupled operators, it is sufficient to prove bounds on the overlap function provided the operator is bounded below; this is proven in Lemma~\ref{lem:semibounded}.
\item There is a lower bound on the overlap function for uncoupled semi-bounded elliptic operators; this is Proposition~\ref{thm:uncoupledbs}.
\item Given an elliptic operator $\BA$, we find an uncoupled elliptic operator $\BA'$ so that the overlap function of $\BA'$ provides a lower bound for the overlap function of $\BA$ at large values of $\lambda$; this is the content of Lemma~\ref{lem:uncoupledlower}.
\end{itemize}

\begin{lemm}\label{lem:semibounded}
Let $\BA \in \BU\BS^\alpha_m \cap \BS\BE\BS^\alpha_m$ is self-adjoint and periodic, and suppose that
\begin{equation}
\BA = \BA_+ \oplus \BA_-
\end{equation}
with $\BA_+$ semi-bounded below and $\BA_-$ semi-bounded above. Then, there is $\lambda_0 > 0$ so that for every $\lambda > \lambda_0$, $\zeta(\pm \lambda;\BA) = \zeta(\pm \lambda; \BA_\pm)$.
\end{lemm}

\begin{proof}
Since $\BA_-$ is semi-bounded above, there is some $\lambda_0 > 0$ such that $\spec(\BA_-) \cap (\lambda_0,\infty) =
\varnothing$. It therefore follows that for $\lambda > \lambda_0$, and every $\bk \in
\CO^\dagger$
\begin{equation}
\spec(\BA_+(\bk)) \cap [\lambda,\infty) = \spec(\BA(\bk)) \cap [\lambda,\infty)
\end{equation}
and $\zeta(\lambda;\BA) = \zeta(\lambda;\BA_+)$ follows from the definition of the overlap function. Replacing $\BA$ with $- \BA$ in the argument, this also proves that $\zeta(-\lambda;\BA) = \zeta(-\lambda;\BA_-)$ as soon as $\lambda >
\lambda_0$, up to maybe increasing the value of $\lambda_0$.
\end{proof}


\begin{prop}\label{thm:uncoupledbs}
Let $\BA \in \BU\BS^\alpha_m \cap \BE\BS_m^\alpha$, $\alpha > 0$ be periodic, essentially self-adjoint, semi-bounded below operator with principal part $\BA_0^\CD$ of the form~\eqref{eq:defh0}. Then, there exist $\tilde \lambda, S,c > 0$ such that the interval $[\tilde \lambda,\infty) \subset \spec(\BA)$ and for every $\lambda \ge \tilde \lambda$, $\zeta(\lambda;\BA) \ge c \lambda^{-S}$. The overlap exponent $S$ depends only on $\alpha$ and $d$. The parameters $\tilde \lambda$ and $c$ can be chosen uniformly in the symbol norms of $\BA$ and $\BA^{\CO\CD}$.
\end{prop}


The proof of this proposition is very involved technically and uses some precise estimates in combinatorial geometry. Section~\ref{sec:cg} is entirely dedicated to the proof of this statement. We can deduce immediately from it and Lemma~\ref{lem:semibounded} the following corollary.


\begin{coro}\label{cor:uncoupledbs}
Let $\BA \in \BU\BS^\alpha_m \cap \BE\BS_m^\alpha$, $\alpha > 0$ be periodic, essentially self-adjoint operator with principal part $\BA_0^\CD$ of the form~\eqref{eq:defh0}. Then, there exist $\tilde \lambda, S,c > 0$ such that
\begin{enumerate}
\item \label{it:bsposuc}
if $\BA$ is unbounded above, the interval $[\tilde \lambda,\infty) \subset
\spec(\BA)$ and for every $\lambda \ge \tilde \lambda$, $\zeta(\lambda;\BA) \ge c
\lambda^{-S}$;
\item \label{it:bsneguc}
if $\BA$ is unbounded below, the interval $(- \infty,-\tilde \lambda] \subset
\spec(\BA)$ and for every $\lambda \ge \tilde \lambda$, $\zeta(-\lambda;\BA) \ge c
\lambda^{-S}$.
\end{enumerate}
The \emph{overlap exponent} $S$ depends only on $\alpha$ and $d$. The parameters $\tilde \lambda$ and $c$ can be chosen uniformly in the symbol norms of $\BA$ and $\BA^{\CO\CD}$.
\end{coro}



\begin{lemm}\label{lem:uncoupledlower}
Suppose that $\BA \in \BE\BS^\alpha_m$, $\alpha > 0$ is a periodic, self-adjoint, completely uncoupleable operator with principal part $\BA_0^\CD$ of the form~\eqref{eq:defh0} and $a_j \ne a_k$ whenever $j \ne k$. Then, for every $K
\in \R$ there exists an operator $\BA_K \in \BU\BS_m^\alpha \cap
\BS\BE\BS_m^\alpha$ periodic, essentially self-adjoint and with principal part $\BA_{K,0}^\CD = \BA_0^\CD$ such that for every $\abs \lambda$ large enough, we have that
\begin{equation}\label{eq:bsequiv}
\zeta(\lambda;\BA)
\ge
\zeta(\lambda;\BA_K) + \bigo{\abs{\lambda}^{-K}}.
\end{equation}
\end{lemm}

\begin{proof}
From Remark~\ref{rem:Iinfinitybs} and Theorem~\ref{thm:ctou}, for any $K$ there are $c_K, \lambda_K > 0$ and an operator $\BA_K \in
\BU\BS_m^\alpha$ such that for all $\lambda > \lambda_K$ and for $\eps_K := c_K \lambda^{-\alpha - K}$ and any interval $I \subset [\lambda - 2 \zeta(\lambda;\BA),\lambda + 2
\zeta(\lambda;\BA)]$ we have that for every $\bk \in \CO^\dagger$,
\begin{equation}\label{eq:idsclose2}
N(I;\BA(\bk)) \le N(I_{\eps_K};\BA_K(\bk)) \quad \text{and} \quad N(I;\BA_K(\bk)) \le N(I_{\eps_K};\BA(\bk)).
\end{equation}
Furthermore, we observe that the gauge transform leaves the principal part of elliptic operators untouched, so that the principal part of $\BA$ and $\BA_K$ coincide, as required.


As in~\eqref{eq:smalldist}, put
\begin{multline}
\delta_K =\\
\min_{\bk\,\in\,\CO^\dagger} \max\big\{\dist(\mu;\spec(\BA_K(\bk))) : \mu \in [\lambda - 2\zeta(\lambda;\BA_K),\lambda + 2\zeta(\lambda;\BA_K)]\big\} \\
 \ge C_K \frac{\zeta(\lambda;\BA_K)}{\max_\bk N([\lambda - 2\zeta(\lambda;\BA_K),\lambda + 2 \zeta(\lambda;\BA_K)];\BA_K(\bk))}
\end{multline}
for some $C_K > 0$. By Lemma~\ref{lem:overlapidscont} and Proposition~\ref{prop:overlapcontev}, if $\eps_K < \delta_K/4$, then~\eqref{eq:idsclose2} implies~\eqref{eq:bsequiv}. It follows from Corollary~\ref{cor:uncoupledbs} that $\zeta(\lambda;\BA_K) \ge c'_K\lambda^{-S}$ for some $S$ independent of $K$. Weyl's law implies that there is $C'_K$ such that
\begin{equation}
\max_\bk N([\lambda - 2
\zeta(\lambda;\BA_K),\lambda + 2\zeta(\lambda;\BA_K)];\BA_K(\bk)) \le C'_K \lambda^{d/\alpha}.
\end{equation}
It follows that by choosing $K > S + \frac d \alpha$, we have $\eps_K = c_K
\lambda^{-\alpha - K} < \delta_K/4$ for $\lambda$ large enough depending only on that fixed choice of $K$ and the constants $c_K,c'_K,C_K,C'_K$ encountered along, finishing the proof of Lemma~\ref{lem:uncoupledlower}.
\end{proof}


Before proceeding with the proof of Proposition~\ref{thm:uncoupledbs}, we indicate how it can be used to prove Theorem~\ref{thm:bs}.

\begin{proof}[Proof of Theorem~\ref{thm:bs}]
We only prove the case where $\BA$ is unbounded above, the other case follows by replacing $\BA$ with $-\BA$. Following Lemma~\ref{lem:uncoupledlower}, we find for any $K > - \alpha$ an operator $\BA_K \in \BU\BS^\alpha_m \cap
\BS\BE\BS_m^\alpha$ satisfying the hypothesis of Corollary~\ref{cor:uncoupledbs} and unbounded above, so that $\zeta(\lambda;\BA) \ge
\zeta(\lambda;\BA_K) + (\lambda^{-K})$ as soon as $\lambda$ is large enough. In turn, it follows from Corollary~\ref{cor:uncoupledbs} that $\zeta(\lambda;\BA_K) \ge c_K \lambda^{-S}$ for some $S$ depending only on $d$ and $\alpha$. Choosing any $K > S$ provides exactly the statement of Theorem~\ref{thm:bs}.
\end{proof}



\section{Bethe--Sommerfeld for uncoupled operators} \label{sec:cg}

In this section we prove Proposition~\ref{thm:uncoupledbs} --- that semi-bounded below self adjoint elliptic periodic operators enjoy the Bethe--Sommerfeld property. The strategy is an adaptation of the ideas found in~\cite{barbpar,parnovski,ParSob2010}, adapted to uncoupled systems rather than operators acting on scalar functions. Of course, since the systems are uncoupled we are able to borrow some of the results from the scalar theory and apply them to individual summands. Since the method is somewhat involved, we first give a heuristic description of the various steps; there we also indicate why we cannot recover Proposition~\ref{thm:uncoupledbs} from individually invoking results that are known for each of the summands comprising $\BA = \BA_1
\oplus \dotso \oplus \BA_m$.

\subsection{Heuristic approach}

Since $\BA$ is semi-bounded below, we do not have to worry about where the labelling of the eigenvalues in each fiber starts: for each $\bk \in \CO^\dagger$ we can write $\spec(\BA(\bk))$ as an increasing sequence accumulating only at infinity:
\begin{equation}
\lambda_1(\BA(\bk)) \le \lambda_2(\BA(\bk)) \le \dotso \nearrow \infty,
\end{equation}
with the function $\bk \mapsto \lambda_\ell(\BA(\bk))$ being piecewise analytic. Following the definition of the overlap function in~\eqref{eq:defnoverlap}, finding a lower bound $\delta := \delta(\lambda)$ for $\zeta(\lambda;\BA)$ means finding some $\ell \in \N$ so that the image of the function $\lambda_\ell(\BA(\bullet)) : \CO^\dagger \to \R$ contains an interval of radius $\delta$ around $\lambda$.

In order to find this interval, we aim at using the pigeonhole principle: for every $\lambda$ large enough we find $\bk_1, \bk_2 \in \CO^\dagger$ such that the counting functions at $\bk_1,\bk_2$ satisfy
\begin{equation}\label{eq:characoverlap}
N((-\infty,\lambda + \delta);\BA(\bk_1)) < N((-\infty,\lambda - \delta);\BA(\bk_2)).
\end{equation}
From this we deduce the existence of some $\ell$ such that $\lambda_\ell(\BA(\bk_2)) < \lambda - \delta < \lambda + \delta <\lambda_\ell(\BA(\bk_1))$. For every $\ell \in \N$ the functions $\bk \mapsto
\lambda_\ell(\BA(\bk))$ are continuous, so that we immediately deduce that the $\ell^{\text{th}}$ band has radius $\delta$.

We now exhibit two examples. The first one is there to show that the overlap function does not necessarily remain bounded away from zero when taking the direct sum of two operators with overlap function bounded away from zero.
\begin{exem}\label{ex:overlapcrash}
Consider the $1 \times 1$ family of matrices $A_1(\bk) = (\sin^2 \bk)$ and $A_2 = (\cos^2 \bk)$ indexed by $\bk \in \R/2\pi\Z$. Of course, for each $\bk$ each of those matrices has exactly one eigenvalue, and the overlap function associated to each is
\begin{equation}
\zeta(\lambda;A_j) =
\begin{cases}
1/2 - \abs{\lambda - 1/2}& \text{if } \abs{\lambda - 1/2} < 1/2 \\
0 & \text{otherwise.}
\end{cases}
\end{equation}
However, if we consider the $2 \times 2$ family of matrices given by $A = A_1(\bk)
\oplus A_2(\bk)$; they have two eigenvalues given by $\lambda_1(A) =
\min(\sin^2\bk,\cos^2 \bk)$ and $\lambda_2(A) =
\max(\sin^2 \bk,\cos^2 \bk)$. In particular, $\zeta(1/2,A) = 0$. This shows that a direct sum of operators can have zero overlap function without having a gap in its spectrum.
\end{exem}

The second example is a simple computation of the overlap function for the Laplacian. Some of the features of this computation will be present in the general case.
\begin{exem}\label{ex:heuristic}
This computation was first performed in~\cite{dahlbergtrubowitz} for the 2D Laplacian and can be used to deduce a Bethe--Sommerfeld property when the perturbation has low enough order (in terms of the dimension). We reproduce the computation because our own strategy will exhibit similar salient features. Consider the Laplacian $-\Delta =: \BA \in \BU\BS_1^2$, whose spectrum is of course the interval $[0,\infty)$. It is still an instructive exercise to compute the overlap function to understand its inner spectral structure.

The Laplacian is periodic with respect to any lattice, so let us make our decomposition with respect to $2\pi\Z^d$ which has dual lattice $\Theta := \Z^d$. Then, writing $\Delta_\bk$ for the Laplacian acting in $\ell^2(\Theta;\C^m)$ we can write the spectrum as
\begin{equation}
\spec(-\BA) =
\bigcup_{\bk\,\in\,\R^d/\Z^d} \spec(-\Delta_\bk) = \bigcup_{\bk\,\in\,\R^d/\Z^d} \set{4\pi^2|\bn + \bk|^2 : \bn
\in \Z^d}.
\end{equation}
Writing $\bone_R$ for the indicator of a ball of radius $R$, we can write the counting functions at $\bk \in \R^d/\Z^d$ as
\begin{equation}
N((-\infty,\lambda);\BA(\bk)) = \sum_{\bn \,\in\,\Z^d}
\bone_{\sqrt{\lambda}}\left(4\pi^2 \abs{\bn + \bk}^2\right).
\end{equation}
The counting function $N$ is periodic in the variable $\bk$, so let
\begin{equation}
\widehat N((-\infty,\lambda);\BA)_\bm := \int_{\CO^\dagger} N((-\infty,\lambda);\BA(\bk)) \be_\bm(\bk) \de \bk
\end{equation}
be its Fourier coefficient at $\bm$.

Our first observation is that
\begin{equation}
\begin{aligned}
\hat N((-\infty,\lambda);\BA)_{\boldsymbol 0} &=
\int_{\CO^\dagger}\sum_{\bn\,\in\,\Z^d} \bone_{\sqrt{\lambda}} \left(4
\pi^2\abs{\bn + \bk}^2 \right) \de \bk \\
&= \int_{\R^d} \bone_{\sqrt\lambda}\left(4\pi^2 \abs \bk^2 \right)
\de \bk \\
&= \frac{\operatorname{Vol}(B(0,1))}{(2\pi)^d}\lambda^{d/2}.
\end{aligned}
\end{equation}
We have unfolded the sum of integrals over translates of the fundamental domain $\CO^\dagger$ into an integral over the whole of $\R^d$. A similar idea will be used in our general strategy later. Similarly, using well known formulae for the Fourier transform of the indicator of a unit ball in terms of the Bessel function $J_{d/2}$ we obtain
\begin{multline}
\hat N((-\infty,\lambda);\BA)_{\bm}\\
\begin{aligned}
&=\int_{\CO^\dagger}\sum_{\bn\,\in\,\Z^d} \bone_{\sqrt{\lambda}} \left(4
\pi^2\abs{\bn + \bk}^2 \right) \be_\bm(\bk) \de \bk \\
&= \int_{\R^d} \bone_{\sqrt \lambda} \left(4\pi^d \abs \bk^2 \right)
\be_\bm(\bk) \de \bk \\
&= \frac{\lambda^{d/2}}{\abs \bm^{d/2}} J_{d/2}(2 \pi \abs \bm
\lambda) \\
&= \frac{\lambda^{\frac{d-2}{4}}}{2 \pi \abs \bm^{\frac{d+1}{2}} }\sin\left(2 \pi
\lambda \abs \bm + \frac{1-d}{4} \pi\right) +
\bigo{\lambda^{\frac{d-4}{4}} \abs \bm^{\frac{-d-3}{2}}}.
\end{aligned}
\end{multline}
If $d \not \equiv 1 \mod 4$, the $\sin$ in the first term in the last equality is always bounded away from $0$, whereas if $d \equiv 1 \mod 4$ it can be made larger than $\lambda^{- \eps}$ for any $\eps > 0$ with an appropriate choice of $\bm$, see~\cite{ParSob2001}, so that in the end we have that for every $\lambda > 0$ and $\eps > 0$ there exists $\bm$ so that
\begin{equation}
\hat N((-\infty,\lambda);\BA)_\bm \gg \lambda^{\frac{d-2}{4} - \eps_d},
\end{equation}
where $\eps_d = \eps$ if $d \equiv1 \mod 4$ and $0$ otherwise. A simple application of the triangle inequality tells us that for every periodic function $f$,
\begin{equation}
\norm{f}_{\RL^1(\CO^\dagger)} \ge \norm{\hat f}_{\RL^\infty(\Lambda)}
\end{equation}
so that
\begin{equation}
\int_{\CO^\dagger} \abs{N((- \infty, \lambda);\BA(\bk)) -
\frac{\operatorname{Vol}(B(0,1))}{(2\pi)^d} \lambda^{d/2}} \de \bk \gg
\lambda^{\frac{d-2}{4} - \eps_d}.
\end{equation}
Since
\begin{equation}
\int_{\CO^\dagger}\left(N((-\infty,\lambda);\BA(\bk)) -
\frac{\operatorname{Vol}(B(0,1))}{(2\pi)^d} \lambda^{d/2}\right) \de \bk = 0,
\end{equation}
this means that there is $C > 0$ so that for every $\lambda > 0$ there exists $\bk_1,\bk_2$ so that
\begin{equation}
N((-\infty,\lambda);\BA(\bk_1)) \ge
\frac{\operatorname{Vol}(B(0,1))}{(2\pi)^d} \lambda^{d/2} + C
\lambda^{\frac{d-2}{4} - \eps_d}
\end{equation}
and
\begin{equation}
N((-\infty,\lambda);\BA(\bk_2)) \le
\frac{\operatorname{Vol}(B(0,1))}{(2\pi)^d} \lambda^{d/2} - C
\lambda^{\frac{d-2}{4} - \eps_d}.
\end{equation}
In particular, for any $S < \frac{2-d}{4}$, taking $\delta = \lambda^{S}$ we have the existence of $\bk_1, \bk_2 \in \CO^\dagger$ so that
\begin{equation}
N((-\infty, \lambda + \delta);\BA(\bk_1)) < N((-\infty, \lambda -
\delta);\BA(\bk_2)),
\end{equation}
in other words the Laplacian has overlap exponent at least $\frac{2-d}{4} -
\eps_d$.
\end{exem}

Before delving into the more general case, let us discuss a key feature of the previous example. In order to describe the spectrum of each $\BA_j(\bk)$, we found functions $f_j : \R^d \to \R$ so that for every $\bk \in \CO^\dagger$ the restriction of $f_j$ to the fibre $\bk + \Theta$ is a bijection to $\spec(\BA_j(\bk))$. This works as easily for other diagonal operators, but not for operators whose symbol depend on the space variable $\bx$. One of the main purposes of using the gauge transform is to make it so that the operator depends, in a way, as little on $\bx$ as possible.

\begin{rema}[Remark on asymptotic notation]
In the remainder of this section we make extensive use of the Landau asymptotic notation:
\begin{itemize}
\item $f \ll g$ or $g \gg f$ to mean that there exists a constant $C$ such that $\abs f \le C \abs g$;
\item $f \asymp g$ to mean $f \ll g$ and $g \ll f$.
\end{itemize}
The implicit constants will always be allowed to depend on the numbers $a_1,\,\dotsc,\,a_m$ in the definition of the principal symbol of $\BA$, on the dimension and on the symbol norms of the subprincipal part. Other dependencies will be mentioned explicitly.
\end{rema}

\subsection{Reduction of the operator and description of the spectrum} \label{sec:constr}

Our first step towards proving Proposition~\ref{thm:uncoupledbs} is to use the gauge transform to reduce the operator to a form that is more manageable. We use $\rho := \lambda^{1/\alpha}$ as a radius in momentum space. In~\cite[Theorem~4.3 and Lemma~10.4]{ParSob2010}, large frequencies are cut out and a gauge transform is applied to deduce that for each summand $\BA_j$, for every $L > 0$, there is an operator $\tilde \BA$ with symbol
\begin{equation}\label{eq:model}
\tilde \ba =
\tilde \ba_0^\CD + \tilde\bbb^\CD + \tilde \bbb^\CR,
\end{equation}
where $\tilde \bbb^\CD$ is a diagonal subprincipal part and $\tilde \bbb^\CR$ is the resonant part, so that:
\begin{itemize}
\item The Fourier coefficients $\tilde \bbb_{\btheta}(\bxi)$ are supported in an annulus $c \rho < \abs \bxi < c\rho$, and in resonant regions $\Lambda_{\btheta}^{\nu,2^{-4}}$ as defined in Section~\ref{sec:gt}. Here, $c,C > 0$ depend on the numbers $a_j$ and $\nu \in (0,1)$ depends on the order of the perturbation. Furthermore, the frequency set $\Theta_\kappa$ for $\tilde \bbb$ is contained in a ball of radius $\rho^\kappa$ for some $0 < \kappa <
\min(d^{-2},d^{-1}(1 -\nu))$.
\item The overlap functions satisfy $\zeta(\BA;\lambda) > \zeta(\tilde \BA) + (\lambda^{-L})$.
\end{itemize}
This means that it is sufficient to prove that $\tilde \BA$ with symbol as in~\eqref{eq:model} has an overlap function bounded below by a power of $\lambda$. From now on, we assume that this reduction has been done and abuse notation by omitting the tildes and writing the model operator as $\BA$.

\begin{rema}\label{rmk:kappanufixed}
We assume for the rest of the section that $\nu$ and $\kappa$ as obtained above are fixed, and that $\kappa d < 1 - \nu$. We also always fix $\Theta_\kappa$ to be the frequency set for $\bbb$.
\end{rema}


Just as in Example~\ref{ex:heuristic}, we want to 'unfold' the spectrum of each fibre operators $\BA_j(\bk)$ by assigning one eigenvalue to each translate of $\CO^\dagger$ by an element of $\Gamma^\dagger$. This gives us $m$ functions $g_j : \R^d \to \spec(\BA_j)$ so that the restriction of $g_j$ to every fibre $\bk + \CO^\dagger$ gives a bijection to $\spec(\BA_j(\bk))$. Of course, outside the union of the resonant regions, we see that
\begin{equation}
\BA (e_{\bxi} \otimes v_j) = \left(a_j \abs \bxi^\alpha + \bbb_j^\CD(\bxi)\right) e_{\bxi} \otimes v_j
\end{equation}
so that it makes sense to put there $g_j(\bxi) = a_j \abs \bxi^\alpha +
\bbb_j^\CD(\bxi)$. The definition of these functions $g_j$ is given in~\cite[Section~7]{ParSob2010}, where some of their properties are studied. We collect the important ones for our purpose here.

\begin{lemm}\label{lem:descriptiongj}
For every $1 \le j \le m$ there exists a function $g_j : \R^d \to
\spec(\BA_j)$ so that the restriction of $g_j$ to every fibre $\bk +
\CO^\dagger$ is a bijection (respecting multiplicity) satisfying the following properties:
\begin{itemize}
\item We can write $g_j(\bxi) = a_j \abs \bxi^\alpha + G_j(\bxi)$, with $\abs{G_j(\bxi)} \ll \rho^\beta$ whenever $\abs \xi \asymp \rho$, for some $\beta < \alpha$.
\item For $\delta \in (0,\rho^\alpha/4)$ define the ``annular'' regions
\begin{equation}
\CA_j := \CA_j(\rho;\delta) := g_j^{-1}\left([\rho^\alpha - \delta,
\rho^\alpha + \delta]\right),
\end{equation}
which by the previous item we know is contained in a genuine annulus $\abs \bxi \asymp \rho$. Put
\begin{equation}
\CR_j := \CA_j \cap \bigcup_{\theta\,\in \,\Theta\,\cap\,B(\rho^\kappa)} \Lambda_{\btheta}^{\gamma,1} \quad \text{ and }
\quad
\CB_j := \CA_j \setminus \CR_j.
\end{equation}
Then, on $\CB_j$, the functions $G_j$ are of class $\RC^2$ and for every $\bxi \in \CB_j$ we have the estimate
\begin{equation}
\abs{\nabla G_j(\bxi)} + \rho \abs{\nabla^2 G_j(\bxi)} \ll
\rho^\gamma
\end{equation}
for some $\gamma< \alpha - 1$.
\end{itemize}
\end{lemm}

\subsection{Geometry of the resonant regions}

The main takeaway from Lemma~\ref{lem:descriptiongj} is that although the perturbations $G_j$ are small with respect to $a_j \abs \bxi^\alpha$ everywhere, their derivatives are only controlled within the non-resonant region. As such, we would like the non-resonant regions to be as large as possible, and as such we now turn our attention to studying their geometry.

We first introduce some notation. The sets we describe depend on both parameters $\rho$ and $\delta$ used to define $\CA_j$; however we often drop the (explicit) dependence on these parameters to make notation lighter. For every $\bxi \in \R^d \setminus \set 0$, let $\bu_{\bxi} := \abs \bxi^{-1}\bxi$ be the unit vector in the direction of $\bxi$. For any subset $\CU$ of the sphere $\BBS^{d-1}$, we denote its radial extension by
\[
\CU_{rd} := \set{\bxi \in \R^d : \bu_{\bxi} \in \CU}.
\]


For $T \in (0,1)$ and $\btheta \in
\Theta \setminus \set 0$, we define spherical resonant regions as
\begin{equation}
\CS(\btheta;T) := \set{\bzeta \in \BBS^{d-1} : \abs{\bzeta \cdot \bu_{\btheta}} < T}, \qquad \text{and} \qquad \CS(T) = \bigcup_{\theta\,\in\,\Theta_\kappa} \CS(\btheta;T).
\end{equation}
The name is justified from the fact that by elementary trigonometry, we have the inclusion for the resonant regions
\begin{equation}\label{eq:crjsub}
\CA_j \cap \Lambda^{\nu,2^{-4}}_{\btheta} \subset \CS\left(\btheta,\rho^{\nu - 1}\right)_{rd} \qquad
\text{so that} \qquad
\CR_j(\rho;\delta) \subset \CS\left(\rho^{\nu - 1}\right)_{rd},
\end{equation}
where $\nu$ is as fixed in Remark~\ref{rmk:kappanufixed}. In particular, we also define a spherical non-resonant region as
\begin{equation}
\CT(\rho) := \BBS^{d-1} \setminus \CS\left(\rho^{\nu - 1}\right)\qquad \text{ so that } \qquad
\tilde \CB_j := \CT_{rd} \cap \CA_j \subset \CB_j.
\end{equation}

The objective of the next lemma is to prove that small enough neighborhoods of the spherical resonant regions have small volume.
\begin{lemm} \label{lem:enlargedresonant}
Let $\varkappa$, $\Theta_\varkappa$ and $\nu$ be as fixed by Remark~\ref{rmk:kappanufixed} and define
\begin{equation}
\tilde \CT(\rho) := \BBS^{d-1} \setminus \CS\left(2 \rho^{\nu-1}\right),
\end{equation}
\begin{equation}\label{eq:czjsub}
\CZ_j(\rho;\delta) := \CA_j \cap \CS_{rd}\left(2\rho^{\nu-1}\right),
\end{equation}
and
\begin{equation}\label{eq:cgjsub}
\CG_j(\rho;\delta) := \CA_j \cap \tilde \CT_{rd} = \CA_j \setminus \CZ_j.
\end{equation}

Then, $\CG_j \subset \CB_j$, and for all $\bxi \in \CG_j$, $\dist(\bxi,\CR_j)\gg
\rho^{\nu}$. Furthermore, for $\eps_0 = 1 - \nu - d \varkappa > 0$,
\begin{equation}\label{eq:volCZ}
\vol(\CZ_j(\rho;\delta)) \ll \delta \rho^{d - \alpha - \eps_0},
\end{equation}
and
\begin{equation}\label{eq:voltCG}
\vol(\CG_j(\rho;\delta)) \asymp \delta \rho^{d - \alpha}.
\end{equation}
\end{lemm}

\begin{proof}
It is clear from the definition that $\CG_j \subset \tilde \CB_j \subset
\CB_j$. For $\bzeta \in \BBS^{d-1}$, define the sets
\begin{equation}
\CI_j(\bzeta) := \CA_j \cap \set{\bzeta}_{rd}.
\end{equation}
It follows from~\cite[Equation (8.5)]{ParSob2010} that for all $\bzeta \in
\CT$, the interval $\CI_j(\bzeta)$ is an interval of length $\abs{\CI_j} \ll \delta
\rho^{1 - \alpha}$ (uniformly in $\bzeta$), and
\begin{equation}
\CI_j \subset \set{\bxi : \abs \bxi \asymp \rho },
\end{equation}
furthermore, by definition,
\[
\dist\left(\tilde \CT(\rho), \CS\left(\btheta;\rho^{\nu -1}\right)\right) > \rho^{\nu -1 }.
\]
It therefore follows from~\eqref{eq:crjsub}, \eqref{eq:cgjsub} and basic trigonometry that $\dist(\CG_j,\CR_j) \gg \rho^{\nu}$. For the volume estimate for $\CZ_j$, we compute
\begin{equation}\label{eq:volczj}
\begin{aligned}
\vol(\CZ_j) &\le \sum_{\btheta\,\in\,\Theta_\varkappa} \int_{
\CS\left(\btheta;2 \rho^{\nu - 1}\right)} \int_{\CI_j(\bzeta)} t^{d-1} \de t \de \bzeta \\
& \ll \#(\Theta_\varkappa)\max_{\btheta}\vol_{d-1}\left(
\CS\left(\btheta;2\rho^{\nu-1}\right)\right)\delta\rho^{d - \alpha}.
\end{aligned}
\end{equation}
Uniformly in $\btheta$ we have $\vol_{d-1}(
\CS(\btheta;2\rho^{\nu-1}) \ll \rho^{\nu-1}$. We also have that $\#\Theta_\varkappa \ll
\rho^{d\varkappa}$. Putting these two estimates in~\eqref{eq:volczj} yields~\eqref{eq:volCZ}. For the estimate on $\vol(\CG_j(\rho;\delta))$, we observe that $\vol(\CG_j) = \vol(\CA_j) - \vol(\CZ_j)$ and that by Lemma~\ref{lem:descriptiongj},
\begin{equation}
\vol(\CA_j) \asymp \delta \rho^{d - \alpha}.
\end{equation}
Estimate~\eqref{eq:voltCG} then follows from the fact that $\vol(\CZ_j) =
\smallo{\vol(\CA_j)}$.
\end{proof}

For our purposes, we need not only to have volume estimates on the resonant and non-resonant regions, but also on intersections of their translates. For $\bb_1,\bb_2 \in \R^d$ and $i,j \in \set{1,\,\dotsc,\,m}$ we define the
\emph{crossing sets}
\begin{equation}\label{eq:defCX}
\CX_{ij}(\rho,\delta,\bb_1,\bb_2) := (\CA_i(\rho;\delta)+\bb_1)\cap(\CA_j(\rho;\delta)+\bb_2).
\end{equation}
We are interested in volume estimates, and since
\begin{equation}
\vol(\CX_{ij}(\rho,\delta,\bb_1,\bb_2)) =
\vol(\CX_{ij}(\rho,\delta,\bzero,\bb_2-\bb_1)),
\end{equation}
we restrict ourselves to sets of the form
\begin{equation}
\CX_{ij}(\bb) := \CX_{ij}(\rho,\delta,\bb) :=
\CX_{ij}(\rho,\delta,\bzero,\bb).
\end{equation}
Denote by $\phi(\ba,\bb)$ the (smaller) angle between $\ba$ and $\bb$. For any angle $\omega \in [0,\pi]$, we define the set
\begin{equation}\label{eq:badvol}
\CX_{ij,\omega}(\bb) := \set{\bxi \in \CX_{ij}(\bb) : \phi(\bxi,\bxi-\bb) >
\omega}.
\end{equation}

In our application, we need to control the volume of crossing sets for angles bounded away from zero. The next proposition tells us that unless $\bbb$ is comparable in size to $\rho$, this volume is zero.
\begin{prop}\label{rem:xemptyint}
Let $\delta > 0$ and $\omega \in (0,\pi)$. There are $c$ and $C$, also depending on numbers $a_j$ such that $\CX_{ij,\omega}(\rho,\delta,\bb) \ne \varnothing$ implies that $c \rho \le \abs\bb \le C\rho$.
\end{prop}

\begin{proof}
We first make the observation that there exists $C > 0$, depending on $\alpha$ and the numbers $a_j$ such that if $\abs \bbb > C\rho$, then for $\rho$ large enough $\CA_i \cap (\CA_j + \bbb) = \varnothing$. On the other hand, it follows from basic planar trigonometry that for every $\omega$, there exists $c$, depending on the constants in $\abs\bxi \asymp \rho$, such that if $\abs \bbb < c \rho$, then $\phi(\bxi, \bxi-\bbb)\le \omega$.
\end{proof}

It follows from Proposition~\ref{rem:xemptyint} and the results in~\cite[Section~9]{ParSob2010} that the following holds: for any $\omega \in (0,\pi)$, $\eps > 0$, if $\delta \rho^{2 - \alpha + 2 \eps} \to 0$ as $\rho \to \infty$, then
\begin{equation}\label{eq:nokissmu}
\vol\left(\CX_{ij,\omega}(\rho,\delta,\bb)\right) \ll \delta^2 \rho^{4 - 2\alpha + d + 6\eps} + \delta \rho^{1 - \alpha - \eps(d-1)},
\end{equation}
uniformly in $\bb$. We now define crossing sets for the non-resonant sets $\CG_j$. For $\bbb \in
\R^d$, let
\begin{equation}
\CY_{ij}(\bb) := \CG_{i} \cap (\CG_j +
\bb)
\end{equation}
and for any angle $\omega \in (0,\pi)$,
\begin{equation}
\CY_{ij,\omega}(\bb) := \set{\bxi \in \CY_{ij}(\bbb) : \phi(\bxi,\bxi-\bb) > \omega} = \CX_{ij,\omega}(\bbb) \cap \CY_{ij}(\bbb).
\end{equation}

Before going on, let us make the following notational convention.

\begin{enonce*}
{Convention}
For any family of subsets $\CE(\delta) \subset \R^d$ depending on the parameter $\delta > 0$, we denote
\[
\CE'(\delta) := \CE'(\delta,Z) := \CE(Z\delta),
\]
where $Z$ is some large constant to be determined later and depending only on the dimension $d$, the order $\alpha$ and the numbers $\set{a_1,\,\dotsc,\,a_m}$.
\end{enonce*}

We need the following lemma.
\begin{lemm} \label{lem:totalintersection}
For any $\omega \in (0,\pi)$ and $\eps > 0$, the condition $\delta \rho^{2 -
\alpha + 2\eps} \to 0$ as $\rho \to \infty$ implies
\begin{equation}
\vol\left(\bigcup_{i,j = 1}^m \bigcup_{\btheta\,\in\,\Theta}
\CY_{ij,\omega}'(\btheta) \right) \ll
\delta^2 \rho^{4 - 2\alpha + 2 d +6\eps} + \delta \rho^{1- \alpha + d - \eps(d-1)},\label{eq:sobranie}
\end{equation}
the implicit constants depending only on $\delta$, $\omega$, $Z$, and the coefficients $a_j$.
\end{lemm}

\begin{proof}
It is sufficient to prove the result for a single pair $i,j$, then sum the estimates over all $m^2$ of those pairs. From Proposition~\ref{rem:xemptyint}, there are constants $c$ and $C$ depending only on $\omega,\delta,T$ and the numbers $a_j$ such that
\begin{equation}
\begin{aligned}
\vol\left(\bigcup_{\btheta\,\in\,\Theta}\CY_{ij,\omega}'(\btheta)\right) &\le \sum_{\substack{\btheta\,\in\,\Theta \\ c \rho \le \abs\btheta
\le C\rho } }\vol\left(\CY_{ij,\omega}(\btheta) \right)\\
&\ll\delta^2 \rho^{4 - 2\alpha + 2 d + 6 \eps} + \delta \rho^{1-
\alpha + d - \eps(d-1)},
\end{aligned}
\end{equation}
where the last line comes from $\CY_{ij,\omega}(\btheta) \subset
\CX_{ij,\omega}(\btheta)$, estimate~\eqref{eq:nokissmu}, and the fact that
\begin{equation}
\#\set{\btheta \in \Theta^\dagger : c \rho \le \abs \btheta \le C \rho} \ll
\rho^d.\qedhere
\end{equation}
\end{proof}


\subsection{Estimating the overlap function}
We are now ready to provide the estimate~\eqref{eq:characoverlap}. For this, we will find $\bk_1, \bk_2 \in \CO^\dagger$ and three types of eigenvalue branches $\lambda_\ell(\bullet)$:
\begin{enumerate}
\item\label{sect11.4.1} branches so that $\lambda_\ell(\bk_1) < \rho^\alpha - \delta$ and $\lambda_\ell(\bk_2) > \rho^\alpha + \delta$; these branches go across $[\rho^\alpha - \delta, \rho^\alpha + \delta]$;
\item\label{sect11.4.2} branches which we cannot control in any way, they may be above $\rho^\alpha - \delta$ at $\bk_1$ and below $\rho^\alpha + \delta$ at $\bk_2$;
\item\label{sect11.4.3} branches so that $\lambda_\ell(\bk_1)$ and $\lambda_\ell(\bk_2)$ are not in the interval $[\rho^\alpha - \delta,\rho^\alpha + \delta]$; this will follow from them being far enough from the interval at some midpoint.
\end{enumerate}
Branches of type~\eqref{sect11.4.1} contribute to inequality~\eqref{eq:characoverlap} whereas branches of type~\eqref{sect11.4.2} are the adversary; therefore, we aim at proving that there are strictly more branches of type~\eqref{sect11.4.1} than branches of type~\eqref{sect11.4.2}. Branches of type~\eqref{sect11.4.3} either contribute to both sides of inequality~\eqref{eq:characoverlap} or to neither, so that their number is irrelevant.

In order to achieve this goal, it is convenient to ``refold'' the functions $g_j$. Indeed, rather than considering them as global functions of $\bxi$, for each $\btheta \in \Theta$ we consider $g_j(\bk+\btheta)$ as a function of $\bk
\in \CO^\dagger$. More precisely, put $\tilde \Theta := \set{1,\,\dotsc,\,m}
\times \Theta$, and for every $\bp = (j,\btheta) \in \tilde \Theta$, put
\begin{equation}
g_\bp(\bk) := g_j(\bk + \btheta).
\end{equation}
By definition of $g$, for any $\rho > 0$ and any $\bk \in \CO^\dagger$ we have
\begin{equation}
N((-\infty, \rho^\alpha;\BA(\bk)) = \#\set{\bp \in \tilde \Theta : g_\bp(\bk) \le \rho^\alpha},
\end{equation}
so that we can directly study the functions $g_\bp$; it will be useful to go back and forth between their geometric description and the fact that these are still the eigenvalues of some operator.

The following lemma explains why it was important to show that intersections of non-resonant zones with large angles had small volume: when the angle is small, the functions $g_j$ are increasing and we will use this fact to construct the branches of type~\eqref{sect11.4.1}.
\begin{lemm} \label{lem:increasing}
Let $\nu$ be as fixed in Remark~\ref{rmk:kappanufixed}. For all $\bxi \in
\CG_j$, and all $\bbb$ such that $\abs{\bxi + \bbb} \asymp \rho$ and $\phi(\bxi,\bxi+\bbb) \le \pi/4$ there is a $t_0 \gg \rho^{-\nu}$ such that for all $t \in [-t_0,t_0]$:
\begin{itemize}
\item the point $\bxi + t(\bxi + \bbb)$ is in $\CB_j'$;
\item the function $t \mapsto g_j(\bxi + t(\bxi + \bbb))$ is increasing;
\item the derivative satisfies
\begin{equation}
\frac{\de}{\de t}g_j(\bxi + t(\bxi + \bbb)) \gg \rho^\alpha.
\end{equation}
\end{itemize}
The implicit constant in $t_0 \gg \rho^{-\nu}$ depends only on the functions $g_j$ and the implicit constants in $\abs{\bxi + \bbb} \asymp
\rho$.
\end{lemm}

\begin{proof}
Since $\bxi \in \CG_j$, we have that not only $\bxi \in \CB_j$ for some $j$, but also, by Lemma~\ref{lem:enlargedresonant}, that there exists $r > 0$ such that for all $j'$, we have $\dist(\bxi;\CR_{j'}) > r \rho^{1-\nu}$. Therefore, for $\abs t \le t_0 := r \rho^{-\nu}$, we have that $\bxi + t(\bxi + \bbb) \in \CB'_j$. By Lemma~\ref{lem:descriptiongj}, we have that
\begin{equation}
\abs{\frac{\de}{\de t} G_j(\bxi + t(\bxi + \bbb))} \ll \rho^{\gamma+1} = \smallo{ \rho^\alpha}.
\end{equation}
On the other hand,
\begin{align}
\frac{\de}{\de t}& \abs{\bxi + t(\bxi + \bbb)}^\alpha\\ &= \alpha
\abs{\bxi + t (\bxi + \bbb)}^{\alpha - 2} \left(\abs \bxi \abs{\bxi +
\bbb}\cos\left(\phi(\bxi,\bxi+\bbb)\right) + t \abs{\bxi+\bbb}^2 \right)\nonumber\\
& \gg \rho^\alpha,
\end{align}
where the last line holds from the fact that $\cos \phi(\bxi,\bxi + \bbb) > \sqrt{2}/2$.
\end{proof}

Branches of type~\eqref{sect11.4.2} will come from either $\bxi = \bk + \btheta$ which are in a resonant region, or where the angle between $\bxi$ and $\bxi + \btheta$ is large. The next two lemmas aim at controlling these two types of situations. But first, we introduce even more notation: define
\begin{equation}
\CZ := \bigcup_{1\,\le\,j\,\le\,m} \CZ_j, \qquad \qquad \text{and} \qquad \qquad
\CG := \bigcup_{1\,\le\,j\,\le\,m} \CG_j.
\end{equation}
It follows directly from the definitions of $\CZ_j$~\eqref{eq:czjsub} and $\CG_j$~\eqref{eq:cgjsub} that $\CG \cap \CZ = \varnothing$. Furthermore, for every $\bxi \in \CG$ and every $1 \le j \le m$, $\bxi \in \CG_j$ or $\bxi \not \in
\CA_j$. For every $\bxi \in \R^d$, and any subset $E \subset \R^d$, we define
\begin{equation}
n(\bxi,E) = \#\set{\btheta \in \Theta : \bxi + \btheta \in E}.
\end{equation}
\begin{lemm}\label{lem:numinter}
Let
\begin{equation}
\CN := \set{\bxi \in \R^d : n(\bxi;\CG) \le m n(\bxi;\CZ')}
\end{equation}
and $\CN_\CG = \CN \cap \CG$. Then, we have that
\begin{equation}
\vol(\CN_\CG) \ll \vol(\CZ') \ll
\delta\rho^{d - \alpha - \eps_0},
\end{equation}
where $\eps_0$ is as in Lemma~\ref{lem:enlargedresonant}.
\end{lemm}
\begin{proof}
Observe first that for any $E \subset \R^d$, the function $n(\bxi;E)$ is constant on the fibres $\bxi \mod \Theta$, in other words it depends only on the fractional part $\set \bxi$. This means that it is well defined on $\CO^\dagger$ and $\CN$ is invariant under the action of $\Theta$.\linebreak We therefore have that
\begin{align}
\vol(\CN_\CG) &= \int_{\CN/\Theta} n(\bxi;\CG) \de \bxi \\
&\le m \int_{\CN/\Theta} n(\bxi;\CZ') \de \bxi \\
& = m \vol(\CN \cap \CZ') \\
&\le m \vol(\CZ').
\end{align}
The claim now follows from Lemma~\ref{lem:enlargedresonant}.
\end{proof}

\begin{lemm}\label{lem:goodcrossings}
Let
\begin{multline}\label{eq:CU}
\CU :=\\
\big\{\bxi \in \CG \setminus \CN : \bxi + \btheta_1 \in
\CY_{ij,\pi/4}'(\btheta_1 - \btheta_2) \text{ for some } 1 \le i,j \le m
\text{ and } \btheta_1,\btheta_2 \in \Theta\big\}.
\end{multline}
Then, if $\eps$ and $\delta$ are such that $\delta \rho^{d - \alpha + 2 \eps} \to 0$, we have
\begin{equation}
\vol(\CU) \ll \delta^2 \rho^{4 - 2 \alpha + 2 d + 6\eps} + \delta \rho^{1 - \alpha + d - \eps(d-1)}.
\end{equation}
\end{lemm}

\begin{proof}
Suppose that $\bxi \in \CU \subset \CG$, so that $\bxi \in \CG_k$ for some $1 \le k \le m$. Consider the lattice elements $\btheta_1, \btheta_2 \in
\Theta^\dagger$ such that $\bxi + \btheta_1 \in \CY_{ij,\pi/4}'(\btheta_1 -
\btheta_2)$. By definition of $\CY_{ij,\pi/4}'$ and translation, this means that
\begin{equation}
\bxi \in (\CG_i' - \btheta_1) \cap (\CG_j' - \btheta_2),
\end{equation}
and therefore that
\begin{equation}
\bxi \in \CY_{ki}'(-\btheta_1) \cap \CY_{kj}'(- \btheta_2).
\end{equation}
Furthermore, $\phi(\bxi+\btheta_1,\bxi+\btheta_2) > \pi/4$. As such,
\begin{equation}
\max \set{\phi(\bxi,\bxi+\btheta_1),\phi(\bxi,\bxi+\btheta_2)} > \pi/8.
\end{equation}
Combining the previous two displays yields
\begin{equation}
\bxi \in \CY_{ki,\pi/8}'(-\btheta_1) \cup \CY_{kj,\pi/8}(-\btheta_2).
\end{equation}
Therefore,
\begin{equation}
\begin{aligned}
\vol(\CU) &\le \vol\left(\bigcup_{i,j = 1}^m \bigcup_{\btheta \in
\Theta} \CY'_{ij,\pi/8}(\btheta) \right) \\
&\ll \delta^2 \rho^{4 - 2 \alpha + 2d + 6\eps} + \delta \rho^{1 -
\alpha + d - \eps(d - 1)},
\end{aligned}
\end{equation}
the last line holding by virtue of Lemma~\ref{lem:totalintersection}.
\end{proof}

The next proposition indicates that the sets $\CU$ and $\CN_\CG$ are thin relative to $\CG$.

\begin{prop}\label{prop:overlap}
Let
\begin{equation}\label{eq:propoverlap}
s := \min \set{\frac{\alpha d - d^2 - 3d - \alpha - 2}{2(d+2)}, \alpha - d + \frac{\alpha - d - 2}{2(d+2)}}.
\end{equation}
For $\rho$ large enough and $\delta = \smallo{\rho^s}$, the set
\[
\CK:= \CG \setminus(\CN_\CG \cup \CU)
\]
is non empty.
\end{prop}
\begin{proof}
Recall from Lemma~\ref{lem:enlargedresonant} that $\vol(\CG) \asymp \delta
\rho^{d-\alpha}$. On the other hand, Lemma~\ref{lem:numinter} implies that there is $\eps_0 > 0$ such that
\begin{equation}
\vol(\CN_\CG) \ll \delta\rho^{d- \alpha - \eps_0},
\end{equation}
and from Lemma~\ref{lem:goodcrossings} that as soon as $\delta \rho^{d -
\alpha + 2 \eps} \to 0$ we have that
\begin{equation}
\vol(\CU) \ll \delta^2 \rho^{4 - 2 \alpha + 2 d + 6 \eps} + \delta \rho^{1 - \alpha + d - \eps(d-1)}.
\end{equation}
Take
\begin{equation}\label{eq:freefixed}
\eps := \frac{\alpha - d - 2}{2(d+2)}.
\end{equation}
Observe that indeed when $\delta = \smallo{\rho^s}$ we have $\delta \rho^{2 - \alpha + 2 \eps} \to 0$ as $\rho \to \infty$. We also observe that with that choice of parameters $\vol(\CU) + \vol(\CN_\CG) =
\smallo{\vol(\CG)}$ and hence, for large enough $\rho$, $\CK$ is not empty.
\end{proof}

We are now in a position to place our final building block towards the proof of Proposition~\ref{thm:uncoupledbs}.

\begin{prop}
There exists $\rho_0 > 0$ and $S \in \R$ (depending on $\{a_1,\dotsc,\linebreak a_m\}$, $\alpha$, and the implicit constants in Lemma~\ref{lem:descriptiongj}) so that
\begin{equation}\label{eq:countinggj}
\#\set{\bp \in \tilde \Theta : g_\bp(\bk_1) \le \rho^\alpha + \rho^S} <
\# \set{\bp \in \tilde \Theta : g_\bp(\bk_2) \le \rho^\alpha - \rho^S}.
\end{equation}
\end{prop}

\begin{proof}
Let $s$ be defined as in~\eqref{eq:propoverlap}. For any $\eps > 0$, set $S = \min \{s - \eps, \alpha + \nu\linebreak - 1\}$, where $\nu < 1$ is fixed in Remark~\ref{rmk:kappanufixed}. By Proposition~\ref{prop:overlap}, for $\rho$ large enough the set $\CK$ is not empty; so fix $\bxi_0 \in \CK$. For $1 \le j \le m$, let $\Gamma_j, \Gamma_j' \subset \Theta$ be defined as
\begin{equation}
\Gamma_j := \set{\btheta \in \Theta : \bxi_0 + \btheta \in \CG_j}
\qquad
\text{and} \qquad \Gamma_j' := \set{\btheta \in \Theta : \bxi_0 +
\btheta \in \CG_j'}.
\end{equation}
It follows from the definition of $\CK$ that
\begin{equation}
\sum_{j=1}^m \#\Gamma_j \ge n(\bxi_0;\CG).
\end{equation}
Since $\bxi_0 \not \in \CU$, for all $\btheta \in \Gamma_j$ we have $\phi(\bxi_0,\bxi_0 + \btheta) \le \pi/4$ for all $\btheta \in \Gamma_j$. It follows from Lemmas~\ref{lem:descriptiongj} and~\ref{lem:increasing}, since $\delta \ll \rho^{\alpha + \nu - 1}$, that there exist $t \ll \rho^{\nu - 1}$ and $Z_0$ independent of $\rho$ such that for all $1 \le j \le m$ and $\btheta
\in \Gamma_j$ we have
\begin{equation}\label{eq:botovlp}
\rho^\alpha - Z_0\delta \le g_j\big((1 - t) \bxi_0 + \btheta \big) \le
\rho^\alpha - \delta
\end{equation}
and
\begin{equation}\label{eq:topovlp}
\rho^\alpha + \delta \le g_j\big((1 + t) \bxi_0 + \btheta \big) \le
\rho^\alpha Z_0 \delta.
\end{equation}
Of course, if these estimates hold, they also hold replacing $Z_0$ with any $Z > Z_0$. The precise value we assign to $Z$ may change as the proof goes along but remains independent of $\rho$ and $\delta$. We denote by $\CJ$ the radial interval of length $\abs \CJ \ll \rho^{\nu - 1}\linebreak = \smallo 1$:
\begin{equation}
\CJ := [(1-t) \bxi_0, (1+t) \bxi_0],
\end{equation}
and put $\bk_0, \bk_1$, and $\bk_2$ to be the projections on $\CO^\dagger$ of $\bxi_0$, $(1-t)\bxi_0$. and $(1+t)\bxi_0$, respectively. We now restrict ourselves to the operator
\begin{equation}
\BA^\CJ := A(P_{\CJ + \Theta} \otimes \Id),
\end{equation}
where the projection $P_{\CJ + \Theta}$ is defined in Proposition~\ref{prop:normorder}. We also put $Q_j$ to be the projection on the $j$th coordinate in $\C^m$, so that $(\Id \otimes Q_j)(e_{\bxi}\otimes v_j) = e_{\bxi} \otimes v_j$ and $(\Id \otimes Q_j)(e_{\bxi} \otimes v_\ell) = 0$ for all $j \ne \ell$. From Lemma~\ref{lem:increasing}, for every $\bxi \in \CJ$ and $\btheta \in \Gamma_j$ we have that $\bxi + \btheta \in \CB_j'$. Since $\BA^\CJ$ acts diagonally on $\spann\set{e_{\bxi} \otimes v_j : \bxi \in \CJ +
\Gamma_j}$, it commutes with the projection
\begin{equation}
\BP = \sum_{j=1}^m P_{\CJ + \Gamma_j} \otimes Q_j.
\end{equation}
In particular, for every $\bk \in [\bk_1,\bk_2]$, the spectrum of $\BA(\bk)$ decomposes (respecting multiplicity) into
\begin{equation}
\spec(\BA(\bk)) = \spec(\BA(\bk) \BP(\bk)) \sqcup \spec(\BA(\bk) (\Id -
\BP(\bk))).
\end{equation}
It follows from the definition of $\Gamma_j$ and inequalities~\eqref{eq:botovlp} and~\eqref{eq:topovlp} that every eigenvalue of $\BA(\bk_1) \BP(\bk_1)$ is smaller than $\rho^\alpha - \delta$ and every eigenvalue of $\BA(\bk_2)
\BP(\bk_2)$ is larger than $\rho^\alpha + \delta$.

It follows from~\cite[Theorem~3.6]{ParSob2010} that there is $Z \ge Z_0$ large enough so that for every $\ell \in \N$,
\begin{multline}
\abs{\lambda_\ell(\BA(\bk_0) (\Id - \BP(\bk_0)) - \rho^\alpha } > Z
\delta \\
\text{implies} \quad
\abs{\lambda_\ell(\BA(\bk) (\Id - \BP(\bk)) - \rho^\alpha } >
\delta
\end{multline}
for every $\bk \in \CJ$. In particular, all of those eigenvalue branches which are at distance larger than $Z\delta$ from $\rho^\alpha$ at $\bk_0$ stay on the same side of the interval $[\rho^\alpha - \delta, \rho^\alpha + \delta]$ and contribute the same quantity to both sides of inequality~\eqref{eq:countinggj}. Finally, if
\begin{equation}
\abs{\lambda_\ell(\BA(\bk_0) (\Id - \BP(\bk_0)) - \rho^\alpha } \le Z
\delta,
\end{equation}
this means that $\bxi_0 + \btheta \in \CZ'$, and we cannot know their values at $\bk_1$ or $\bk_2$. However, since $\bxi_0 \in \CK$, there are at most $m n(\bxi;\CZ') < n(\bxi;\CG)$ values of $\bp = (\btheta,j)$ for which this holds. In the end, this means that
\begin{equation}
\begin{aligned}
1 &\le n(\bxi_0;\CG) - m n(\bxi_0;\CZ') \\&\le
\#\set{\bp \in \tilde \Theta : g_\bp(\bk_2) \le \rho^\alpha - \rho^S } -
\#\set{\bp \in \tilde \Theta : g_\bp(\bk_1) \le \rho^\alpha + \rho^S}
\end{aligned}
\end{equation}
which is our claim.
\end{proof}

We can now finish with the proof of Proposition~\ref{thm:uncoupledbs}.

\begin{proof}
{Proof of Proposition~\ref{thm:uncoupledbs}}
By construction of the functions $g_\bp$, inequality~\eqref{eq:countinggj} is equivalent to inequality~\eqref{eq:characoverlap} with $\delta = \rho^S$. The overlap exponent $S$ is seen from~\eqref{eq:propoverlap} to depend only on $\alpha$ and $d$, whereas the other parameters $\tilde \lambda$ and $c$ depend on the implied constants in Lemma~\ref{lem:descriptiongj}. This Lemma was taken from the constructions of~\cite[Section~7]{ParSob2010} where the constants are shown to depend only on the symbol norms of $\BA$ and $\BA^{\CO\CD}$.
\end{proof}



\section{The Dirac Operator} \label{sec:dirac}

In this section, we aim to get conditions on perturbations of the Dirac operator so that the gauge transform and, more importantly, all the theorems from Part II can be applied. Basic facts and theorems on the Dirac operator are found in~\cite{gilbertmurray,thaller}. We consider Dirac operators built through Clifford algebras, of which the usual two- and three-dimensional cases are examples. We are then able to explicitly describe perturbations to which we can apply the gauge transform method and recover the results of Sections~\ref{sec:besicovitch}--\ref{sec:cg}.

\subsection{Clifford algebras} \label{sec:clifford}

We give here basic facts about Clifford algebras used to construct the Dirac operator in the flat setting. They can be found in~\cite[Section~7]{gilbertmurray}. Let $\R^{p,q}$ be the euclidean space of dimension $p+q$ equipped with the canonical quadratic form $\eta$ of signature $(p,q)$. In our applications, we consider only the cases $\R^{0,d}$ (Euclidean) and $\R^{1,d}$ (Minkowski). We denote their orthonormal bases respectively $\set{\bv_1,\,\dotsc,\,\bv_d}$ and $\set{\bv_0,\bv_1,\,\dotsc,\,\bv_d}$. Consider the algebra $\FA_{p,q}$ generated by $\set{1,\bv_1,\,\dotsc,\,\bv_d}$ or $\{1,\bv_0,\,\dotsc,\bv_d\}$ with the relations
\begin{equation}\label{eq:anticommute}
\bv_j \bv_k + \bv_k \bv_j = - 2\eta_{jk}.
\end{equation}
It is easy to see that $\FA_{p,q}$ has dimension~$2^{p+q}$. For any subset $S := \set{s_1,\,\dotsc,\,s_k} \subset \set{0,\,\dotsc,\,d}$ (or of $\set{1,\,\dotsc,\,d}$ in the euclidean setting), we denote by $\bv_S$ the element $\bv_{s_1} \cdots \bv_{s_k} \in \FA_{p,q}$, where by convention $\bv_\varnothing = 1$. The Clifford algebra on $\R^{p,q}$ is isomorphic to the exterior algebra $\Lambda^*(\R^d)$.

From the anticommutation relation~\eqref{eq:anticommute}, we deduce that each pair of the $2^{p+q}$ generators of $\FA_{p,q}$ either commutes or anticommutes, according to the rule
\begin{equation}
\begin{cases}
\bv_j \bv_S = (-1)^{\abs S} \bv_S \bv_j & \text{if } j \not \in S, \\
\bv_j \bv_S = (-1)^{\abs S - 1} \bv_S \bv_j & \text{if } j \in S. \\
\end{cases}
\end{equation}
When $p+q$ is even, there is a faithful representation of $\FA_{p,q}$ acting on the spinor space $\C^{2^{(p+q)/2}}$. A specific representation by matrices constructed recursively is given in~\cite{Upmeier} in the Euclidean and Minkowski cases. This representation $\gamma$ has the property that for all $1 \le j \le d$, the matrix $\gamma_j := \gamma(\bv_j)$ is skew-hermitian and squares to $-\Id_{p+q}$, $\gamma_0 := \gamma(\bv_0)$ is hermitian and squares to the identity, and there is some $\abs c = 1$ so that the grading operator $\Gamma := c \prod_j \gamma_j$ is a diagonal matrix of the form
\begin{equation}
\Gamma =
\begin{pmatrix}
\Id_{(p+q)/2} & 0 \\
0 & - \Id_{(p+q)/2}
\end{pmatrix}.
\end{equation}
We can observe that for all $j$, $\Gamma \gamma_j = -
\gamma_j \Gamma$. The operator $\Gamma$ is called ``grading'' because it induces a $\Z_2$ grading on $\FA_{p,q}$. The even subalgebra of $\FA_{p,q}$ consists of all the elements commuting with $\Gamma$, while the odd subspace consists of all the anti-commuting elements. In particular, all the $\gamma_j$ are in the odd subspace, which is characterised as a product of an odd number of generators, while the even subalgebra is characterised as products of even number of generators.

\begin{lemm}\label{lem:anticom}
Let $\gamma$ be an element of the odd subspace. Then, as a matrix it has the form
\begin{equation}
\gamma :=
\begin{pmatrix}
\boldsymbol 0 & \BX \\
\BY & \boldsymbol 0
\end{pmatrix},
\end{equation}
where each of the blocks is a $m/2 \times m/2$ matrix.
\end{lemm}
\begin{proof}
This follows from a simple computation of the relation $\gamma \Gamma +
\Gamma \gamma = 0$ on the matrix elements.
\end{proof}

The representation $\gamma$ also allows us to see that, as a $C^*$-algebra, $\FA_{p,q}$ is naturally isomorphic to an algebra of operators on a Hilbert space $\FS_{p+q}$, which is called the \emph{spinor space}. When $m$ is even, we have that $\FS_{p+q} \cong \C^{2^{(p+q)/2}}$. Therefore, setting $m = 2^{(p+q)/2}$, we can use this representation to obtain operators in $\BS^\infty_m$.


\subsection{Dirac operators}

We define (spatial) Dirac operators differently depending on whether the number of spatial dimensions is even or odd.
\begin{defi}
Let $d$ be odd. The $d$-dimensional free Dirac operator $\BA_d$ is the first order system acting on spinors in $\RL^2(\R^d;\C^m)$, for $m = 2^{\frac{d+1}{2}}$ given by
\begin{equation}
\BA_d = \sum_{j=1}^d \gamma_j \del_j,
\end{equation}
where the $\gamma_j$ are given by the representation of $\FA_{1,d-1}$ in $\CL(\C^{m})$.
\end{defi}
\begin{defi}
Let $d$ be even. The $d$-dimensional free Dirac operator $\BA_d$ is the first order system acting on spinors $\RL^2(\R^d;\C^m)$, for $m = 2^{\frac d 2}$ given by
\begin{equation}
\BA_d = \sum_{j=1}^d \gamma_j \del_j,
\end{equation}
where the $\gamma_j$ are given by the representation of $\FA_{0,d}$ in $\CL(\C^{m})$.
\end{defi}

It is easy to see in both cases that $\BA_d^2 = - \Delta \Id_{m}$.
\begin{exem}
The two-dimensional Dirac operator with mass $M$ is given in~\cite[Equation 1.14]{thaller} as
\begin{equation}
\BA_{2,M} = -i\left(\sigma_1 \del_{x_1} + \sigma_2 \del_{x_2} \right) + \sigma_3 M,
\end{equation}
where $\sigma_1, \sigma_2,
\sigma_3$ are the Pauli matrices
\begin{align}
\sigma_1=
\begin{pmatrix}
0 & 1\\
1 & 0
\end{pmatrix},\
\sigma_2=
\begin{pmatrix}
0 & -i\\
i & 0
\end{pmatrix},\
\sigma_3=
\begin{pmatrix}
1 & 0\\
0 & -1
\end{pmatrix}.
\end{align}
It is a perturbation of order $0$ of the free Dirac operator. Indeed, the Pauli matrices can be used for a representation of the Clifford algebra $\FA_{0,2}$, and $\sigma_3$ corresponds to the grading operator $\Gamma$.
\end{exem}

\begin{exem}
The three-dimensional Dirac operator with mass $M$ from~\cite[Equation 1.11]{thaller} given by
\begin{equation}
\BA_{3,M} = -i\left(\gamma_1 \del_{x_1} + \gamma_2 \del_{x_2} + \gamma_3
\del_{x_3}\right) + \Gamma M
\end{equation}
is also a perturbation of order $0$ of the free Dirac operator. Here, the matrices $\gamma_j$ are the Dirac $\gamma$-matrices used as a representation of $\FA_{1,3}$, and our notation generalises this notion, following~\cite{Upmeier}.
\end{exem}

We now show that the operators $\BA_d$ are elliptic in the sense of Section~\ref{sec:perturbation}.

\begin{prop}\label{prop:diagdirac}
Let $m:= m(d)$ be the dimension of the spinor space on which $\BA_d$ acts. The operator $\BU \in \BS^0_m$ with symbol
\begin{equation}\label{eq:explicitunitary}
\bu(\bx,\bxi) := \frac{\bone_{\set{\abs \bxi\,\ge\,1}}(\bxi)}{\sqrt 2}\left(\Id_m + \frac{i}{\abs \bxi} \Gamma
\sum_{j=1}^d \bxi_j \gamma_j\right)+
\bone_{\set{\abs{\bxi}\,<\,1}}(\bxi)\Id_m
\end{equation}
is unitary. Furthermore, $\BU \BA_d \BU^* \in \BD\BE\BS_m^1$ and there is $\BR
\in \BS_m^{-\infty}$ such that the symbol of $\BU \BA_d \BU^* - \BR$ is $\abs{\bxi} \Gamma$.
\end{prop}

\begin{proof}
The symbol of the adjoint of $\BU$ is given, following~\eqref{eq:adjoint symbol}, by
\begin{equation}
\bu^\dagger(\bx,\bxi) = \frac{\bone_{\set{\abs \bxi\,\ge\,1}}(\bxi)}{\sqrt 2}\left(\Id_m - \frac{i}{\abs \bxi} \Gamma
\sum_{j=1}^d \bxi_j \gamma_j\right) + \bone_{\set{\abs{\bxi}\,<\,1}}(\bxi)
\Id_m
\end{equation}
and we can compute that
\begin{equation}
\begin{aligned}
\left[\bu\circ \bu^\dagger \right] (\bxi) &= \frac{\bone_{\set{\abs \bxi
\,\ge\,1}}(\bxi)}{2}\left(\Id_m -
\frac{1}{\abs{\bxi}^2} \sum_{j,k = 1}^d \Gamma^2 \gamma_j \gamma_k \bxi_j
\bxi_k\right) + \bone_{\set{\abs \bxi\,< \,1}}(\bxi) \Id_m \\
&= \frac{\bone_{\set{ \abs \bxi\,\ge\,1}}(\bxi)}{2} \left(\Id_m - \frac{1}{\abs{\bxi}^2} \sum_j \gamma_j^2
\bxi_j^2\right) + \bone_{\set{\abs \bxi\,< \,1}}(\bxi) \Id_m\\
&= \Id_m.
\end{aligned}
\end{equation}
In a very similar fashion, we see that the symbol of $\BU \BA_d \BU^*$ is given by
\begin{equation}
[\bu \circ \ba_d \circ \bu^\dagger](\bxi) = \bone_{\set{\abs \bxi\,\ge\,1 }}(\bxi) \abs \bxi\Gamma + \bone_{\set{\abs \bxi\,<\,1}}(\bxi) \ba_d(\bxi).
\end{equation}
This proves our claim where $\BR \in \BS_m^{-\infty}$ has symbol
\begin{equation}
\br(\bxi)= \bone_{\set{\abs \bxi\,<\,1}} (\bxi) (\ba_d(\bxi) - \abs \bxi
\Gamma).\qedhere
\end{equation}
\end{proof}

We now see that for $d = m = 2$, the operators $\BA_2 + \BB$, $\BB \in \BS^\beta_m$, $\beta < 1$ are unitarily equivalent to an operator satisfying the hypotheses of Theorems~\ref{thm:aexpconcrete} and Theorem~\ref{thm:bs}, which proves that we generically have a complete asymptotic expansion for the density of states, and that if $\BB$ is periodic then $\BA$ has the Bethe--Sommerfeld property. In other words, the following two theorems are proved, which are more precise reformulations of Theorems~\ref{thm:dirac2dids} and~\ref{thm:dirac2dbs}.

\begin{theo}
Let $\beta < 1$ and $\BA = \BA_2 + \BB$, where $\BB \in \BS_2^\beta$ satisfies the generic conditions \textbf{A}, \textbf B and \textbf C. Then, for every $K > - 2$ there is a finite set $L \subset (0,2+K)$ so that for every $j \in L
\cup \set{0}$ there are constants $C_{j}^\pm$, $C_{j,\log}^\pm$ such that
\begin{equation}
N^\pm(\BA; \lambda) = C_0^\pm \lambda^2 + \sum_{j\,\in\,L} \left(C_j^\pm
\lambda^{2 - j} + C_{j,\log}^\pm \lambda^{2-j} \log \lambda \right) +
\bigo{\lambda^{-K}}
\end{equation}
as $\lambda \to \infty$.
\end{theo}


\begin{theo}
Let $\beta < 1 $ and $\BA = \BA_2 + \BB$, where $\BB \in \BS_2^\beta$ is periodic. Then, $\BA$ has the Bethe--Sommerfeld property, i.e. there exists $\lambda_0 > 0 $ such that the spectrum of $\BA$ contains intervals $(-\infty, -\lambda_0]$ and $[\lambda_0,\infty)$.
\end{theo}

We now want to address the question of the perturbations that are allowed whenever $d \ge 3$.


\begin{prop}\label{prop:permittedperturbations}
For $\beta < 1$, and $0 \le j \le d$ (with $0$ omitted when $d$ is even), let $B_{\Id}, B_\Gamma, B_j \in \BS_1^\beta$ be scalar pseudo-differential operators of order $\beta$, and put
\begin{equation}
\BB = B_{\Id} \Id_m + B_\Gamma \Gamma + \sum_j B_j \gamma_j.
\end{equation}
Then, there are operators $\BB' \in \BU\BS^\beta_m$, $\BR \in \BS_m^{\beta -1 }$ and $\tilde \BB \in
\BS_m^\beta$ whose symbol has image in the odd subspace of $\FA_{p,q}$ such that
\begin{equation}\label{eq:afterunit}
\BU (\BA_d + \BB) \BU^* = \Op(\abs \bxi) \Gamma + \BB' + \tilde \BB + \BR
\end{equation}
\end{prop}

\begin{proof}
The unitary operator $\BU$ from~\eqref{eq:explicitunitary} can be written as
\begin{equation}
\BU = \frac{1}{\sqrt 2} \left(\Id_m + \sum_{j=1}^d U_j \Gamma \gamma_j\right)
\mod \BS_m^{-\infty}.\label{eq:Udecomposed}
\end{equation}
Here, $U_j \in \BS_1^0$ are scalar pseudo-differential operators given by
\begin{equation}
U_j = \Op\left(\frac{i \bxi_j \chi(\bxi)}{\abs \bxi}\right),
\end{equation}
where $\chi$ is a smooth function supported in $\set{\abs{\bxi} \ge 1/2}$ and $\chi(\bxi) \equiv 1$ for all $\abs \bxi \ge 3/4$. We now compute $\BU B_\gamma \gamma \BU^*$ for different values of $\gamma$. All the sums range from $1$ to $d$ with additional restrictions, we have only written the restrictions to make notation lighter. For $1 \le j \le d$ we have
\begin{multline}\label{eq:unittrans}
\BU B_{j} \gamma_j \BU^*=\\
 \frac{1}{2}\Bigg(B_{j} \gamma_j + \sum_{k
\neq j}[U_k;B_{j}]\Gamma \gamma_k \gamma_j - \left(U_j B_{j} + B_{j} U_j\right) \Gamma + \sum_{k} \left(U_k B_{j} U_j + U_j B_j U_k\right) \gamma_k \\
- \sum_{k} U_k B_{j} U_k \gamma_j -
\sum_{\substack{\ell \neq j \\ k \neq j \\k\,<\,\ell}}\left([U_\ell;B_{j} U_k] + [B_{j};U_k]U_\ell \right) \gamma_k \gamma_\ell
\gamma_j
\bigg) \mod \BS_m^{- \infty}.
\end{multline}
Let us have a careful look at each of the six terms in Equation~\eqref{eq:unittrans}. The second and the last terms involve commutators of operators with scalar-valued symbols, they are in $\BS_m^{\beta - 1}$ and we put $\BR_{j}$ as their sum. The third term is in $\BU\BS^\beta_m$, and we denote it $\BB'_j$. Finally, the first, fourth and fifth term are readily seen to have symbols in the odd subspace, we put $\tilde \BB_{j}$ as their sum.

The operator $\BU B_{0} \gamma_0 \BU^*$ is computed similarly as in~\eqref{eq:unittrans} with some of the terms vanishing. It is given by
\begin{multline}
\BU B_0 \gamma_0 \BU^* =
\frac 1 2 \Bigg(B_0 \gamma_0 + \sum_{k} [U_k;B_0]
\Gamma \gamma_k \gamma_0 - \sum_k U_k B_0 U_k \gamma_0  \\- \sum_{k<\ell} ([U_\ell;B_0 U_k] + [B_0;U_k]U_\ell) \gamma_k \gamma_\ell \gamma_0\Bigg) \mod
\BS_m^{-\infty}.
\end{multline}
The first and third term have image in the odd subspace, we put $\tilde
\BB_0$ as their sum. The second and last terms involve commutators of operators with scalar-valued symbols, as such they are in $\BS_m^{\beta - 1}$ and we put $\BR_0$ as their sum. We note that there are no uncoupled terms.

The operator $\BU B_\Gamma \Gamma \BU^*$ is given by
\begin{multline}\label{eq:Gammatrans}
\BU B_\Gamma\Gamma \BU^* =
\frac{1}{2}\Bigg(B_\Gamma \Gamma - \sum_{k}(U_k B_\Gamma + B_\Gamma U_k) \gamma_k \\
- \sum_{k} U_k B_\Gamma U_k \Gamma +
\sum_{k < \ell}\left([U_\ell;B_\Gamma U_k] + [B_\Gamma;U_k]U_\ell \right)\Gamma
\gamma_\ell \gamma_k
\bigg) \mod \BS_m^{-\infty}.
\end{multline}
This time, the first and third terms are seen to be in $\BU\BS^\beta_m$ and we put their sum as $\BB'_\Gamma$. The second term has symbol in the odd subspace and we denote it by $\tilde
\BB_\Gamma$. The last term can be seen to be in $\BS_m^{\beta - 1}$ and we denote it by $\BR_{\Gamma}$.

Finally, the operator $\BU B_{\Id} \Id_m \BU^*$ is given by
\begin{multline}\label{eq:Idtrans}
\BU B_{\Id} \Id_m \BU^* =
\frac{1}{2}\Bigg(B_{\Id} \Id_m +
\sum_{k}[U_k;B_{\Id}]\Gamma \gamma_k + \sum_{k} U_k B_{\Id} U_k \Id_m \\
+\sum_{k\,<\,\ell}\left([U_\ell;B_{\Id}U_k] + [B_{\Id};U_k]U_\ell \right)
\gamma_\ell \gamma_k
\Bigg) \mod \BS_m^{-\infty}.
\end{multline}
This time, we see that the first and third terms are in $\BU\BS_m^\beta$, we put their sum as $\BB'_{\Id}$, while the second and last terms are in $\BS_m^{\beta - 1}$ and we put their sum as $\BR_{\Id}$.

Finally, put $\tilde \BR \in \BS_m^{-\infty}$ as the sum of the remainders $\mod \BS_m^{-\infty}$ obtained at every step. Combining all our computations and Proposition~\ref{prop:diagdirac} gives us that~\eqref{eq:afterunit} holds with
\begin{equation}
\begin{aligned}
\BB' &= \BB'_{\Id} + \BB'_{\Gamma} + \sum_{j=1}^d \BB'_j \\
\tilde \BB &= \tilde \BB_{\Id} + \tilde \BB_{\Gamma} + \sum_{j=0}^d \tilde
\BB_j \\
\BR &= \tilde \BR + \BR_{\Id} + \BR_\Gamma + \sum_{j=0}^d \BR_j.
\end{aligned}\qedhere
\end{equation}
\end{proof}

The next theorem follows and includes Theorem~\ref{thm:dirac3dids} as a special case when $d = 3$.
\begin{theo}
Let $m(d)$ be the dimension of the spinor space on which $\BA_d$ acts. For $\beta \le 1/2$ and $0 \le j \le d$ (with $0$ omitted when $d$ is even) let $B_\Gamma, B_j, B_{\Id} \in \BS^\beta$ be scalar pseudo-differential operators satisfying Conditions~\ref{condI}--\ref{condIV}, and put
\begin{equation}
\BB = B_{\Id} \Id_m + B_\Gamma \Gamma + \sum_{j=0}^d B_j \gamma_j,
\end{equation}
and $\BA = \BA_d + \BB$. Then, putting $\gamma^* = \max \set{\beta - 1, 2
\beta - 1}$, there exists a finite set $L \subset (0, 1 - \gamma^*)$ and constants $C_0^\pm$ and $C_{j,q}^\pm$, $0 \le q \le d-1$, $j \in L$ such that
\begin{equation}
N^\pm(\BA;\lambda) = C_0^\pm \lambda^d + \sum_{j \in L} \sum_{q = 0}^{d-1} C_{j,q}^\pm \lambda^{d-j} \log^q \lambda + \bigo{\lambda^{d-1 + \gamma^*}}
\end{equation}
as $\lambda \to \infty$.
\end{theo}

\begin{proof}
It follows from Proposition~\ref{prop:permittedperturbations} that $\BU \BA \BU^*$ satisfies the hypotheses of Theorem~\ref{thm:aexpconcretecut} with $\gamma^* =
\max\set{\beta - 1, 2 \beta - 1}$. In particular, the restricted asymptotics of the IDS given in that theorem are true for such operators with $\alpha = 1$.
\end{proof}

Finally, in some highly non-generic cases we can get complete asymptotic expansions and the Bethe--Sommerfeld property for $d$-dimensional Dirac operators with $d \ge 3$. We state both results and observe that they follow directly from the fact that after conjugation by $\BU$, these operators are uncoupled.

\begin{theo}
Let $m(d)$ be the dimension of the spinor space on which $\BA_d$ acts, $\beta < 1$ and $\BB \in \BU\BS_m^\beta$ satisfying Conditions~\ref{condI}--\ref{condIV}. Put $\BA = \BA_d + \BU^* \BB \BU$. Then, $N^\pm(\BA;\lambda)$ satisfies the complete asymptotic expansion~\eqref{eq:asymptoticsconcrete} with $\alpha = 1$.
\end{theo}

\begin{theo}
Let $m(d)$ be the dimension of the spinor space on which $\BA_d$ acts, $\beta < 1$ and $\BB \in \BU\BS_m^\beta$ be periodic. Put $\BA = \BA_d + \BU^* \BB
\BU$. Then, $\BA$ has the Bethe--Sommerfeld property.
\end{theo}




\begin{thebibliography}{99}
\bibitem{barbpar}Barbatis, Gerassimos; Parnovski, Leonid Bethe–Sommerfeld conjecture for pseudodifferential perturbation, Commun. Partial Differ. Equations, Volume 34 (2009) no. 4-6, pp. 383-418 | DOI | MR | Zbl
\bibitem{BirmanSolomyak}Birman, Mikhail S.; Solomjak, Mikhaĭl Spectral Theory of Self-Adjoint Operators in Hilbert Space, Mathematics and its Applications. Soviet Series, Kluwer Academic Publishers, 1987 | Zbl | DOI
\bibitem{CoburnMoyerSinger}Coburn, Lewis A.; Moyer, R. D.; Singer, Isadore M. ${C}^*$-algebras of almost periodic pseudo-differential operators, Acta Math., Volume 130 (1973), pp. 279-307 | DOI | MR | Zbl
\bibitem{LaurentVuNgoc}Charles, Laurent; Vũ Ngọc, San Spectral asymptotics via the semiclassical Birkhoff normal form, Duke Math. J., Volume 143 (2008) no. 3, pp. 463-511 | DOI | MR | Zbl
\bibitem{Dixmier}Dixmier, Jacques Von Neumann Algebras, North-Holland Mathematical Library, 27, North-Holland, 1981 (with a preface by E. C. Lance, Translated from the second French edition by F. Jellett) | MR | Zbl
\bibitem{dahlbergtrubowitz}Dahlberg, Bjorn E. J.; Trubowitz, Eugene A remark on two-dimensional periodic potentials, Comment. Math. Helv., Volume 57 (1982) no. 1, pp. 130-134 | DOI | MR | Zbl
\bibitem{gilbertmurray}Gilbert, John E.; Murray, Margaret A. M. Clifford algebras and Dirac operators in harmonic analysis, Cambridge Studies in Advanced Mathematics, 26, Cambridge University Press, 1991 | DOI | MR | Zbl
\bibitem{HormanderIII}Hörmander, Lars The analysis of linear partial differential operators. III Pseudo-differential operators, Classics in Mathematics, Springer, 2007 (reprint of the 1994 edition) | DOI | MR | Zbl
\bibitem{Ivrii2018}Ivrii, Victor Complete Semiclassical Spectral Asymptotics for Periodic and Almost Periodic Perturbations of Constant Operators, Springer (2019), pp. 583-606 | Zbl
\bibitem{Kuchment}Kuchment, Peter Floquet theory for partial differential equations, Operator Theory: Advances and Applications, 60, Birkhäuser, 1993 | DOI | MR | Zbl
\bibitem{MorParSht2014}Morozov, Sergey; Parnovski, Leonid; Shterenberg, Roman Complete asymptotic expansion of the integrated density of states of multidimensional almost-periodic pseudo-differential operators, Ann. Henri Poincaré, Volume 15 (2014) no. 2, pp. 263-312 | DOI | MR | Zbl
\bibitem{Naimark}Naĭmark, Mark A. Normed algebras, Wolters-Noordhoff Series of Monographs and Textbooks on Pure and Applied Mathematics, Wolters-Noordhoff Publishing, Groningen, 1972 (translated from the second Russian edition by Leo F. Boron) | MR | Zbl
\bibitem{parnovski}Parnovski, Leonid Bethe–Sommerfeld conjecture, Ann. Henri Poincaré, Volume 9 (2008) no. 3, pp. 457-508 | DOI | MR | Zbl
\bibitem{ParSob2001}Parnovski, Leonid; Sobolev, Alexander V. On the Bethe–Sommerfeld conjecture for the polyharmonic operator, Duke Math. J., Volume 107 (2001) no. 2, pp. 209-238 | DOI | MR | Zbl
\bibitem{ParSob2010}Parnovski, Leonid; Sobolev, Alexander V. Bethe–Sommerfeld conjecture for periodic operators with strong perturbations, Invent. Math., Volume 181 (2010) no. 3, pp. 467-540 | DOI | MR | Zbl
\bibitem{ParSht2012}Parnovski, Leonid; Shterenberg, Roman Complete asymptotic expansion of the integrated density of states of multidimensional almost-periodic Schrödinger operators, Ann. Math., Volume 176 (2012) no. 2, pp. 1039-1096 | DOI | MR | Zbl
\bibitem{ParSht2016}Parnovski, Leonid; Shterenberg, Roman Complete asymptotic expansion of the spectral function of multidimensional almost-periodic Schrödinger operators, Duke Math. J., Volume 165 (2016) no. 3, pp. 509-561 | DOI | MR | Zbl
\bibitem{ParSht2019}Parnovski, Leonid; Shterenberg, Roman Perturbation theory for almost-periodic potentials I: one-dimensional case, Commun. Math. Phys., Volume 366 (2019) no. 3, pp. 1229-1257 | DOI | MR | Zbl
\bibitem{Rozenbljum1978}Rozenbljum, Grigorii V. Near-similarity of operators and the spectral asymptotic behavior of pseudodifferential operators on the circle, Tr. Mosk. Mat. O.-va, Volume 36 (1978), p. 59-84, 294 | MR | Zbl
\bibitem{Shubin1978}Shubin, Mikhail A. Almost periodic functions and partial differential operators, Usp. Mat. Nauk, Volume 33 (1978) no. 2, p. 3-47, 247 | MR
\bibitem{Shubin1979a}Shubin, Mikhail A. Pseudodifferential almost-periodic operators and von Neumann algebras, Trans. Mosc. Math. Soc., Volume 35 (1979), pp. 103-166 | Zbl
\bibitem{Shubin1979}Shubin, Mikhail A. Spectral theory and the index of elliptic operators with almost-periodic coefficients, Usp. Mat. Nauk, Volume 34 (1979) no. 2, pp. 95-135 | MR | Zbl
\bibitem{Skriganov}Skriganov, Maxim M. Geometric and arithmetic methods in the spectral theory of multidimensional periodic operators, Tr. Mat. Inst. Steklova, Volume 171 (1985), pp. 3-122 | MR
\bibitem{Sobolev2005}Sobolev, Alexander V. Integrated density of states for the periodic Schrödinger operator in dimension two, Ann. Henri Poincaré, Volume 6 (2005) no. 1, pp. 31-84 | DOI | MR | Zbl
\bibitem{Sobolev2006}Sobolev, Alexander V. Asymptotics of the integrated density of states for periodic elliptic pseudo-differential operators in dimension one, Rev. Mat. Iberoam., Volume 22 (2006) no. 1, pp. 55-92 | DOI | MR | Zbl
\bibitem{TaylorII}Taylor, Michael E. Partial differential equations II. Qualitative studies of linear equations, Applied Mathematical Sciences, 116, Springer, 2011 | DOI | MR | Zbl
\bibitem{thaller}Thaller, Bernd The Dirac equation, Texts and Monographs in Physics, Springer, 1991 | DOI | MR | Zbl
\bibitem{Upmeier}Upmeier, Harald Dirac operator and real structure on Euclidean and Minkowski spacetime, Noncommutative geometry and the standard model of elementary particle physics (Hesselberg, 1999) (Scheck, Florian et al., eds.) (Lecture Notes in Physics), Volume 596, Springer, 2002, pp. 136-151 | DOI | MR | Zbl
\bibitem{Weinstein1977}Weinstein, Alan Asymptotics of eigenvalue clusters for the Laplacian plus a potential, Duke Math. J., Volume 44 (1977) no. 4, pp. 883-892 | MR | Zbl
\end{thebibliography}
\end{document}
